% to cite: https://www.cambridge.org/core/services/aop-cambridge-core/content/view/78B4021417116DCC65F05177645AA34F/S0003055422000156a.pdf/indias-farmers-protest-an-inclusive-vision-of-indian-democracy.pdf


\documentclass[12pt]{article}
% Self-contained local render. Empirical assets are resolved only from the
% repository-level figures/ and tables/ folders.
\usepackage[margin=1in]{geometry}
\usepackage{amsmath,amssymb}
\usepackage{natbib}
\usepackage{hyperref}
\usepackage{placeins}
\usepackage{chngcntr}
\usepackage{comment}
\usepackage{fancyvrb}
\usepackage{listings}
\usepackage{subcaption}
\usepackage{color}
\usepackage{graphicx}
\usepackage{xcolor}
%\usepackage{adjustbox}
\usepackage[export]{adjustbox}
% \usepackage{natbib}
\usepackage{setspace}
\usepackage{float}
\usepackage{verbatim}
\usepackage{booktabs}
\usepackage{array}
\usepackage{etoolbox}
\usepackage{todonotes}
\usepackage[english]{babel}
% \usepackage[backend=biber, style=nature, autocite=superscript, defernumbers=true]{biblatex}
% \usepackage[natbib=true]{bibentry} 
\newcommand{\sym}[1]{\ensuremath{^{#1}}}
\newcommand{\gray}[1]{\textcolor{gray}{#1}}
\newcommand{\red}[1]{\textcolor{red}{#1}}
\providecommand{\textt}[1]{\texttt{#1}}
\providecommand{\printbibliography}{}
\providecommand{\omicron}{\mathrm{o}}
\newcommand{\subfloat}[2][]{\subcaptionbox{#1}{#2}}
% Optional result suffix supplied by a render driver. The ordinary report uses
% an empty suffix; the above-median report defines this as _rice_high.
\providecommand{\resultsuffix}{}

% Resolve every table call against the repository root. Missing legacy assets
% remain visible as named placeholders and are never sourced elsewhere.
\let\SourceInput\input
\renewcommand{\input}[1]{%
  \IfFileExists{../../#1\resultsuffix.tex}{\SourceInput{../../#1\resultsuffix.tex}}{%
  \IfFileExists{../../#1}{\SourceInput{../../#1}}{%
    \IfFileExists{../../#1.tex}{\SourceInput{../../#1.tex}}{%
      \fbox{\scriptsize Missing table: \detokenize{#1}}}}}}

% Resolve figures only from the repository-level figures folder.
\let\SourceIncludeGraphics\includegraphics
\newcommand{\MissingGraphic}[2][]{%
  \adjustbox{#1}{\fbox{\parbox[c][2.2cm][c]{4cm}{\centering\scriptsize
  Missing figure: \detokenize{#2}}}}}
\makeatletter
\renewcommand{\includegraphics}[2][]{%
  \filename@parse{#2}%
  \edef\ResultGraphic{../../\filename@area\filename@base\resultsuffix.\filename@ext}%
  \IfFileExists{\ResultGraphic}{\SourceIncludeGraphics[#1]{\ResultGraphic}}{%
  \IfFileExists{../../#2}{\SourceIncludeGraphics[#1]{../../#2}}{%
  \IfFileExists{../../#2.pdf}{\SourceIncludeGraphics[#1]{../../#2.pdf}}{%
  \IfFileExists{../../#2.png}{\SourceIncludeGraphics[#1]{../../#2.png}}{%
  \IfFileExists{../../#2.jpg}{\SourceIncludeGraphics[#1]{../../#2.jpg}}{%
  \IfFileExists{../../#2.jpeg}{\SourceIncludeGraphics[#1]{../../#2.jpeg}}{%
  \MissingGraphic[#1]{#2}}}}}}}}
\makeatother
%\addbibresource{all_references.bib} %\addbibresource{appendix_references.bib} 
\sloppy
\AtBeginEnvironment{table}{\singlespacing}
\AtBeginEnvironment{minipage}{\singlespacing}


\title{Democratic Accountability and Environmental Regulation}

\author{
Gemma Dipoppa \thanks{Columbia University, Department of Political Science (gd2733@columbia.edu)} 
\and Saad Gulzar \thanks{Notre Dame University, Department of Political Science, (sgulzar@nd.edu)} \and Annalisa Pezone \thanks{New York University, Department of Politics, (ap8979@nyu.edu)} 
\and Anzony Quispe \thanks{Sciences Po University, Department of Economics, (anzony.quisperojas@sciencespo.fr)} 
}


%note: add in the thanks of the paper “Mart Buntaine, and participants to EGAP 2026 and Climate Pipeline 2026 Columbia
\date{\today}


\begin{document}
\setcounter{page}{0}
	\maketitle
	\thispagestyle{empty}
	
	
	\begin{singlespace}
\begin{abstract}

%Environmental governance presents a challenge for democratic accountability: pollution yields concentrated benefits while distributing costs broadly, creating weak incentives for politicians to act. We examine the conditions under which political action comes together and falls apart in the context of air pollution from crop-residue burning in South Asia -- one of the world’s largest public-health crises responsible for millions of premature deaths each year. Using a ten-year panel of 12 million 5 km² grid-cells linking satellite-detected fires, wind fields, and political data, we exploit changes in wind direction that temporarily alter which constituencies bear the costs of pollution. Polluting behavior falls by 15 percent when smoke is poised to affect a politician’s own constituents, demonstrating that electoral accountability can prompt environmental action even in hard cases. However, these accountability effects erode under competing pressures: we show that lobbying from the 2020–2021 farmers’ protests and politicians’ private agricultural interests lead to increased burning by 10 and 13 percent, respectively. 
%Reflecting the pressures that politicians face, we further show that their social media communication 
%We further show that policiains' social media farmer policiains talk more about farmers, but less about paddy or air pollution
%Contrary to democratic skeptics, these results show that environmental protection can emerge even under diffuse electoral incentives, yet such regulation remains acutely vulnerable to organized interests that stand to benefit from pollution.

%\bigskip
%\textbf{Gemma:} 
Environmental governance presents a core challenge for democracies: pollution yields concentrated benefits and diffuse harms, weakening incentives for political action. We investigate when politicians act to curb pollution, and what breaks this accountability, in the context of crop-residue burning in South Asia, which kills an estimated 2 million people annually. Linking satellite-detected fires, wind fields, and political data across 808 rural constituencies in northern India over a decade, we use a difference-in-differences design that exploits wind-direction shifts altering whether smoke from a given location drifts into a politician's own constituency or a neighbor's. When their constituents bear the pollution costs, politicians curb fires by 13 percent. Yet two forces systematically undermine this accountability: sanctioning weakens as organized farmer protests pressure politicians to tolerate burning, and selection fails when constituencies elect politicians with private agricultural interests, both increasing fires even where politicians' own voters are harmed. These findings demonstrate that democratic accountability can generate environmental protection under diffuse incentives, but sustaining it requires institutions robust to the concentrated interests that benefit from pollution.

%Environmental governance presents a challenge for democratic accountability: pollution yields concentrated benefits and diffused harms, weakening incentives for politicians to act. We examine the conditions under which political action comes together and falls apart in the context of air pollution from crop-residue burning in South Asia -- one of the world’s largest public-health crises. Using a ten-year panel of 12 million 5 km² grid-cells linking satellite-detected fires, wind fields, and political data, we exploit changes in wind direction that temporarily alter which constituencies bear the costs of pollution. Polluting behavior falls by 15 percent when smoke is poised to affect a politician’s own constituents, demonstrating that electoral accountability can prompt environmental action even under diffused incentives. However, these accountability effects erode under multiple principals and adverse selection: we show that lobbying from the 2020–2021 farmers’ protests and politicians’ private agricultural interests offset the gains in accountability, leading to more fires. Contrary to democratic skeptics, these results show that environmental protection can emerge even under diffuse electoral incentives, yet such regulation remains acutely vulnerable to organized interests that stand to benefit from pollution.
 %182 words



\end{abstract}

%


\end{singlespace}

\clearpage
\section{Introduction}


Few areas of governance reveal the limits of democratic accountability more clearly than environmental management. A growing body of evidence shows that voters rarely prioritize environmental issues \citep{singh_thachil_2023_airpollution}, leaving politicians with little electoral incentive to address pollution. Even when citizens do care about environmental problems, these preferences seldom translate into policies addressing environmental damage \citep{dechezlepretre2025fighting, hsiao2025climate}.%meaningful political action \citep{hsiao2025climate}.
\footnote{Examples of successful action on pollution come from China (see \citet{bar_gulzar2025} for a review), where, however, politicians are unelected and face different incentives than in democracies \citep{buntaine2024does}. Moreover, scholars increasingly document that a transition to a democracy may be associated with poorer environmental outcomes \citep{sanford2023democratization, chesler2023democracy}.} Instead, environmental regulation often reflects the influence of concentrated and organized lobbies \citep{mcalexander2020elections, stokes2020short, harding2024buying, gonzalez2024money}, reinforcing a growing pessimism about the prospects for democratic political action against environmental degradation. 

Yet this pessimism raises a crucial question: do politicians ever act to protect the environment, and if so, under what conditions?\footnote{In contrast to the political context, scholars have recently examined these questions within bureaucracies \citep{dipoppa2024bureaucrat, buntaine2024does}.} Studying this question is challenging because environmental degradation, unlike most policy domains, is a slow-moving process whose harms unfold gradually and often escape citizens’ immediate perception. Environmental humanities scholars have described this temporal disjuncture as ``slow violence'' \citep{Nixon2013}, and political scientists argue that this very disconnect between environmental harm and lived experience helps explain why environmental issues remain politically marginal and why democratic systems so often struggle to generate sustained action against them \citep{grossman2026politics}.

% summary of results
In contrast to the powerful incentives for inaction on the environment described above, we instead document that there is still a promise of accountability: we show that politicians do take meaningful steps to curb pollution when their own constituents bear its toxic effects. However, we also demonstrate that this control is systematically undermined. Building on classic literature, we examine two key channels of political accountability whose erosion can break down politicians' responsiveness to citizens' welfare: sanctioning and selection. Specifically, we highlight how sanctioning weakens under pressure from multiple principals, as citizens’ interests compete with those of concentrated lobbies, and how political selection falters when politicians’ self-interest diverges from citizens' well-being \citep{Fearon_1999}. Understanding these mechanisms is crucial for designing institutions that can channel democratic accountability towards durable environmental improvements. 



% context
We study these dynamics in the context of South Asia's air pollution crisis, one of the largest public health emergencies on the planet estimated to kill 2 million people every year, with children as one of its most vulnerable victims \citep{dandona2020subnational}.\footnote{In comparison, COVID-19 at its peak killed 6 million people in the first two years worldwide.}  Each winter, a thick smog so vast it is visible from space, blankets northern India as millions of farmers burn leftover crop residue after rice harvest. Roughly 60 percent of this seasonal pollution originates from these fires \citep{bikkina2019air, lan2022air}.  Figure \ref{fig:delhi} shows Delhi during smog season. Despite the scale of the crisis, the visibility and consequentiality of the problem, and the illegality of crop burning, the practice continues today. This paper studies how politics explains variation in this persistence. %This persistence reflects powerful incentives for political inaction: monitoring millions of farmers is costly, pollution drifts across regions obscuring political responsibility, and organized farmer groups lobby for leniency (New Indian Express 2023)\todo{cite}. 




\begin{figure}
    \centering
    \includegraphics[width=0.5\linewidth]{figures/cnn.png}
    \caption{Gate of India, Delhi, during and after the smog season (Source: \citealp{cnn2020lockdown})}
    \label{fig:delhi}
\end{figure}

%Data and estimation strategy
The first part of the paper examines how democratic accountability shapes political action towards environmental regulation. We propose that politicians will be most incentivized to stop crop burning in places from which smoke pollutes their own constituents, rather than a neighbors'. We leverage wind as a source of variation to resolve the challenge of identifying when politicians' incentives bind as wind direction determines where smoke from a burning field would travel. Because seasonal wind patterns in the region are broadly predictable, politicians will have reasonable expectations about them. When wind direction shifts in a month, the same 5km$^2$ grid-cell in a politician's constituency will switch from sending smoke into neighboring constituencies to sending it into the politician's own, raising the political cost of inaction. We exploit these within-location changes in a difference-in-differences design, combining satellite-detected fires from VIIRS and MODIS \citep{VIIRS_Instrument_Overview}, and wind direction and speed from \citet{ERA5_monthly_pressure_levels_2023} across 808 rural political jurisdictions in northern India over ten years (2012–2022), for approximately 9 million observations.


We find that fires decrease by 13 percent right after a change in wind direction switches a location from polluting neighboring jurisdictions to a politician's own constituency, an effect that is 34 percent larger in rice producing areas. These results are robust to a host of tests that relax the identifying assumption of the DiD design and vary how we define the outcome, the treatment, and the sample.

We further bolster our interpretation by ruling out that these findings are not explained by the incentives of two key non-politician actors: farmers and bureaucrats. First, if the decline were driven by farmers’ pro-sociality rather than by political accountability, we should observe fewer fires when more people live downwind from farmers. In fact, the opposite holds: farmers tend to burn more when their crop burning is likely to affect more people, and the reduction in fires emerges only when political incentives are active. Second, to test if bureaucrats' incentives explain our results instead of politicians', we use the fact that the jurisdictions for the two actors overlap imperfectly in India, allowing us to examine the impact of the same location being a source of pollution for a politician's constituents, the bureaucrat's constituents, both, or neither. We show that our results persist when a location is a source of pollution for politicians but not for bureaucrats, confirming that the logic of political incentives drives fires.\footnote{Interestingly, we also find that when the location pollutes both the politician's and bureaucrat's jurisdictions, the total reduction in fires is lower than when the location is a source of pollution only for politicians or only for bureaucrats. This is consistent with the interpretation that political effort is sensitive to free-riding considerations in the presence of incentives for action for other government agents.} 

Ultimately, this first finding provides powerful evidence that there are prospects for political accountability in environmental management even in a setting where action would seem unlikely: indeed, our evidence suggests that politicians are sensitive to environmental pollution when the well-being of their own constituents is at stake. We use this accountability effect as a benchmark to study its limits, focusing on the channels of weakening sanctions via multiple principals, farmers and citizens, and adverse selection through the election of self-interested politicians. 


%Result 2: lobbying
The second part of the paper examines how accountability weakens when it competes with the interests of actors who benefit from pollution-intensive activities: farmers organized into lobbying groups, and politicians with agricultural ties. %The second part of the paper begins by considering how accountability weakens when citizen sanctioning in democracies competes with organized interests. 
Around the world, from truck driver strikes in South Korea, Canada, and the UK, to teacher protests during Covid, sectors have lobbied for state intervention amid rising economic pressures. We investigate this dynamic in the context of the largest protest movement in the world's history, the Farmers’ Movement of 2020–21 in India, during which more than 25 million farmers mobilized to secure policy concessions, including the decriminalization of crop burning. We examine the impacts of 390 farmers' protest events (2020-2021) on subsequent crop burning in a stacked difference-in-differences framework.

%We investigate this dynamic in the context of the largest protest movement in the world's history, the Farmers’ Movement of 2020–21 in India, where over 25 million people took to the streets against three agricultural laws perceived to undermine government support for agriculture through the Minimum Support Price (MSP) system. After about 18 months of protests, the movement secured major concessions from the government: the repeal of the three contentious laws, which were never implemented effectively \citep{FarmLawsRepeal}; and, importantly for this paper, the decriminalization of crop burning \citep{newindianexpress2021centre}, a long-standing concern and grievance among farmers toward the state.


%result 2
We show that crop‐residue fires increase by 21 percent following protests, consistent with politicians relaxing environmental control in response to intense farmer lobbying. Having demonstrated that protests increased the production of pollution, we next show that they also erode the accountability incentives to politicians' own constituents that we document in the first part of this paper. To do so, we compare the reduction in fires produced by accountability effects before and after protests, and show that in places that did not have protests, the accountability-driven decrease in fires persists. In contrast, places that switch to having protests show significantly more crop burning afterwards. This suggests that politicians relax environmental control in response to concentrated lobbying pressure, even when such inaction harms their own constituents.


%Result 3: personal interest
In addition to lobbying, we examine another set of actors who may benefit from pollution: politicians with agricultural ties. To capture these interests, we use variation in politicians’ professional involvement in agriculture, which may give them a direct economic stake in practices associated with crop burning. We proxy these ties by collecting the candidacy papers of all 2,565 elected politicians in our ten-year sample and coding whether they are agriculturalists by profession. We then estimate the impact of a constituency switching from a non-agriculturalist to an agriculturalist politician on crop burning using a stacked difference-in-differences framework.
%In addition to lobbying, we examine another lever that undermines political accountability: adverse selection of self-interested politicians. To gain traction on politician's self-interest, we use the idea that in their private lives, politicians have varying degrees of financial interests in the production of crop burning. We proxy for self-interest by gathering the candidacy papers of all 2,565 elected politicians in our ten year sample and coding if politicians are agriculturalists by profession. 

%3: selection+downup
We find that when a constituency switches to electing an agriculturalist politician, the rate of crop burning increases by 24 percent, indicating a more permissive policy environment when politicians have personal stakes in agriculture.  We next test whether such self-interest also weakens the accountability mechanisms documented in the first part of the paper. Specifically, we compare the reduction in fires driven by accountability before and after a constituency switches to electing an agricultural politician. When non-agricultural politicians remain in office, fires decline when politicians' own constituents are polluted, as expected. In contrast, after constituencies switch to agricultural politicians, crop burning increases even where accountability incentives to reduce it are present. This pattern indicates that, as with lobbying,  self-interest can override the pressure of political accountability, even when citizens directly experience the consequences.

%%By interacting the incentives that trigger political responsiveness with those that amplify  inaction, we can observe when and how accountability erodes, giving way to competing interests.
%paraphrase: Understanding these dynamics is crucial for designing institutions that can channel democratic accountability towards durable environmental improvements. 

\todo{Saad/Gemma write new mechanisms here}

%Contributions
\paragraph{Contributions:}

This paper makes several contributions. First, we contribute to the literature on how governments regulate the environment, where previous research has tended to focuse on the bureaucracy.\footnote{See, for instance, \citet{burgess2012political, lipscomb2016decentralization, shen2022political, dipoppa2024bureaucrat}.} We expand this literature by centering the role of politicians: we ask whether and under what conditions politicians in democracies translate citizen interests into better environmental regulation. This contrasts with most studies that portray politicians as obstacles to, not agents of, environmental protection.\footnote{For instance, scholars have shown that campaign donations and lobbying can depress pro-environment votes and legal enforcement \citep{mcalexander2020elections, stokes2020short, harding2024buying, gonzalez2024money}, elections can exacerbate environmental harms \citep{cao2019environmental, balboni2021cycles, sanford2023democratization} (but c.f. \citet{list2006elections}), political control of bureaucratic assignments weakens implementation \citep{xu2025bureaucratic}, and parties in power can offload conservation costs onto their opponents \citep{mangonnet2022playing}. Even informational nudges about voter preferences often fail to lead politicians to adopt pro-environmental policies \citep{hsiao2025climate}.} In doing so, we join very recent work that has begun to probe these questions in the context of political selection by asking if certain politician identity, such as gender and ethnic minorities, improves environmental governance \citep{mavisakalyan2019gender, baragwanath2024greener, gagliarducci2019hurricanes,bragancca2022cutting, gulzar2024representation, jagnani2021women}. In contrast to the role of identity, we focus on accountability incentives and document its promise as one of the few levers that can improve environmental governance.

The second strand of literature we contribute to is on the politics of air pollution. Despite its scale and urgency, air pollution has received little systematic attention in political science.\footnote{See \citet{bar_gulzar2025} and \citet{grossman2026politics} for recent reviews on environment and climate change more broadly} This deficit is striking because transboundary air pollution presents a particularly poignant challenge to democratic accountability: \citet{monogan2017gone} show that polluting plants are more likely to be located at US state boundaries to potentially offset localized pollution impacts, while \citet{LeeVoeten2023} show that international Chinese air pollution shapes blame attribution in South Korea.\footnote{A related literature in economics also studies transboundary pollution in the context of industrial emissions \citep{sigman2002international, kahn2015water, lipscomb2016decentralization}.} Probing if regulatory accountability can be leveraged to control air pollution, scholars have thus far tended to focus on the demand side, examining conditions under which citizens consider air pollution to be a political problem. Instead, we turn attention towards the supply side, showing that politicians are sensitive to transboundary incentives, take action to prevent harm to their own constituents, yet are swayed by weakening sanctions and selection. By doing so, we bring the study of air pollution into the core of political economy and contribute to a broader understanding of when and why politicians act to confront diffuse environmental harms.


Finally, we also make a contribution to the study of environmental politics in South Asia. Nowhere is the problem of air pollution more severe than in South Asia, home to one in four humans, which now bears the world’s highest levels of ambient particulate pollution \citep{pullabhotla2023global}. Scholars across disciplines have started investigating the causes and consequences of crop burning, one of the largest contributors to this crisis.\footnote{For a selection, see: \citet{shyamsundar2019fields, dipoppa2024bureaucrat, garg2023, singh_thachil_2023_airpollution, jack2022money, bhuvaneshwari2019crop}} We provide the first systematic, large-scale analysis of the political accountability incentives that drive both the persistence and reduction of crop burning in Northern India, the region where most of the subcontinent’s air pollution originates and where political action carries profound implications for the lives of millions.
%Most closely related to our focus on politicians, \citet{jagnani2021women} shows that electing women leaders reduces crop burning rates in India.

%%%%%%%%%%%%%%%%%

% =========================================================================
% Section 2: The Management of Air Pollution (Revised)
% =========================================================================
% This file contains Section 2 only. Include it in the main document via
% \input{section2} or paste it in place of the existing Section 2.
% Requires: \usepackage{natbib} and section2_references.bib
% =========================================================================

\section{The Management of Air Pollution}\label{sec:context}
 
 
\subsection{Crop Burning in South Asia}
 
Air pollution is one of the foremost threats to human capital and development worldwide. Exposure to fine particulate matter (PM2.5) causes strokes, chronic respiratory disease, diabetes, and premature mortality, with cascading effects on labor productivity, educational attainment, and hospitalization rates \citep{lelieveld2015contribution, deryugina2019mortality, heo2025international}. The global economic cost of air pollution is estimated at nearly \$2.9 trillion annually, or roughly 3.3 percent of global GDP \citep{myllyvirta2020quantifying}. This burden falls disproportionately on the developing world: South Asia now bears the highest levels of ambient PM2.5 on the planet \citep{pullabhotla2023global}, and recent epidemiological work attributes between 1 and 2 million premature deaths per year in India alone to ambient air pollution \citep{pandey2021health, lancet2024pm25india}.
 
Agricultural emissions, and crop-residue burning in particular, are among the largest contributors to seasonal PM2.5 spikes in South Asia \citep{lelieveld2015contribution, bikkina2019air}. Crop burning is the practice of setting fire to the stubble left in fields after mechanical harvesting to clear the land quickly and cheaply before the next sowing cycle. As farmers themselves acknowledge, burning is a rational response to tight economic constraints: clearing a field with a matchbox costs virtually nothing compared to the thousands of rupees required to rent a mechanized residue-management alternative \citep{jitendra2017, jack2022money}. Yet the aggregate consequences are severe. Each post-monsoon season, millions of fires across Punjab, Haryana, and western Uttar Pradesh generate a vast plume of smoke visible from space, blanketing northern India in smog that persists for weeks \citep{lan2022air, bikkina2019air}.
 
The roots of this crisis trace to the Green Revolution. The introduction of high-yielding rice and wheat varieties, combined with public procurement guarantees and subsidized inputs, entrenched a rice-wheat double-cropping system across the Indo-Gangetic Plain that leaves farmers with as little as ten to fifteen days between the rice harvest and wheat sowing \citep{kumar2015crop, bhuvaneshwari2019crop}. The Punjab Preservation of Subsoil Water Act of 2009, which delayed paddy transplanting until mid-June to conserve groundwater, compressed this window further, increasing farmers' reliance on fire as the fastest residue-clearing method \citep{maue2021climate, tripathi2016subsoilwater}. 
%Prior to the late 1960s, Punjab and Haryana were not major rice-producing regions, and large-scale residue burning was essentially unknown \citep{singh2018great}. The introduction of high-yielding rice and wheat varieties, combined with public procurement guarantees and subsidized inputs, entrenched a rice-wheat cropping system across the Indo-Gangetic Plain that today covers roughly 13.5 million hectares \citep{kumar2022emerging}. This system demands two harvests per year in rapid succession, leaving farmers with as little as ten to fifteen days between the rice harvest and wheat sowing \citep{kumar2015crop, sahai2011assessment}. The adoption of combine harvesters, which leave 20--30 centimeters of standing stubble, further raised the volume of residue and the cost of mechanical removal \citep{grover2017crop, sidhu2015crop}. The problem was compounded by the Punjab Preservation of Subsoil Water Act of 2009, which mandated that paddy transplanting not begin before mid-June in order to conserve groundwater. While successful in slowing aquifer depletion, the law shifted the rice cycle later into the calendar, compressing the post-harvest window even further and increasing farmers' reliance on fire as the fastest residue-clearing method \citep{maue2021climate, tripathi2016subsoilwater}. Finally, sustained out-migration from rural areas has tightened the agricultural labor market, making manual residue collection increasingly impractical and reinforcing the incentive to burn \citep{garg2023, singh2018great}.


\subsection{Political Salience of Crop Burning}\label{ssec:salience}
Crop residue burning is increasingly a politically salient environmental issue, creating electoral incentives for politicians to claim credit for reducing burning or to avoid blame when it persists. In a representative survey of more than 2,200 residents across all districts of Punjab conducted in 2020–21, \citet{yang2025perceptions} asked respondents, ``What do you think about stubble burning in Punjab?'' Nearly two-thirds (65.1\%) described stubble burning as a ``Big problem, and have to stop now.'' \textcolor{red}{In addition, in a survey of 3,116 Delhi residents, \citep{singh_thachil_2023_airpollution} find that citizens blame the government for air pollution more than any other actor, 
and that 75 percent hold the state government of Punjab, where the burning occurs, responsible for solving crop-burning pollution, rather than the government of the jurisdiction where the smoke lands. Together, these two studies indicate broad public concern about air pollution, as well as sophisticated blame attribution.}\todo{Gemma have a look at this}

In Punjab, state-level political discourse tends to center on the scale of subsidies for CRM machinery, the construction of biomass and bio-CNG plants, and the allocation of budgetary resources to stubble management. For example, the Governor of Punjab's address to the state legislature in 2024 highlighted that the government had provided subsidies of 80 percent to farmer groups and 50 percent to individual farmers since 2018, distributing over 117,000 machines, alongside investments in compressed biogas projects designed to absorb surplus paddy straw.\footnote{Address of the Governor of Punjab, March 1, 2024. See also the Punjab Governor's Address of January 2020, which detailed the distribution of 28,609 machines through 5,439 custom hiring centers.} In Haryana's Vidhan Sabha, legislators have highlighted both punitive and incentive-based measures, from per-acre fines imposed by the Pollution Control Board to subsidies for straw management projects and briquette-making units.\footnote{Haryana Vidhan Sabha, Bulletin No.\ 3, September 11, 2018; and the 45th Report of the Committee on the Welfare of SC/ST/OBC, 2021--22.}
 
Importantly, opposition politicians use the same issue to hold incumbents accountable for implementation failures. In Punjab's Vidhan Sabha debates, opposition members have alleged that funds earmarked for stubble management were diverted, and that promised per-acre support never materialized on the ground.\footnote{Punjab Vidhan Sabha Debates, November 29, 2017.} This pattern of credit-claiming by incumbents and blame-attribution by opponents is consistent with a setting in which public opinion creates electoral incentives for responsiveness on environmental governance. 

Crop burning appears regularly in wider public discourse as well: journalists, courts, and civil society groups routinely pressure the central and state governments to demonstrate progress, particularly during the winter smog season when air quality indices reach hazardous levels.
 
Taken together, the institutional architecture of anti-burning policy in India creates a setting in which politicians are important actors. They occupy the implementation node where national mandates meet local realities: they influence the deployment of machinery, the intensity of enforcement, the distribution of subsidies, and the political messaging around environmental governance. This makes the study of their incentives, and the conditions under which those incentives bind or break down, central to understanding the persistence and variation of crop burning across northern India.
 

 
 
\subsection{State Politicians' Role in Curbing Crop Burning}\label{ssec:politicians_role}
 
The governance of crop burning in India involves multiple tiers of government. At the national level, the Commission for Air Quality Management (CAQM), established as a statutory body in 2021, coordinates anti-pollution frameworks across the National Capital Region and adjoining states, combining in-situ residue management support, such as the subsidized distribution of Happy Seeders, mulchers, and balers, with ex-situ programs to channel crop residue toward industrial uses in power plants and brick kilns.\footnote{See the Commission for Air Quality Management in National Capital Region and Adjoining Areas Act, 2021; and \citealp{pib_caqm_2024}.} The CAQM directs states to prepare district-level and block-level action plans targeting hotspot villages, and its 2024 amendment rules substantially increased the environmental compensation penalties for violations.\footnote{See CAQM Direction No.\ 84 (October 10, 2024) and the revised schedule of penalties notified in December 2024.} Nevertheless, because air pollution crosses jurisdictions, effective control requires not just higher-level coordination but local-level implementation that is responsive to political incentives \citep{bar_gulzar2025}.
 
State legislators, Members of the Legislative Assembly (MLAs), are important actors in translating these national frameworks into local enforcement and service delivery. In India's federal system, MLAs elected from single-member assembly constituencies (ACs) exercise considerable influence over district administrations. They shape the allocation and deployment of subsidized crop-residue management (CRM) machinery through custom hiring centers, direct the attention of nodal officers and village-level functionaries toward enforcement and monitoring, and serve as important interlocutors between farmers and the state bureaucracy on agricultural policy. This role is reinforced by the structure of India's anti-burning policy architecture. Since the central government launched the Promotion of Agricultural Mechanisation for In-Situ Management of Crop Residue scheme in 2018, over 100,000 CRM machines have been distributed across Punjab, Haryana, Uttar Pradesh, and Delhi, with subsidies of 50 percent to individual farmers and 80 percent to cooperatives and panchayats.\footnote{Ministry of Agriculture and Farmers Welfare, Government of India. Between 2018-19 and 2025-26, cumulative releases under the scheme totaled approximately Rs.~4,237 crore. See \citealp{pib_crm_2024}{PIB press release, October 2024}.} The utilization of this machinery, however, depends on local-level logistics, such as timely delivery, operator training, equitable access across landholding sizes, and the resolution of disputes over service charges, tasks that are fundamentally shaped by politicians' convening power and constituency service incentives \citep{jack2022money, shyamsundar2019fields}.

Beyond this informal influence, elected representatives also have a formal institutional role in overseeing scheme implementation at the district level. The District Development Coordination and Monitoring Committee (DISHA), constituted by the Ministry of Rural Development since 2016, brings together Members of Parliament, MLAs, and the district bureaucracy in a structured oversight body.\footnote{DISHA superseded the earlier District Vigilance and Monitoring Committees. See \citealp{disha_mord}.} Each district's DISHA committee is chaired by the senior-most Lok Sabha MP, with all MLAs from the district serving as members and the District Collector acting as convenor. Meeting quarterly, the committee reviews the implementation of central government schemes, now covering over 100 programs across 37 ministries, including agricultural mechanization and crop residue management, and has formal powers to seek information, identify bottlenecks, and demand follow-up action from district officials.\footnote{As of 2024, DISHA chairpersons have been appointed in 689 of 718 districts. See \citealp{pib_disha_2018}.} This institutional arrangement gives MLAs a formal channel, alongside their informal leverage, through which they can press district bureaucracies on CRM machinery utilization, anti-burning enforcement, and alternative residue-management programs, or conversely, signal leniency when facing pressure from farmer constituents.

This combination of formal and informal influence extends to enforcement. District magistrates and pollution control boards impose environmental compensation penalties on farmers who burn residue, ranging from Rs.~2,500 to Rs.~15,000 per incident depending on landholding size.\footnote{See the Haryana Vidhan Sabha Bulletin No.\ 3 (September 11, 2018) for penalty schedules. \textt{https://haryanaassembly.gov.in/wp-content/uploads/2019/01/Bulletin-No.-3-119181st-Sitting.pdf}} MLAs can intensify or relax the pressure behind this enforcement: they may direct district officials to step up patrols or put pressure on state level bodies to use remote-sensing dashboards to target hotspot villages during peak burning season. Alternatively, they may signal tolerance for burning in response to farmer grievances. \textcolor{blue}{Enforcement priorities can be set at the start of the burning season and adjusted as conditions evolve, making political action responsive on a time-sensitive scale.}
This discretion makes politicians a critical node in the implementation chain, both potentially enabling and constraining the effectiveness of national policy.
 
 



\begin{comment}
    

\section{OLD The Management of Air Pollution}\label{sec:context}



\subsection{Crop Burning in South Asia}

Air pollution is increasingly a global challenge with important consequences for human capital and development. High levels of PM particles have been shown to cause strokes, diabetes, and respiratory diseases, leading to lower work productivity and educational attainment as well as to more hospitalizations and premature deaths \citep{lelieveld2015contribution, deryugina2019mortality,heo2025international}. The cost of poor air for governments and society is immense and estimated at 2.9 trillion US dollars, corresponding to 3.3 percent of the global GDP \citep{myllyvirta2020quantifying}. 

This burden is not equally distributed across countries: levels of particulate matter in the air are highest in the developing world, particularly in South Asia. One of the largest contributors to air pollution in this setting is PM2.5 from agricultural emissions \citep{lelieveld2015contribution}, and in particular crop-residue burning. Used as a fast and inexpensive method for clearing crop residue after harvest and before the next cropping cycle, crop burning has costs that far outweigh its convenience \citep{jack2022money}. Farmers echo this sentiment: ``If I can clear my farm using a one-rupee matchbox, why will I spend thousands?'' \citep{jitendra2017}, where the latter is referring to alternative methods of clearing stubble such as renting mechanized harvesters.

In India, crop residue burning is used extensively.  Several factors have contributed to the severity of the problem. First, the Green Revolution enabled farmers in in the northern sub-continent to cultivate multiple crops every year. This practice is called the `rice-wheat cropping system' \citep[p. 150]{singh2018great}, and is characterized by rapid harvest and sowing cycles so that crops can catch prime conditions for their growth. Given that India has a large agricultural sector that focuses on rice and wheat production, which generate the largest crop residues (\href{https://agricoop.nic.in/sites/default/files/NPMCR_1.pdf}{NPMCR, 2014}), the problem of stubble burning is acute. Second, crop burning has increased with mechanization technologies, which produce larger crop residues and increase incentives to use fire to clear the fields \citep{maue2021climate}. An increase of labor wage due to sustained out-migration from rural areas has further increased the use of these mechanized tools (see \citet{garg2023} and \cite[p. 156-164]{singh2018great} for more details).


\subsection{Constituency Politicians' Role in Curbing Crop Burning}

The management of crop burning is regulated by national rules and state schemes, but ultimately its implementation is based on constituency-level politicians. Nationally, the Commission for Air Quality Management (CAQM) coordinates a common framework across states, combining in-situ support, such as the provision of machinery to remove the stubble and bio-decomposers, with ex-situ linkages to power. States are directed to prepare district- and block-level action plans and to enforce them through local administrations \href{https://www.pib.gov.in/PressReleaseIframePage.aspx?PRID=2076918}{(PIB/CAQM press release)}. This multi-level action is necessary as, given that air pollution crosses jurisdictions, control requires both higher-level coordination \emph{and} effective local action \href{https://documents1.worldbank.org/curated/en/701891468285328404/pdf/890650PUB0Clea00Box385269B00PUBLIC0.pdf}{(World Bank, 2014)}.

State politicians are fundamental in the control of crop-burning. First, they can exert pressure on district bureaucrats for pulling levers that matter for implementation, such as the allocation of CRM machines, staffing, and direction of nodal officers, village-level awareness and monitoring campaigns, and rapid response in hotspot villages \href{https://www.pib.gov.in/PressReleaseIframePage.aspx?PRID=2076918}{(PIB/CAQM press release)}. Second, programs hinge on constituency-facing delivery: since 2018, large volumes of machines have been distributed, and their utilization depends on local mobilization, training, logistics, and accountability, classic tasks where elected representatives shape outcomes via convening power and constituent services \href{https://www.pib.gov.in/PressReleaseIframePage.aspx?PRID=2076918}{(PIB/CAQM press release)}. Finally, action plans target specific hotspot villages and enforce harvester rules and EC penalties, steps that require day-to-day political attention to work \href{https://www.pib.gov.in/PressReleaseIframePage.aspx?PRID=2076918}{(PIB/CAQM press release)}\todo{We cite the same thing multiple times in these two paragraphs.. need to diversify}

Elected officials face incentives for action from a mix of electoral accountability at the constituency level, political credibility more broadly, and distributive politics. First, severe post-harvest smoke is visible and salient; when fire counts spike, constituents and courts blame governments, while documented declines are readily claimed as achievements \href{https://www.pib.gov.in/PressReleaseIframePage.aspx?PRID=2076918}{(PIB/CAQM press release)}. Second, CAQM’s directions and the strengthened EC schedule raise reputational and administrative costs for non-enforcement, giving politicians a hard incentive to be seen as enforcing rules and fixing delivery gaps \href{https://www.pib.gov.in/PressReleaseIframePage.aspx?PRID=2076918}{(PIB/CAQM press release)}. Third, subsidies, CHCs, and ex-situ market linkages generate tangible benefits (machines, contracts, straw purchase), which legislators can channel to local groups, building coalitions for alternatives to burning \href{https://www.pib.gov.in/PressReleaseIframePage.aspx?PRID=2076918}{(PIB/CAQM press release)}.


\paragraph{Legal and policy backdrop}
Stubble burning has been proscribed under air-pollution statutes for years; in 2021, Parliament created the CAQM with a regional mandate for the NCR and adjoining areas, emphasizing coordination and enforcement over criminalization. The 2024 CAQM Amendment Rules revised (and effectively doubled) EC penalties and mandated their immediate imposition by states—clarifying that while burning is handled administratively rather than criminally, the financial and administrative consequences have been strengthened \href{https://www.pib.gov.in/PressReleaseIframePage.aspx?PRID=2076918}{(PIB/CAQM press release)}.


State politicians are pivotal because they turn national frameworks into local-level actions. Their incentives are now sharper—courts, CAQM schedules, visible health harms, and distributive benefits all push toward prevention. And the tool-kit is concrete: equip and supervise CRM use, enforce Super SMS and EC fairly, stand up ex-situ straw markets, surge resources to hotspots, and coordinate across borders on an airshed calendar. Where these pieces move together, fire counts \emph{do} fall \href{https://www.pib.gov.in/PressReleaseIframePage.aspx?PRID=2076918}{(PIB/CAQM press release)}.

\paragraph{State politicians' actions on crop burning}

In practice, MLA politicians can take a range of concrete actions to reduce crop burning, often via district coordination meetings and field visits:

\begin{itemize}
\item \textbf{Deploy and monitor CRM capacity:} push timely delivery and equitable mapping of Happy Seeders, Super SMS, mulchers, balers/rakes; ensure CHCs are functional, operators trained, and service prices posted; redirect idle inventory to hotspot villages \href{https://www.pib.gov.in/PressReleaseIframePage.aspx?PRID=2076918}{(PIB/CAQM press release)}.
\item \textbf{Enforce rules fairly:} back the district administration and police to enforce Super SMS mandates on harvesters and apply EC penalties proportionately; support grievance redress to prevent rent-seeking in fine collection and to keep enforcement legitimate \href{https://www.pib.gov.in/PressReleaseIframePage.aspx?PRID=2076918}{(PIB/CAQM press release)}.
\item \textbf{Build ex-situ markets:} convene straw buyers (power/CBG plants, industrial boilers, brick kilns) and farmer groups to set up purchase schedules, aggregation points, and transport; use state/central schemes for viability-gap funding and supply-chain equipment \href{https://www.pib.gov.in/PressReleaseIframePage.aspx?PRID=2076918}{(PIB/CAQM press release)}.
\item \textbf{Target hotspots:} resource additional patrolling, remote-sensing feed review, and last-mile outreach in high-incidence villages; hold nodal officers to performance with transparent dashboards \href{https://www.pib.gov.in/PressReleaseIframePage.aspx?PRID=2076918}{(PIB/CAQM press release)}.
\item \textbf{Promote agronomic shifts:} amplify short-duration paddy varieties to widen the harvest window; expand DSR incentives to relax turnaround pressure; run farmer field-schools on soil/yield benefits of non-burning \href{https://www.pib.gov.in/PressReleaseIframePage.aspx?PRID=2076918}{(PIB/CAQM press release)}.
\item \textbf{Coordinate airshed responses:} align burn-prevention schedules with neighboring constituencies for co-timed enforcement and straw off-take, consistent with airshed management principles \href{https://documents1.worldbank.org/curated/en/701891468285328404/pdf/890650PUB0Clea00Box385269B00PUBLIC0.pdf}{(World Bank)}.
\end{itemize}


\paragraph{State politicians discourse on crop burning:} Politicians  also advertise their action in public discourse, which indicates there are political incentives to claim credit for environmental action. Politicians' discourse on crop burning is tied to the instruments of implementation discussed above. We document politicians' discourse via tweets as well as statements made by them in the state assembly, the Vidhan Sabha.

Politicians in Punjab, India, for example, tweeted about their action to distribute crop-residue management machines through custom hiring centers and to individual farmers, the scale of the subsidies attached to these programs, and the development of biomass, bio-CNG, and bio-ethanol projects intended to create off-take for paddy straw. 

\begin{quote}
\emph{“My Government has been providing 80\% subsidy to farmer groups/panchayats and 50\% subsidy to individual farmers... From FY 2018–19 to FY 2022–23 a total of 1{,}17{,}000 machines has been provided for stubble management... [and] Compressed Biogas (CBG) projects... shall utilize surplus paddy straw to the tune of 1.87 million tons per annum.”} \href{2519.pdf}{(Address of the Governor of Punjab, Mar.\ 1, 2024)}
\end{quote}

\textbf{i would keep only first and drop next 3}

\begin{quote}
\emph{“Under the In-Situ Crop Residue Management scheme... 28{,}609 residue management machines have been distributed through 5{,}439 custom hiring centers at 80\% subsidy and 12{,}075 to individual farmers at 50\% subsidy... approximately 24{,}000 machines [in 2019–20] have been provided to the farmers at subsidized rates.”} \href{1213.pdf}{(Punjab Governor’s Address, Jan.\ 2020; Vol.\ X—Special Session—No.\ 1)}
\end{quote}
\begin{quote}
\emph{“Presently, 10 Biomass power projects of 7350 MWs are in operation and two paddy straw based biomass power projects of 36 MWs have recently been commissioned to utilize the paddy stubble... Private investment is being facilitated for setting up Bio-CNG/Bio-ethanol projects based on paddy stubble.”} \href{1213.pdf}{(Punjab Governor’s Address, Jan.\ 2020; Vol.\ X—Special Session—No.\ 1)}
\end{quote}
\begin{quote}
\emph{“Checking Stubble Burning Practices... this Government would only undertake positive targeted actions for controlling it… An allocation of ₹200 crore is hereby proposed in this budget.”} \href{2367.pdf}{(Punjab Budget Speech, 2022–23, Finance Minister Harpal Singh Cheema)}
\end{quote}


In Haryana, politicians emphasized both carrots and sticks: one statement in the Vidhan Sabha referred to the size of penalties imposed on farmers who burned residue, while another stressed subsidies for residue-management projects. 

\begin{quote}
\emph{“The Pollution Control Board has imposed penalty @ ₹3{,}000 per acre on those farmers, who have burnt the husk (Parali)...”} \href{Bulletin-No.-3-119181st-Sitting.pdf}{(Haryana Vidhan Sabha, Bulletin No.\ 3, Sept.\ 11, 2018)}
\end{quote}
\begin{quote}
\emph{“It has been also proposed to provide subsidy under state scheme (SB-89) to straw management projects like briquette/pellet/fodder making @ 50\% subsidy (capped at ₹50 lakhs)... for establishment of CHC/Crop Residue Management Projects.”} \href{45th-Report-2021-22.pdf}{(Haryana Vidhan Sabha, 45th Report of the Committee on the Welfare of SC/ST/OBC, 2021–22)}
\end{quote}

Opposition politicians used the same issue to criticize incumbents for failures of implementation, arguing that funds allocated for stubble management had been diverted and that promised support had not materialized on the ground. 


\begin{quote}
\emph{“The government spent ₹97 crore for stubble management, but that money was diverted to other issues… They talk about allocating ₹5{,}000 per acre for stubble management, but nothing concrete is happening.”} \href{460.pdf}{(Punjab Vidhan Sabha Debates, Nov.\ 29, 2017)}
\end{quote}

Crop burning also appears in wider public discourse as an issue on which governments are expected to respond visibly, as illustrated by Vishnu Som’s appeal to the Prime Minister to end the ``annual nightmare'' of air pollution.

\begin{quote}
\emph{``
@PMOIndia, Its time we ended the annual nightmare on air pollution.  The loss of life is now quantified - this cannot be a scenario that repeats year after year.  We need technology solutions to stubble burning and need to end the battle between the state and centre.  Please help us.''
(\href{https://x.com/vishnundtv/status/1991143509156598228?s=46&t=zmR0TSx-gxK3_KXZkFwBJQ}{Vishnu Som, November 19th, 2025, X}
})
\end{quote}

Taken together, these examples indicate that crop burning is politically salient: politicians seek credit for subsidies, machines, and alternative-use projects, and they are exposed to blame when implementation is perceived as inadequate. This is consistent with a setting in which public opinion creates incentives for responsiveness in environmental governance.



%The management of air pollution involves all levels of government. At the national and state level, goverments in India have launched a series of initiatives to lead farmers to switch to alternative methods than burning: schemes to promote in-situ management technologies, such as happy seeders, mulchers, rotavators, or financial support for farmers to adopt alternatives, as well as campaigns to raise awareness about the consequences of burning. Importantly, governments also made the practice illegal.\footnote{In response to the farmers' protests, crop burning has been decriminalized in December 2021, such that it is not a crime anymore. However, fines are still imposed for crop burning and they have been doubled in 2024. The bulk of our data overlap with a time in which the practice was a crime.} Implementation of both bans and subsidies involves local administrations. Politicians frame these actions as part of broader commitments to reduce air pollution and protect public health


%\textcolor{blue}{can use these sources: \url{https://www.pib.gov.in/PressReleaseIframePage.aspx?PRID=2076918&utm_source=chatgpt.com} and \url{https://documents1.worldbank.org/curated/en/701891468285328404/pdf/890650PUB0Clea00Box385269B00PUBLIC0.pdf?utm_source=chatgpt.com}}

% Summary of the Delimitation Commission (Generated using AI assistance)
\end{comment}


\section{Politicians' Incentives to Control Fires}

Our first objective is to establish whether  accountability incentives drive politicians' action to control environmental damage. 
In this section, we describe the data, empirical strategy, and results showing that responsiveness to constituents' exposure to pollution shapes politicians’ behavior to control crop related fires.

%We have in Punjab 234 ACs. 486 Bihar. 806 UP. Haryana 90 ACs
\subsection{Data\label{ssec:data1}}

We assemble data on the identity and characteristics of the politicians in charge of each constituency in the sample, as well as on the behavior of farmers, and spatially merge these with information on fires and wind direction. We begin by describing the political data.

\paragraph{Assembly Constituencies:} We focus on the geographical jurisdiction of the State Legislative Assembly constituency, defined according to the 2008 delimitation, which was implemented under the Delimitation Act, 2002 \citep{delimitation_act_2002} and
based on the 2001 Census. We focus on the states where crop burning is a major problem ---Punjab, Haryana, Uttar Pradesh, and Bihar--- restricting the sample to rural areas by dropping urban grids. We classify a grid as urban if it contains at least one pixel labeled ``Urban" or ``Dense Urban Cluster" according to the GHS-SMOD R2023A dataset \citep{ghs2023smodR2023}, which classifies
$1km^2$ pixels according to population size and density thresholds to capture the full settlement
hierarchy. Constituencies that consist entirely of urban grids are dropped from the sample, leaving 808 of the 853 assembly constituencies. 

%These representatives hold significant authority over local governance, including implementing measures to address air pollution from crop residue burning and directing the allocation of development funds and public services. 
%\footnote{These are the same states considered in other papers focusing on air pollution in South Asia, for example, \citet{liu2020crop, dipoppa2024bureaucrat}.} 


\paragraph{Politicians' Incentives:} We proxy politicians' incentives to curb fires using variation in wind direction, which determines whether smoke from fires is likely to affect politicians' own constituency or a neighboring one. Wind direction and speed data come from the \citet{munozsabater2019era5landmonthly}, which is precise at a 25x25 km resolution. To approximate what politicians might expect for a given month–year, we compute ten-year rolling averages of wind direction using data from the early 2000s onward. Figure \ref{fig:overtime} (right panel) shows that actual wind directions and their rolling averages closely track each other, substantiating the expectations that politicians are aware of prevalent wind directions in a season. In this setting, indeed, the two major seasonal wind shifts occur in June and August each year. Section~\ref{ssec:empirical1} explains in more detail how we define political incentives by combining data on wind direction and constituency shapes.

\paragraph{Fires:} Our main outcome of interest is the presence and extent of crop stubble burning. We measure crop burning with data on fires from Nasa satellites Moderate Resolution Imaging Spectroradiometer (MODIS, 1km x 1km resolution) and Visible Infrared Imaging Radiometer Suite (VIIRS, 375m x 375m resolution) from September 2012, the first year in which VIIRS starts operating, to August 2022. The satellite records fires from thermal signatures, such that the presence of cloud cover or smoke in the atmosphere does not prevent its detection. The absence of forests in the area of India we study means that we can safely attribute fires to crop burning activity rather than to forest fires. We plot the total number of fires over time in Figure~\ref{fig:overtime} (left panel): each year, there are two peaks in the months of April-May and October-November, in correspondence with the crop burning seasons of wheat and rice respectively \citep{singh2014crop}.\footnote{Data on fires and wind are taken from \citet{dipoppa2024bureaucrat}.}

\paragraph{Crops:} Rice is the crop responsible for the largest share of stubble burning, followed—though to a lesser extent—by wheat. To focus on areas most affected by crop burning, we collect data on average crop area and production in 2010, the latest year with micro-level crop data available before the start of our analysis period. The data is sourced from the \citet{fao_maps}, it has a resolution of 5 x 5 arc/min, corresponding to approximately twice the size of a grid in the final database, and is plotted in Figure~\ref{fig:map_crops}.

%\textcolor{red}{Gemma: i propose to move the next paragraph to the appendix and only leave in the part about dropping urban and giving total n obs.}
%\paragraph{Combining all the data at a 5 square km grid-cell level:}
%We combine all the data into a single dataset at the $5x5km$ grid-cell level. Wind and fire grids do not overlap precisely: to reduce errors in the assignment of fires to wind grids, we split the $25km$x$25km$ wind grid-cells into squares of $5x5k$ each. This guarantees that fires are assigned to the wind grid-cell which contains most of the fire grid-cell. Splitting wind grids in smaller parts also allows us to address another concern: some grids might overlap across jurisdiction boundaries. By reducing their size, we minimize the amount of misattribution and we are able to assign each grid to the appropriate jurisdiction. Any remaining error should bias our findings toward zero, as errors are uncorrelated with the dependent and independent variables. In case of overlap, we assign a grid to the jurisdiction that contains a larger area of the grid. For crop data, we assign the pre-period crop production to each $5km^2$ grid-cell to classify what crops are produced in each location.  To exclude areas in which crop burning could not occur, we drop urban grid-cells using the GHS-SMOD dataset \citep{schiavina2026ghssmod}, which classifies $1km^2$ pixels according population size and density thresholds to capture the full settlement hierarchy. We define a grid-cell as urban if there is at least on pixel within it classified as Urban Center or Dense Urban Cluster.
Our full sample includes 76,800 grid-cells which we observe in each month and year from September 2012 to August 2022, for a total of 9,216,000 observations. We map each constituency in the sample with an illustration of the gridded data in Figure~\ref{fig:gridsize}.



% \begin{figure}[htbp!]   
% \caption{\textbf{How Changes in Wind Direction Affect Political Calculus and Crop Burning}\label{fig:empirics_illustration1}}
% \centering\subfloat[Gridded Dataset]{\includegraphics[width=.3\textwidth]{figures/map_grids.png}\label{fig:gridsize} }\hspace{.05\textwidth} \centering\subfloat[Treatment Definition and Variation Example]{\includegraphics[width=.18\textwidth]{figures/panel_A_downwind.png}
%     \includegraphics[width=.18\textwidth]{figures/panel_C_upwind.png}
%     \includegraphics[width=.18\textwidth]{figures/panel_B_downwind.png} \label{fig:example_downwind}} \\
% %  \setcounter{subfigure}{1}
% \centering\subfloat[Variation in Treated Grids Over Time Across the Sample]{\includegraphics[width=.45\textwidth]{figures/march2013_plot.png}\label{fig:map_treat}
%     \includegraphics[width=.45\textwidth]{figures/october2013_plot.png}} 
%   %%%% Eventstudy %%%%%%%%
% \hspace{\fill}
% \begin{minipage}{\textwidth}
%     {\textit{Notes: } \textcolor{red}{Remove the districts borders, only states and highlight only one district. Change the aesthetic of the figures in general} Figure \ref{fig:gridsize} illustrates the 5km$^2$ gridded dataset in green with constituency boundaries in blue. It shows the Bachhrawan assembly constituency magnified as an example. Figure \ref{fig:example_downwind} illustrates the coding of treatment. The blue arrow represents wind direction in a given month-year, black dots represent population. In the left panel, the population in the  downwind area from a potential fire is larger than the population in the upwind area from the fire. We define this grid as treated. In the central panel, the area downwind from the fire is larger than the upwind are, but the population in the downwind area is smaller: we define this grid as control. In the right panel, a change in wind direction converts the same grid into the treatment group because a potential crop burning would now become a source of pollution for the majority population of the home constituency. 
%     Figure \ref{fig:map_treat} displays the spatial distribution of this treatment and its temporal variation in two different months. In dark blue shading are treated grids and in light blue are control grids.   
%         }
% \end{minipage}
% \end{figure}

\begin{figure}[htbp]
    \centering
    \caption{\textbf{How Changes in Wind Direction Affect Political Calculus and Crop Burning}}
    \label{fig:empirics_illustration1}

    % ---------- Row 1: subfigure (a) centered alone ----------
    \subcaptionbox{Gridded Dataset\label{fig:gridsize}}[0.45\textwidth]{%
        \includegraphics[width=.78\linewidth]{figures/map_grids.png}%
    }

    \vspace{0.5em}

    % ---------- Row 2: subfigures (b) ----------
    \subcaptionbox{Treatment Definition and Variation Example\label{fig:example_downwind}}[0.75\textwidth]{%
    \includegraphics[width=.32\linewidth,valign=t]{figures/panel_A_downwind.png}%
    \includegraphics[width=.32\linewidth,valign=t]{figures/panel_C_upwind.png}%
    \includegraphics[width=.32\linewidth,valign=t]{figures/panel_B_downwind.png}%
}%

    \vspace{0.5em}

    % ---------- Row 3: subfigures (c)  ----------
    
    \subcaptionbox{Variation in Treatment Over Time Across Units\label{fig:map_treat}}[0.75\textwidth]{%
        \includegraphics[width=.49\linewidth]{figures/march2013_plot.png}%
        \includegraphics[width=.49\linewidth]{figures/october2013_plot.png}%
    }

    \begin{minipage}{\textwidth}
        {\textit{Notes: } \textit{(a)} The map illustrates the 5km $\times$ 5km gridded dataset in green with constituency boundaries in blue. It shows the Bachhrawan assembly constituency magnified as an example. 
        \textit{(b)} The figures illustrate the coding of treatment. The blue arrow represents wind direction in a given month-year, and the numbers represent population in a grid. In the left panel, the population in the downwind area from a potential crop burning is larger than the population in the upwind area: we define this grid as treated. In the central panel, the downwind area is larger than the upwind area, but the population in the downwind area is smaller: we define this grid as control. In the right panel, a change in wind direction converts the same grid into the treatment group because a potential crop burning would now pollute the majority of the constituency's population. 
        \textit{(c)} The maps display the spatial distribution of this treatment and its temporal variation in two different months. Treated grids are shaded in dark blue and control grids in light blue. An image covering all periods of analysis is located in \href{https://drive.google.com/file/d/1_3vuvipO7Us2vbe2WF1VRdrgGGI9jsbP/view?usp=share_link}{this link}.}
    \end{minipage}
\end{figure}

\subsection{Empirical Strategy}\label{ssec:empirical1}


\paragraph{Treatment Definition:}

Our treatment is designed to capture a specific accountability incentive. A politician has stronger reason to prevent crop burning in a location when the resulting smoke is likely to settle over more of their own constituents rather than blow into a neighboring constituency. Whether a given location does so depends on four things: where the location sits within the constituency, how population is distributed inside the constituency, the shape of the constituency, and which way the wind is blowing.

We operationalize this at the grid-cell level. For each 5km$^2$ grid-cell centroid, we draw a line through the centroid perpendicular to the prevailing wind direction vector for that month. This line partitions the constituency into a downwind area, which would receive smoke from a fire in that grid-cell, and an upwind area, which would not. We then compute the constituency population living in each area. 

We define a grid-cell as \textit{treated} in a given month-year if the constituency population downwind of the grid-cell exceeds the constituency population upwind of it: a fire in a treated cell would impose pollution costs on most of the politician's own voters, while a fire in an untreated cell would carry most of its smoke past them.\footnote{We show robustness to other definitions of treatment that are continuous or otherwise transformed in Appendix.}

For the prevailing wind direction we use ten-year rolling averages by month rather than realized contemporaneous wind. Seasonal wind in northern India is governed by the monsoon system, with two major directional reversals occurring at approximately known points in the calendar. Politicians can therefore anticipate, before the burning season, which parts of their constituency will bear the pollution costs of a fire in a given month. The monthly rolling average is intended to approximate that expectation, not to measure the wind on any particular day.

Because wind direction shifts across the calendar, the same grid-cell can be treated in some months and untreated in others: a fire there would pollute its own constituency's voters in one month and a neighbor's in another. This within-grid-cell, across-month switching is the source of identifying variation for our first set of analyses.

We present an example of treatment definition and change in status for a grid in our sample in Figure \ref{fig:empirics_illustration1}b, which shows how the downwind by upwind population fractions change across different grids and months under the maps. In Figure \ref{fig:empirics_illustration1}c we also illustrate the spatial distribution of treatment in the full sample over two different months. 

\paragraph{Estimation Strategy:}
Using this treatment definition to capture politicians' accountability incentives, we estimate a stacked-by-event difference-in-differences model for the impact of a grid-cell switching to polluting most of its jurisdiction's population on fires, leveraging within-grid changes in treatment status over time \citep{cengiz2019effect}. 
We build a separate dataset for each cohort, defined by month of treatment, restricted to a window of 5 months before to 6 months after the switch. In each cohort, treated units are those that begin treatment in that month; controls units are grids never treated during the sample window, together with units not yet treated at the switching month, included only for the periods in which they are in the control condition. We then stack each cohort-dataset into a pooled sample. 
For a 5km$^2$ grid-cell $i$ in assembly constituency $j$, month-year $t$, and cohort $c$, we estimate:
\begin{equation}\label{eq1}
 \text{Fires }_{ict}=\beta \text{{Down $>$ Up}}_{ict}+ \alpha_{ic}+ \gamma_{jct} + \delta X_{it} + \epsilon_{ict}
\end{equation}

where our dependent variable is the number of fires in a gridcell month-year. The coefficient of interest, $\beta$, captures the effect of treatment, as described above, on the number of fires in that grid. 
The high-dimensionality of our data allows for the inclusion of fine-grained fixed effects. First, $\alpha_{ic}$ are grid-cell-cohort fixed effects that absorb time-invariant factors for over 76,800 grid-cells in our data, such as crop suitability or economic characteristics of small locations.  Second,  $\gamma_{jct}$ are jurisdiction by month-year-cohort fixed effects that account for time-varying shocks at the constituency level. $X_{it}$ is a vector of grid-level controls that includes actual wind direction and average wind speed. Given the stacked nature of our data, we interact each of the fixed effects with cohort fixed effects. Standard errors are clustered at the grid-cell and constituency by month-year level. The first accounts for temporal dependence of observations at the treatment status level (grid-cell). The second accounts for potential spatial correlation in wind patterns. 

A typical concern in the standard difference-in-differences setting is that treatment might be assigned based on pre-existing characteristics of the unit endogenous to the outcome: for example, poorer areas might be more likely to receive welfare interventions. Our setting has instead a unique feature: treatment, measuring which locations is more likely to be a source of pollution for a jurisdiction, is assigned based on changes in wind direction interacted with the shape of the assembly constituency polygon and its population, a quantity which is plausibly exogenous. This feature makes the parallel trends assumption  ---that treated grid-cells would have displayed a similar trend in fires as the control grid-cells had the wind not shifted, exposing a larger area of the jurisdiction to smoke--- considerably more credible, as treatment is assigned quasi-randomly. 

We nonetheless provide evidence consistent with the parallel trends assumption with a stacked-by-event event study. For relative time $r$ from the switch to treatment, $Switch$, we estimate:
\begin{equation}\label{eq2}
 \text{Fires }_{ict}= \sum_{r=-n}^{N} \zeta_r Switch_{ir} + \alpha_{ic}+ \gamma_{jtc}+ \delta X_{it} + \epsilon_{ict}
\end{equation}
using the same cohort structure and fixed effects as the specification above, where $Switch$ is a dummy that takes value one when unit $i$ is $r$ periods away from its treatment switch, and zero otherwise. 
Our coefficients of interest, $\zeta_r$, represents the effect of treatment status in each relative month before and after treatment, where the month prior to the switch is the omitted category. 


\subsection{Results\label{ssec:results1}}

Figure \ref{fig:es_main_rotated} plots the event study coefficients. Before switching to polluting more of the constituency, grids that will switch experience a similar number of fires as others. From the month they switch to being a source of exposure for constituents, however, the number of fires significantly drops. This effect correspond to a 13\% reduction in fires with respect to baseline. The result is mainly driven by rice-producing areas where we observe 33\% less fires with respect to baseline (Figure \ref{fig:es_rice}), confirming the main result that politicians are significantly more likely to reduce fires when pollution would affect their own constituency. 


Table \ref{tab:main_DID} reports the results. When wind is expected to blow smoke into a grid's assembly constituency rather than away from it, we observe 5.435 fewer fires per 1,000 in each grid-month in our preferred specification (column 4), corresponding to a reduction of approximately 5.9\% relative to the average number of fires when a grid is not a source of exposure.


 These results are consistent with two interpretations, both of which we classify under the broader umbrella of political accountability. First, it could reflect politicians enforcement effort toward the locations whose smoke would fall on their own voters, reflecting either constituents' direct demand or  politicians' own calculus of citizen welfare. Another possibility is that farmers anticipate that enforcement will be concentrated where burning would pollute the constituency, and do not burn as a result. This is a deterrence channel represents an accountability incentive because the latter is what generates variation in behavior.



To give a sense of the public health stakes of this political action, we translate the 13 percent fire reduction into air quality and mortality estimates using atmospheric modeling from \citet{lan2022air}\todo{Change the numbers}. During the post-monsoon burning season, when rice-burning effects are concentrated, crop residue burning is estimated to contribute 26 $\mu$g m$^{-3}$ to the mean population-weighted PM2.5 concentrations in India. A 13 percent reduction in fires corresponds to a reduction of approximately 3 $\mu$g m$^{-3}$ in burning-season ambient PM2.5 at the national level, and 10-29 $\mu$g m$^{-3}$ in the northwestern constituencies where burning is concentrated. In mortality terms, applying this proportional reduction to Lan et al.'s estimate that Punjab, Haryana, and Uttar Pradesh collectively account for approximately 70 percent of the 69,000 annual premature deaths attributable to crop burning in India would yield roughly 3,100 fewer premature deaths per year. These figures should be interpreted with the caveat that the estimated fire reduction may in part reflect spatial displacement across constituency lines rather than a net decrease in total emissions: if enforcement effort in treated areas is offset by reduced effort elsewhere in the constituency, aggregate pollution and mortality benefits would be correspondingly smaller. Nonetheless, this estimate provides a useful benchmark for the scale of impact that would be achievable if politicians were incentivized to act uniformly as they do in the places and times where their accountability incentives are heightened. Appendix~\ref{sec:mortality_estimates} details the derivation of these figures.

%presents the results when we do not adjust for the pre-treatment linear trends in the post period. The estimated coefficients still show a negative effect of pollution visibility on the number of fires after the switch. The pre-treatment coefficients indicate a trend opposite to the treatment effect, with fires increasing before the switch. Following \citet{rambachan2023more}, we test the sensitivity of our results to pre-trend violations. Results in Figure \ref{fig:es_main_rotated_panel}(c) show that the post-treatment effect we detect is robust to violations of parallel trends up to 0.5 the size of the largest pre-trend violation.




%We conduct an event study analysis around the time in which upwind grids switch to being downwind, using units that never switch and have not yet switched as the control group. %We exclude always-treated observations from the sample because comparisons involving them, i.e. always-treated versus early-treated, are not meaningful and often lead to negative weights in standard DiD estimators with staggered treatment adoption. We include grid and assembly constituency $\times$ month-year fixed effects as in column (4) of Table \ref{tab:main}.
%
%The result of the specification of Equation \ref{eq2} is presented in Figure \ref{fig:combined_es_main}. Given the linear trend in the pre-treatment period, we follow the approach of \citet{rambachan2023more} to test the sensitivity of our results to pre-trend violations. The results show that the pre-treatment trend would have to change by 100\% in the post-treatment period for the estimated effect to be null, \red{which is ?}.


\begin{figure}
    \centering
    \caption{\textbf{Effect of constituency exposure to smoke on number of fires}}\label{fig:main_panel}
    \vspace{-0.5cm}
    \singlespacing

    %%% ---------- Fila 1: (a) izquierda + (c) tabla derecha ---------- %%%
\begin{tabular}{@{}m{0.49\textwidth}@{\hspace{0.02\textwidth}}m{0.48\textwidth}@{}}
  \centering
  \subcaptionbox{Event-study relative to exposure\label{fig:es_main_rotated}}[\linewidth]{%
    \includegraphics[width=\linewidth,height=4.7cm]{figures/stacked_event_study_pop_5pre_rural_1_ori.png}%
  }
  &
  \centering
  \setcounter{subfigure}{2}% el siguiente bloque será (c)
  \subcaptionbox{Effect of constituency and district exposure to smoke on N Fires\label{tab:main_ac_distr}}[\linewidth]{%
    \begin{adjustbox}{width=0.49\textwidth}
      {\def\sym#1{\ifmmode^{#1}\else\(^{#1}\)\fi}
\begin{tabular}{l*{5}{c}}
\hline
 & (1) & (2) & (3) & (4) & (5) \\
 & \multicolumn{5}{c}{Number of Fires (in 1,000 units)} \\ \hline

Down$>$ Up Bureaucrat&     -24.369\sym{***}&     -22.156\sym{***}&      -2.506\sym{**} &      -7.114\sym{***}&     -18.051\sym{***}\\
            &     (2.958)         &     (1.355)         &     (1.018)         &     (1.002)         &     (1.156)         \\
Down$>$ Up Politician (Pop)&       7.819\sym{**} &      -7.902\sym{***}&       4.762\sym{***}&     -10.448\sym{***}&      -6.530\sym{***}\\
            &     (3.110)         &     (2.077)         &     (1.711)         &     (2.059)         &     (2.170)         \\
Down$>$ Up Pol. (Pop) $\times$ Down$>$ Up Bur.&     -21.678\sym{***}&       3.514         &       1.102         &       9.958\sym{***}&       1.616         \\
            &     (3.991)         &     (2.462)         &     (1.885)         &     (2.003)         &     (2.261)         \\
\hline
Observations&  36,758,771         &  36,758,771         &  36,758,771         &  36,758,771         &  36,758,531         \\
N Assembly Constituencies&         808         &         808         &         808         &         808         &         808         \\
N Districts &         160         &         160         &         160         &         160         &         160         \\
Month-Year FE&           N         &           Y         &           N         &           N         &           N         \\
AC FE       &           N         &           Y         &           N         &           N         &           N         \\
AC $\times$ Month-Year FE&           N         &           N         &           Y         &           Y         &           N         \\
District $\times$ Month-Year FE&           N         &           N         &           N         &           N         &           Y         \\
Grid FE     &           N         &           N         &           N         &           Y         &           Y         \\
Mean DV     &      86.353         &      86.353         &      86.353         &      86.353         &      86.353         \\
\hline
\end{tabular}
}
%
    \end{adjustbox}%
  }
  \tabularnewline
\end{tabular}

    \vspace{0.6em}

    %%% ---------- Fila 2: (b) neighbor_output a ancho completo ---------- %%%
    \setcounter{subfigure}{1}% el siguiente bloque será (b)
    \subcaptionbox{Effect of grid distance to the upwind (in blue) and downwind (in red) border on number fires\label{fig:neighbor_analysis}}[\textwidth]{%
      \includegraphics[width=0.55\textwidth]{figures/neighbor_output.pdf}%
    }

    \vspace{-0.3cm}
    \begin{minipage}{\textwidth}
        {\small \textit{Notes: } \textit{(a)} Results from Equation \ref{eq2}.
         % after de-trending for the pre-treatment period linear trend following \citet{dustmann2022reallocation}. 
         The reference category is t=-1, i.e. the month before the switch. We use a stacked diff-in-diff to generate cohorts based on switching date. We restrict the analyses to a 6 months window around the switch to maximize support. The control grids are never treated grids and grids that are upwind around the switching date. 
         %Figure \ref{fig:combined_es_main} reports the not rotated eventstudy as in Equation \ref{eq2} and the sensitivity analysis by \citet{rambachan2023more}. 
          The regression includes grid and assembly constituency $\times$ month-year $\times$ cohort fixed effects and controls for average wind speed and direction. Standard errors are clustered at the grid $\times$ cohort and assembly constituency $\times$ month-year $\times$ cohort level.
          \textit{(b)} (Blue) Effect of being close to an upwind border with another constituency, where smoke would affect the politicians' assembly constituency. (Red) Differential effect (from the upwind border) of being close to a downwind border, where smoke would travel to the neighboring jurisdiction.  All effects are plotted by quintiles of distance to the border, and the 5th quintile, the farthest, is used as omitted category. The regression includes  grid and border $\times$ month-year fixed effects and control for average wind speed and direction. Standard errors are clustered at the grid and assembly constituency $\times$ month-year level.
          %%% Free Riding %%%%%%%
          \textit{(c)} Results from Equation \ref{eq1} with the additional interaction with a dummy (\textit{Down $>$ Up Bureaucrat}) equal to one in months and grids from which most of the district area is exposed to smoke, i.e. the area of the district downwind is larger than the area upwind from the grid.
      The table reports the mean of the dependent variable in grids that are upwind with respect to the borders of both the constituency and the district.   All regressions control for average wind speed and direction. Standard errors are clustered at the grid and district $\times$ month-year level. }   % sin cambios
    \end{minipage}
\end{figure}

 

 \paragraph{Robustness:} We start by testing whether the findings are robust to adopting different definitions of treatment. In Table~\ref{tab:main_by_treat}, we define a grid-cell as treated when most of the area of the constituency (rather than most of its population) lies downwind from the grid, finding effects of similar magnitudes and significance (column 2); we standardize the independent variable, instead of dichotomizing it, so that the treatment coefficient can be interpreted as the effect of a one standard deviation increase in the size of the downwind area (column 3); we adopt two other continuous definitions of treatment, the share of assembly constituency area that is downwind (column 4) and the difference between downwind and upwind area (column 5). In Table~\ref{tab:main_bydv} we vary the definition of the dependent variable, replacing the count of fires with a dummy equal to one when any fire took place in the grid-cell (column 1); with the natural logarithm of the count of fires plus one (column 2); and with a continuous variable capturing the average brightness of fires in a grid-month, a measure of fire intensity (column 3). Results are robust in all cases. When we disaggregate the analysis by year, the effect remains negative in each of the 10 years in our sample, with more imprecisely estimated effects (Table \ref{tab:by_year}). When we subset by state, the magnitude of the effect varies in line with expectations: the drop in fires is largest in Haryana (14\%), Punjab and Uttar Pradesh  (5 and 4\% respectively), where agriculture is most prevalent and crop burning salient, while there is  a positive effect in Bihar (Table \ref{tab:by_province}).
 We test the importance of differential pre-trends in Figure \ref{fig:es_main_rotated_panel} following the approach by \citet{dustmann2022reallocation}. \textcolor{purple}{ Panel (a) presents the standard  event-study estimates, while Panel (b) presents the de-trended event-study estimates, as in Figure \ref{fig:main_panel}b, following \citet{dustmann2022reallocation}. } We follow \citet{dustmann2022reallocation} in also conducting sensitivity analyses for violations of the pre-trends assumption using the smoothness restriction \citep{rambachan2023more}, showing that findings are robust to violations of parallel trends up to 0.75 times the size of the highest observed pre-trend deviations (Figure \ref{fig:es_main_rotated_panel}c).

 \paragraph{Strategic action and inaction near constituency borders:}

Our next question is whether the main effect, a reduction in crop-related fires when constituents are exposed, is driven by politicians becoming more active where pollution harms their own constituency, by politicians becoming less active where pollution is displaced onto neighboring constituencies, or by both. To answer this question, we examine how fire activity varies with strategic considerations about pollution exposure.


When burning occurs close to an upwind border, a larger share of the smoke remains within the politician’s own constituency, strengthening incentives to intervene. When burning occurs close to a downwind border, the smoke is shifted into a neighboring constituency, weakening incentives to act.
To test if politicians respond to both these incentives, we construct a new dataset at the level of grid-cell and each assembly constituency-border pair. Because wind direction varies over time, the same border can be downwind in some periods and upwind in others. We calculate the distance from each grid-cell to each relevant border and use this time-varying proximity to examine whether politicians adjust their behavior strategically as opportunities to internalize or externalize pollution become stronger.


The results in Figure \ref{fig:neighbor_analysis} reveal two patterns. First, fires decline monotonically as one approaches the upwind border, consistent with politicians acting to reduce fires when a larger share of their own constituency would bear the pollution costs. Second, fires increase as one approaches the downwind border, consistent with reduced effort when pollution can be offloaded onto neighboring constituencies. 

Together, these results show that politicians behave strategically along both margins: they increase effort when their own constituents are more exposed, and they relax effort when the costs can be shifted elsewhere. Additionally, politicians are also acting strategically in increasing the intensity of their action (and inaction) depending on the intensity of pollution exposure ---distance to the border.

  %%%%%%%%%%%%%%% Strategic Enforcement %%%%%%%%%%%%%

  \begin{comment}
  
\begin{figure}[htbp!]
    \centering
    \caption{\textbf{Effect of grid distance to the upwind (in blue) and downwind (in red) border on number fires }}
    \label{fig:neighbor_analysis}
        \includegraphics[height = 6cm]{figures/neighbor_output.pdf}
        \begin{minipage}{\textwidth}
            {\small \textit{Notes: }(Blue) Effect of being close to an upwind border with another constituency, where smoke would affect the politicians' assembly constituency. (Red) Differential effect (from the upwind border) of being close to a downwind border, where smoke would travel to the neighboring jurisdiction.  All effects are plotted by quintiles of distance to the border, and the 5th quintile, the farthest, is used as omitted category. The regression includes  grid and border $\times$ month-year fixed effects and control for average wind speed and direction. Standard errors are clustered at the grid and assembly constituency $\times$ month-year level. }
        \end{minipage}
\end{figure}

  \end{comment}

\subsubsection{Alternative Interpretations}\label{sssec:prosociality}

\paragraph{Farmers' Pro-sociality:}

A first alternative explanation is that fires may be reduced by farmers themselves, rather than by politicians, because farmers behave prosocially toward nearby residents who would be exposed to the resulting pollution. Under this account, farmers would care about neighbors as such, regardless of whether they belong to the same electoral constituency, since the constituency is an administrative unit with no obvious social or cultural meaning. 

To assess this possibility, we construct a placebo measure of exposure based on farmers’ local social environment rather than political boundaries. Specifically, around each grid we draw a placebo circle with a radius of 13 km, corresponding to the median size of an assembly constituency, and use it to approximate the set of nearby residents potentially relevant to farmers’ prosocial concerns. Figure \ref{fig:prosociality_ex} illustrates this strategy. In this analysis, a grid is coded as treated when the population living downwind within the placebo circle is larger than the population living upwind.
%\footnote{We compute the median size by creating circles that are tangent to each assembly constituency. Then, we measure the size of radius and obtain the median value.} 

 \ref{tab:main_ac_distr} \todo{Results now changed} Table \ref{tab:prosociality} reports the estimated treatment effect for the full sample (column 1), for locations that are treated both in the placebo specification and from the politician’s perspective (column 2), and for locations that are treated in the placebo-farmer specification but are control areas from the politician’s perspective (column 3). We find no evidence consistent with farmer prosociality. The reduction in fires appears only in locations where politicians have incentives to act, while in places where politicians have weak or no incentives, the estimated effect is even positive. This pattern suggests that the reduction in fires is not driven by farmers concern for downwind residents, but by political incentives operating through constituency boundaries.

\paragraph{Bureaucratic action:}
A second alternative explanation is that the main results are driven by bureaucrats' incentives rather than by politicians'. Indeed, in previous work we show that bureaucrats respond to crop burning when their own districts receive the brunt of air pollution versus their neighbors \citep{dipoppa2024bureaucrat}. The objective here is not to make the case that bureaucrat incentives do not matter; rather, we want to establish that political incentives also matter and explore the interaction of the two. To do so, we examine whether the effect of pollution visibility depends on whether exposure is concentrated only within the politician’s constituency or also within the jurisdiction of the relevant bureaucratic authority. In other words, we test whether politicians still reduce fires when the same pollution would affect not only to their own constituents but also district residents.

The results are reported in Table. We continue to find that both politicians and bureaucrats reduce fires when their respective jurisdictions are exposed to pollution from potential crop burning in a location but not vice versa suggesting that political and bureaucratic incentives matter on their own. However, we also find that the total effect of a location being a source of pollution for both the politician's constituency and the bureaucratic district is smaller, suggesting that the politician’s incentive to reduce visible pollution is weakened once responsibility is shared with the bureaucracy. One interpretation is a coordination failure: when both political and administrative authorities are implicated, each may expect the other to bear the costs of enforcement, resulting in less action overall. The fact that both actors matter reflects the layered governance structure of Indian anti-burning policy, in which politicians and bureaucrats share overlapping responsibility.


\section{Mechanism}

\subsection{Political Incentives Mobilize Government Machinery}

The evidence so far establishes that political incentives reduce fires, and that neither farmer prosociality nor bureaucratic action fully account for this reduction. We now ask whether the reduction can be traced to concrete political action. In the context section, we discussed qualitatively how politicians can concretely affect crop burning by galvanizing the machinery of the state. This section considers this quantitatively. 

Outside of legislation, where political will directly lays the foundation for policy action, most policy change is not implemented directly by politicians and requires that they coordinate with the bureaucracy. Therefore, for outcomes like crime reduction or increases in public spending, direct evidence of politicians' role is difficult to isolate \citep{gulzar2017politicians,agnihotri2026bureaucrat}. 

The same is true of controlling crop burning. However, two sources of data allow us to gain some insights into how political incentives may shape political action via the bureaucracy. 

First, since 2016, politicians across India convene formal meetings with district administrators through the District Development Coordination and Monitoring Committee (DISHA), which we described in the context section. The content of these proceedings offers a direct window into what issues politicians choose to raise and pursue. Second, India's eCourts records contain judicial cases filed under environmental laws, including crop burning offenses, allowing us to study whether political incentives activate enforcement.

\paragraph{DISHA Data}
As the DISHA institution gives politicians formal monitoring authority over bureaucratic action, the documentation it generates, together with the district material posted alongside it, offers a direct window into which issues are raised and pursued through the committee and the district administration that services it. We focus on whether the issues of agriculture and air pollution are tabled, a manifestation of the political will to act on crop burning through institutional channels.

DISHA proceedings and their associated documents are posted on the individual portal of each district, with no common catalog and no standard address. We therefore use a web-crawl strategy. For each district in our sample we generate candidate addresses from the recurring conventions of the \texttt{.nic.in} and \texttt{.gov.in} district domains, probe them in parallel, and retain those that resolve, adding the Haryana state rural development portal, the one state-level site in our sample that hosts DISHA material of its own. The `crawl' part of our strategy is that we do not stop at these entry pages. We follow links whose visible text or address carries a DISHA cue, including the Hindi rendering of DISHA, and  terms like ``Proceeding'', ``Minutes'' and ``Agenda'', one step deeper within the same website, and we retrieve every document posted on the pages we reach, including the further pages of paginated archives and the material that district administrations post alongside the minutes. This is what allows for a deeper understanding of the substance of the discussions rather than only the fact that a meeting took place. Each document is assigned to the district whose portal hosts it and to the date recorded inside it. Finally, we code the topic of each document, for example as relevant to agriculture, welfare, or environment.

\textcolor{red}{For example, from the DISHA page of Hisar district, Haryana (\url{https://hisar.gov.in/disha/}), we retrieve the agenda notes of the DISHA meeting of 20 June 2017, chaired by Shri Dushyant Chautala, Member of Parliament for Hisar, which we assign to Hisar district in June 2017. Its second agenda item is the Pradhan Mantri Fasal Bima Yojana, presented by the Deputy Director of Agriculture, reporting for Kharif 2016 and Rabi 2017 the number of insured loanee and non-loanee farmers, the insured crop area, the premiums collected and the claims paid, which we code as agriculture. The same page yields the meeting record posted for that date, the proceedings of a meeting of 23 October 2017 and of three later meetings, the district's MGNREGA target and achievement statement, the MPLADS progress report to 15 June 2017 and the membership list of the district DISHA committee.}

\textcolor{red}{Crop residue enters the collection through a second route, which illustrates why the crawl follows links rather than stopping at the pages it can address directly. Khagaria district, Bihar, posts its meeting records in a paginated proceedings archive. Following the pagination links beyond the first page, we retrieve from the second page of results a compendium of district task force proceedings of November 2022, chaired by the District Magistrate, which bundles the supply, agriculture, mining and health task forces into a single posted file. Within the proceedings of the District Agriculture Task Force, the item ``Crop Residue Management'' records that the District Agriculture Officer reported the stopping of farmer registration in the state direct benefit transfer system for 47 farmers found burning paddy residue, of whom five in the Alauli block and seven in the Gogri block, that street plays were staged on the subject with geo-tagged photographs filed to the public relations department, and that the incidence of burning had reduced significantly. This is district enforcement of the burning ban being minuted, and it is content that a crawl stopping at the first page of an archive would not reach.} See Appendix \ref{ssec:disha_app} for details of document retrieval and processing.

This procedure yields an unbalanced district-by-month panel of DISHA documents. To ensure that a zero reflects the absence of content related to the topics we are interested in, such as agriculture or air pollution, rather than the absence of any DISHA documentation, we include each district only from the first month-year in which a DISHA document on any topic is available.\footnote{This is because prior to a district's first documented meeting, a zero would be ambiguous: it could reflect that the issue was not raised, or simply that DISHA meetings had not yet been initiated or documented in that area.}


\paragraph{Courts Data} 
We also study enforcement action against crop burning via India's eCourts records. Our outcome is case filings, which measures the moment when the police transfer a case to the judicial system. This is the margin politicians plausibly influence by directing district officials to step up patrols, activating remote-sensing dashboards to identify hotspot villages, or signaling zero-tolerance at the start of the burning season. Politicians can thus prompt the police activity and effort that generates filings. Where DISHA captures what politicians choose to raise in coordination meetings, court filings capture whether the enforcement apparatus is actually engaged.

%Court cases are a measure of the state acting against crop burning, prompted by a combination of bureaucratic mobilization and the political pressure that activates it. Where DISHA captures what politicians choose to raise, court filings capture whether the enforcement apparatus is actually engaged. 

We access judicial case records from the SHRUG 
database \citep{aabbcdgns2025bias}, which digitizes the Indian eCourts platform. We retain cases filed under the Air (Prevention and Control of Pollution) Act, 1981, the statute under which crop burning offenses are prosecuted. \todo{Anzony: please fill in exactly how we filter these cases (act code field, keyword search, or both), how many cases the filter returns, and whether we can separate agricultural residue burning from industrial, vehicular, and construction violations under the same Act. If we cannot separate them, we need a sentence on why the resulting measurement error attenuates rather than confounds.} See Appendix \ref{ssec:courts} for further details. The sample runs from 2012 to 2018, the period covered by the SHRUG judicial extract.


\paragraph{Empirical Strategy}
We test whether DISHA meetings and court cases respond to the same political incentives that reduce fires. Both outcomes are recorded at the district level: DISHA meetings are convened by district, and eCourts records the location of the court rather than the location of the offense. We construct both as district-month panels and link them to our assembly constituency-month database by mapping each constituency to the district that contains it.\footnote{30 out of 853 assembly constituencies in our sample (3.5\%) span more than one district. Because documents and judicial case records are identified only at the district level, matching them to these constituencies would be ambiguous, and we therefore drop them from the analysis.}

To capture politician-level incentives, we aggregate our grid-level treatment into a continuous constituency-level measure which we call \emph{Potential Exposure}. As above, let $i$ index 5km$^2$ grid-cells and $j$ assembly constituencies, and let $\text{Down}_{it}$ denote the constituency population lying downwind of each grid $i$ in month-year $t$, the same quantity that underlies $\text{Down} > \text{Up}_{it}$. Let $\text{Pop}_j$ the constituency's total population. We define:
\begin{equation}
\text{Potential Exposure}_{jt} \;=\; \frac{1}{\text{Pop}_j} \sum_{i \in j} \text{Down}_{it},
\label{eq:potexposure}
\end{equation}

$\text{Potential Exposure}_{jt}$ measures the average share of the constituency's population that lies downwind of a randomly drawn grid-cell, given each grid's prevailing wind direction. Unlike a simple share, it can exceed one because the same constituent may lie downwind of multiple grids simultaneously, appearing in $\textit{Down}_{it}$ for each of them. It is high when, across most locations in the constituency, wind would carry smoke toward the politician's own constituents, making them potentially exposed to fires. %the total population inside the politician's own constituency that \emph{potentially} receives pollution from an upwind fire.
We standardize $\text{Potential Exposure}_{jt}$ over the estimation sample (mean = 60.12, SD = 37.71), so a one-unit increase corresponds to a one-standard-deviation increase in the number of own-constituency potential pollution for an average constituent. 

For outcome $Y$ in constituency $j$ and month-year $t$, we estimate the impact of potential exposure on the likelihood of observing political action in a given constituency-month $Y_jt$:
\begin{equation}
Y_{jt} \;=\; \lambda \, \text{Potential Exposure}_{jt} + \mu_j + \nu_t + \xi X_{jt} + \omicron_{jt},
\label{eq:mech}
\end{equation}
where $\mu_j$ are constituency fixed effects, $\nu_t$ are month-year fixed effects, and $X_{jt}$ is a vector of controls containing the constituency's average wind speed and direction. Standard errors are clustered at the constituency level.



\paragraph{Results}

Table~\ref{tab:mechanism} presents the results, starting from DISHA outcomes (columns 1 and 2). A one-standard-deviation increase in potential exposure reduces the probability that agriculture is mentioned in Disha documents by 62\% relative to a baseline rate of 0.5 mentions per constituency-months, and increases the probability of that an environment-related document is produced by 81\% relative to a baseline of 0.3 mentions. These are large effects given how rarely these topics appear in DISHA documents absent accountability pressure.\footnote{Appendix Table~\ref{tab:court_disha_app} shows that this response is specific to agriculture and the environment because higher potential exposure does not increase discussion of other topics; if anything, documents related to welfare and development declines, suggesting that politicians reallocate their agenda towardenvironmental concerns and actively away from agricultural concerns when constituent exposure to pollution increases.}

Turning to court cases in Columns 3 and 4, we find that a one-standard-deviation increase in politicians' incentives yields an increase of 0.21 more cases within the same month. In the following month,  exposure results in 0.14 additional court cases against crop burning per constituency-month, an 11\% increase relative to the mean ($p<0.01$). This delayed response is consistent with the time required for enforcement actions to materialize as cases are passed to the judicial system only after necessary evidentiary work by the police. Given that severe penalties are rarely issued (on average, less than one per district every month), this represents a substantively meaningful enforcement response.


Taken together, these results show evidence for the mobilization of government machinery in response to changes in political incentives. However, these estimates should be interpreted with two caveats. Because treatment varies continuously and over time, with units switching in and out of treatment, the design is too complex to permit formal validation of difference-in-differences assumptions. The low geographic granularity of the outcome data also means that several coefficients are imprecisely estimated. Notwithstanding these limitations, the findings provide direct evidence that politicians translate accountability incentives into concrete institutional action, both by steering bureaucratic attention toward crop burning in coordination meetings, and by directing subsequent enforcement through the courts.


\begin{table}[htbp!] 
\caption{Pollution Exposure, Politicians' Actions, and Electoral Consequences Pop}\label{tab:mechanism}
        \resizebox{\linewidth}{!}{
        \input{tables/action_final4}
        }
    \begin{minipage}{\textwidth}
        {\footnotesize \textit{Notes:} The table reports politicians' actions (columns 1-4) and electoral outcomes (columns 5-7) by level of constituency exposure.
        \textit{Actual Exposure} is the sum across grids of whether any fire was recorded in that same calendar month of the previous year weighted by the  constituency population lying downwind of that grid, normalized by the  constituency's total population, and standardized (z-scored) within the  estimation sample.
        %
       \textit{Actual Exposure$_{t-1}$} is defined analogously where variables are averaged at  the constituency-by-electoral-term level before standardization.
        Columns (1) and (2) use DISHA documents to report the probability that a document is recorded per 100 in a given month-year, for the following categories: agriculture (column 1) and environmental pollution (column 2). Columns (3) and (4) report the change in probability of filing court cases about pollution in the same month and in the following month, respectively. Column (5) reports the probability that the incumbent politician runs for re-election. Column (6)-(7) reports the incumbent's vote share in the following election. Column (7) subset for districts with at least one court case in the previous term.
        The table reports the mean of the dependent variable in the sample. All regressions in columns (1)-(4) control for wind direction and average wind speed. Standard errors are clustered at the constituency level.
        }
    \end{minipage}
\end{table}



\subsection{Exposure to Fires Drives Electoral Outcomes}

The previous subsection showed that politicians respond to accountability incentives by mobilizing state action against crop burning. To complete evidence on the way this accountability mechanism operates, we now show evidence on voters: constituents should reward politicians who reduce fires and punish those who do not. We examine this directly by asking whether constituents' \textit{actual} exposure to pollution over an electoral term affects incumbent electoral performance.

\paragraph{Empirical Strategy}
To answer this question, we create a modified measure of exposure, which we call $\text{Actual Exposure}_{j}$. While for the first mechanism subsection we tracked month-by-month intensity in politicians' incentives to act using potential exposure, this new variable is constituency-by-electoral-term measure of the cumulative exposure of voters to actual fires. % how much fire-related pollution did constituents actually endure, and was it attributable to their politician? A politician who failed to suppress fires when accountability incentives were present should be electorally punished for it.

For each constituency-term in office $o$, let $M_o$ denote the number of months in the term. We define:
\begin{equation}
\text{Actual Exposure}_{jo} \;=\; \frac{1}{M_o} \sum_{t \in o} 
\left( \frac{1}{\text{Pop}_j} \sum_{i \in j} \text{Down}_{it} \times \text{\#Fires}_{it} \right),
\label{eq:actexposure}
\end{equation}
where $M_o$ is the number of months in term $o$, $\text{Down}_{it}$ equals the number of people living downwind of grid $i$ in month-year $t$, and $\text{\#Fires}_{it}$ is the number of fires in grid $i$ in that month, such that when there are no fires, the quantity inside the parenthesis equals zero. 

$\text{Actual Exposure}_{j}$ therefore represents the total population in the constituency exposed to any fire-month during an electoral term. As before, we regress incumbent electoral outcomes on the standarized version of this actual exposure (mean = 1040.09, sd = 2973.72), using constituency and state times election year fixed effects, and clustering standard errors at the constituency level.

\paragraph{Results}
We find that constituents exposed to more fires during an electoral term punish their incumbent in Table~\ref{tab:mechanism}, Columns 6-8. A one-standard-deviation increase in accountability-relevant fire exposure is positively the probability that the incumbent chooses to run again by 0.054 percentage points, a 57\% decrease relative to the baseline rerunning rate of 0.094 ($p<0.1$). The same increase in exposure reduces incumbents' unconditional vote share by 2.1 percentage points, a 61\% decrease relative to the mean at baseline ($p<0.1$), though this decrease is mostly driven by the reduced re-contestation.

These effects indicate that when voters are exposed to more pollution from fires during when their politician had clear accountability incentives, they respond by withdrawing electoral support validating the mechanism underlying the main analysis: politicians suppress fires because failure to do so is electorally costly.

\section{Countervailing Incentives}

The first part of the paper provided evidence that there are prospects for environmental accountability. In this second part, we consider how accountability can weaken or break down with two fundamental levers that can act as countervailing incentives: organized lobbying, and adverse political selection.

 \subsection{Lobbying\label{ssec:lobbying}}
 
 From June 2020 to December 2021, India experienced the largest protest movement in the world history: 25 million farmers mobilized in large-scale protests against agricultural laws that aimed to deregulate the sale of farm produce, allow contract farming, and remove stockholding limits. One key reason for mobilization was the request to decriminalize crop burning, removing strict penalties such as high fines and jail terms (\citep{guardian2020farmers}). Farmers staged demonstrations, road blockades, and sit-ins, particularly in Punjab, Haryana, and at the borders of New Delhi. The demonstrations were successful: they led to the repeal of the laws and to the depenalization of stubble burning. This section analyzes the impact of all 390 farmers' protests events on crop burning, and asks whether those counteract accountability incentives. \textcolor{blue}{While those are large scale protests, we use the term `lobbying' to refer to the organized pursuit of policy concessions by groups with a concentrated economic interest in a policy, even when this interest conflicts with the diffuse welfare of broader populations, consistent with canonical examples from the literature \citep{olson1971logic}, from taxi driver strikes to energy sector mobilization, in which sector-based collective action secures concessions at public expense.}

\subsubsection{Data\label{sssec:data2}}
We gather information on the date and location of each protest event related to the Farmers' Movement from June 2020 to December 2021 from \citet{acled}. We identify farmers' protests by searching for any conjugation of the word  ``farm’’ in the event description, considering both type and participants to the event. We validate this process by reading the full description for a random sample of events.  Table \ref{tab:_desc_incentive_main} present the descriptive statistics for the sample. Figure \ref{fig:protest_by_time} shows the temporal distribution of protest occurrence, highlighting a peak in late 2020 after the passage of the farm laws (September 2020). Protests were highly prevalent in Punjab but widesperead across all ares in our sample (Figure \ref{fig:map_protest}).

\subsubsection{Empirical Strategy\label{sssec:empirical2}}
\paragraph{Treatment Definition:} We code proximity to a farmer protest by drawing a %Politicians are more likely to respond to farmer lobbying in areas near protest sites, where the political pressure they face is most salient. We construct our treatment variable to capture this geographic exposure to farmer mobilization. First, we draw 
5km-radius buffer areas around the location of the 390 protests occurring between June 2020 and December 2021. We code a grid-cell as treated if it falls within at least one buffer area, from the month-year when a protest first occurs until the end of our sample period.\footnote{We consider as treated grid-cells that fall within a buffer area regardless of whether they fall inside or outside the constituency boundaries. Grids located in overlapping buffer areas are assigned to the buffer associated with the earliest protest. } Figure \ref{fig:protest_construction} provides a visual illustration of the grid-cells coded as treated in a constituency. %why do we choose 5km? 

\paragraph{Estimation Strategy:}
We estimate the impact of proximity to a protest event using a stacked design following \citet{cengiz2019effect}. For each event-time cohort in our data, defined as the month-year in which a location first experiences a protest, we build one cohort-specific dataset. Each dataset includes only never-treated units --grids that never experience a protest-- and units first treated in that event-time --grids experiencing their first protest-- thereby eliminating problematic comparisons in the presence of heterogeneous treatment effects. Each cohort is further restricted to observations within a single electoral term, so that treated and control grids are compared only within the same term. We then stack these cohort-specific datasets into a pooled database. For grid-cell $i$ in constituency $j$ in province $p$, relative time from treatment $r$ and electoral cycle $s$ in cohort dataset $c$, we estimate: %note that time fixed effect is not the same as calendar year, but 12 months befor/after treatment 
\begin{equation}\label{eq_protest}
 { \text{Fires} }_{i c r}=\beta {\text{Post$\times$Protest}}_{i c r}+ \alpha_{ic}+ \gamma_{rc} + \theta_{p \times s}+ X_{ipr} + \epsilon_{icr}
\end{equation}
where the dependent variable is the number of fires per 1,000 units. We include a comprehensive set of fixed effects which account for cohort-specific and constituency specific factors. First, $\alpha_{ic}$ are grid $\times$ cohort fixed effects that absorb time-invariant factors for each grid in our sample within each sub-experiment. Second, $\gamma_{tc}$ are relative time from treatment defined for every cohort that account for time-varying shocks at the national level, including changes in the legislative environment following protest events. Third, $\theta_{ps}$ are province $\times$ month-year $\times$ cohort -specific electoral cycles fixed effects that account for differences across cycles within a province. $X$ is a vector of controls that includes province $\times$ cohort  linear time trends, average wind speed and direction. \textcolor{purple}{While the stacked dataset is constructed at the month-year level, we aggregate relative-time into years (12 months from treatment) when plotting the event-study estimates to reduce noise.} Standard errors are clustered at the buffer area level.

\subsubsection{Results\label{sssec:results2}}

%\textcolor{red}{Take-away: Accountability incentives after protest return more fires -- accountability collapses to organized farmer protest. Accountability incentives after a protest in places without protest also have an effect: there are spillovers of protests in places that didn't protest, as we observe more fires there too in the post period. Overall, there is a reduction of the impact of accountability.}

After the first protest occurs in a given area, we observe 0.105 more fires per grid-month, a 75.1\% increase with respect to an average of 0.139 fires in treated areas before any protest occurs (Table \ref{tab:protest_downup}, columns 1-3). This result is in line with politicians relaxing environmental control in response to concentrated farmer pressure.
Figure \ref{fig:es_protest} plots the event-study results showing an increase in the number of fires after the first protest. This effect corresponds to 92.43\% more fires with respect to baseline. This indicates that politicians relax environmental enforcement in response to farmers organized protests. %This effect is driven by areas with above-median rice production, where we observe 25\% more fires after a protest with respect to places with below median rice production (Figure \ref{fig:es_protest_rice}). 

%Results are robust to using a standard two-way fixed effects specification,  with smaller magnitudes potentially driven by negative weighting from treatment effect heterogeneity (Table \ref{tab:table_protest}). \red{Maybe drop/comment out this as we do not have the same robustness for politicians}

\textcolor{purple}{Figure \ref{fig:protest_main_rotated_panel} presents the test of the importance of differential pre-trends, following \citet{dustmann2022reallocation}. Panel (a) presents the standard event-study estimates, while Panel (b) presents the de-trended estimates. Panel (c) presents the sensitivity analysis for violations of the parallel-trends assumption using the smoothness restriction of \citet{rambachan2023more}, showing that our findings are robust to violations of parallel trends of more than 1 times the size of the largest observed pre-trend deviation.}\todo{Do we have what's the max of M? Do we need it if larger than 1?}

% c a b
\begin{figure}[hbtp!]
    \centering
    \caption{\textbf{Effect of farmers' protests occurrence on number of fires}}
    \label{fig:main_protest}

    % ---------- Row 1: subfigures (a) and (b) side by side ----------
    \subcaptionbox{Protest occurrence by location\label{fig:map_protest}}[0.48\textwidth]{%
        \includegraphics[width=.8\linewidth]{figures/5km_plot.png}%
    }%
    \hfill
    \subcaptionbox{Farmers' Protest\label{fig:protest_photo}}[0.48\textwidth]{%
        \includegraphics[width=.8\linewidth]{figures/2020_Indian_farmers_protest.jpg}%
    }

    \vspace{1em}

    % ---------- Row 2: subfigures (c) and (d) ----------
    \subcaptionbox{By relative time period\label{fig:es_protest}}[0.48\textwidth]{%
        \includegraphics[width=\linewidth]{figures/_app_17_5km_fe12_evst_all_rural_fe03_1_ori.png}%
    }%
    \hfill
    \subcaptionbox{By constituency exposure to smoke\label{fig:protest_downup_coefplot}}[0.48\textwidth]{%
        \includegraphics[width=\linewidth]{figures/_app_18_protest_5km_fe12_did_downup_plot_rural_acpop.png}%
    }
    
    \vspace{1em}
    
    % ---------- Row 2: subfigure (c) centered alone ----------
    
    
    \begin{minipage}{\textwidth}
        {\small \textit{Notes: (a)} The map shows treated (blue) and control (white) grid-cells. A grid-cell is coded as treated if it falls within 5km-radius buffer areas around protest locations. Treatment begins in the month-year of the first protest within each buffer area. Control grid-cells are located outside all buffer areas.
         \textit{(b)} This photograph is called ``Art, pen and people''. It was taken in the 2020 Indian farmers' protest by opposing Modi's government farm laws \citep{maddoke2020art}.
         \textit{(c)} Event-Study results from Equation \ref{eq_protest} %after de-trending for the pre-treatment period linear trend following \citet{dustmann2022reallocation}
         . The reference category is t=-1, a year before the protest occurrence. 
        The regression includes grid-cohort, relative time-cohort, and province linear time trends fixed effects, and controls for  average wind speed and direction. Standard errors are clustered at the buffer area $\times$ assembly constituency level. The control group includes never treated grids only. 
        \textit{(d)} 
        Results from Equation \ref{eq_protest} with the additional interaction of a dummy equal to one in months and grids from which most of the assembly constituency area is exposed to smoke ($Down>Up$, see Section \ref{ssec:empirical1} for details).
        Plotted coefficients report the standarized to the pre-period effect of being in a constituency exposed to smoke in the post-protest period, separately for constituencies that did and did not experience protests. The reference category is non-protest areas, before protests when they are downwind ($Down>Up$=1). 
        Table \ref{tab:protest_downup} (columns 4-6) reports the corresponding regression coefficients.
        The regression includes grid-cohort, relative time-cohort, year-cohort and province linear time trends fixed effects, and controls for average wind speed and direction. Standard errors are clustered at the buffer area $\times$ assembly constituency level.}
    \end{minipage}
\end{figure}






\paragraph{Erosion of Accountability Incentives:}
Organized pressure from farmer lobbies might contribute to influence the accountability incentives we document in the previous section: do politicians continue reducing fires when smoke would impact their constituents, even after farmers' protests pressure them to stop sanctioning them for crop burning? Figure \ref{fig:protest_downup_coefplot} explores this question by interacting the two treatments, the effect of a grid becoming a source of exposure for its constituency, and protests. Before protests, areas whose pollution burden falls on their own constituents experience fewer fires relative to areas that would pollute neighboring jurisdictions, in line with the accountability mechanism. 

In protest-affected areas,  the accountability effect disappears entirely after protests begin: fires increase significantly, offsetting the reductions in crop burning otherwise associated with political accountability. In places without protests, the accountability effect weakens, but this does not translate into higher crop-related fires relative to baseline. This pattern is consistent with organized lobbying increasing fires directly where protests occur, while also weakening politicians’ incentives for environmental enforcement beyond the protest locations.


\FloatBarrier
 \subsection{Adverse Selection}

 
In India, candidates with direct ties to agriculture are often elected to state legislatures, especially in rural constituencies where farming remains economically and politically central. Once in office, these politicians may use their authority to pursue policies that are more aligned with the interests of their sector than with those of other constituents. We therefore examine whether the selection of politicians with agricultural ties shapes the incidence of fires after elections, and whether it alters the accountability incentives documented in the first part of the paper. 

 \subsubsection{Data}

 We build new data on the sectoral identity of elected politicians in each Assembly Constituency in our sample by scraping information about politicians' profession and properties from Myneta.info \citep{myneta}. The website provides background information on candidates contesting elections from 2008, including profession, and details on movable and immovable assets. To capture politicians' self-interest in agriculture, we classify whether each politician's self-reported occupation is agriculture-related, using a large language model (GPT)\footnote{Most politicians in our sample report holding agricultural assets, leaving little variation to exploit. We therefore rely on self-reported profession as our primary measure of self-interest in crop burning.} (See Section \ref{sssec:myneta_app} for additional details.).
Our sample comprises 2,565 politicians, of whom 1,102 (43\%) have agricultural professions.  %To minimize prompt contamination and variance, each attribute is queried independently. Outputs are encoded as binary indicators (1 if agriculture-related profession; 0 otherwise). 
% Inference is run with temperature = 0 to ensure deterministic outputs and reduce hallucinations. 
Table \ref{tab:_desc_adverse_main} presents descriptive statistics for the resulting sample. 

\subsubsection{Empirical Strategy \label{sssec:empirics3}}
\paragraph{Treatment Definition:} 
A comparison of constituencies that ever elect agriculturalist politicians versus not would conflate politician characteristics with place characteristics. Rural localities with strong agricultural ties may be more likely both to elect agriculturalists and to experience more agricultural fires. 
% Similarly, a close election regression discontinuity may not be able to to isolate the effect of agricultural ties: constituencies where an agriculturalist barely wins may differ systematically from those where one barely loses, reflecting underlying local characteristics that determine both electoral outcomes and fire activity.
To isolate the effect of electing an agriculturalist from underlying place characteristics, we exploit within-location variation of winning politicians.  We compare crop related fires before and after the same constituency switches from electing a non-agriculturalist to electing an agriculturalist. %Our control group consists of constituencies that elect non-agriculturalists throughout the sample period.

% Anzony -- I am not sure if this is perfectly stated.

\paragraph{Estimation Strategy:}
We estimate the impact of a constituency switching to electing an agriculturalist politician using a stacked design \citep{cengiz2019effect}. For each switch --- defined as a constituency electing an agriculturalist politician after not having one in the prior election--- we build a cohort-specific dataset spanning 10 years (2 elections) around the event. Never-treated units, i.e. constituencies that elect a non-agriculturalist politician throughout the 10 year window, are the control group.  To avoid compositional changes that might bias our estimates, each dataset is a balanced panel, including only units that we observe for the entire 10 year period. We then stack these cohort-specific datasets into a pooled database. For a 5km$^2$ grid-cell $i$ in constituency $j$ in province $p$, relative time from treatment $r$, election $s$ in cohort-dataset $c$, we estimate:
\begin{equation}\label{eq_agriculture}
 { \text{Number of Fires } }_{i c r}=\beta {\text{ Post$\times$Agriculturalist}}_{i c r}+ \alpha_{ic}+ \gamma_{tc} + \theta_{ps}+ X_{ipr} + \epsilon_{i c r}
\end{equation}

where the dependent variable is the number of fires per 1,000 units; $Post$ is a dummy that takes value one from the cohort-specific year onward; $Agriculturalist$ is a dummy equal to one after the election of a politician with agricultural ties. We include the same set of fixed effects as in the protest specification, which account for cohort-specific and electoral factors. \textcolor{blue}{In particular, the grid x cohort fixed effects guarantees that the estimate is identified from within-location changes over time, rather than from comparing potentially different constituencies that never elect agriculturalist politicians. The province × cohort × electoral cycle FEs ($\theta_{ps}$) absorb the state-level political environment in each legislature term. For example, if BJP swept Uttar Pradesh in 2017, nominated more agricultural candidates across the state, and also changed state-level burning policy, that is absorbed. 
}
 \textcolor{purple}{While the stacked dataset is constructed at the month-year level, we aggregate relative-time into years (12 months from treatment) when plotting the event-study estimates to reduce noise.
Standard errors are clustered at the assembly constituency $\times$ election year level. 
}

\subsubsection{Results\label{sssec:results3}}

%\textcolor{purple}{\ref{tab:agri_profession_rice_table}}

When a constituency switches to electing an agriculturalist politician, fires increase by 6\% with respect to baseline, corresponding to about 7 additional fires per grid-year as shown in Table \ref{tab:agri_profession_downwind}, Column 3. The effect is driven by rice-producing areas, where we observe 50\% increase following the election of a politician with agricultural ties  (Figure \ref{fig:es_agri_rice}). The de-trended event study in Figure \ref{fig:es_agri_profession} confirm this pattern:  fires are 23\% higher after a constituency switches into the electoral term of a agriculturalist politician, relative to the control group. Politicians with agricultural ties behave in line with the interest of their professional group, avoiding enforcement pressure on environmental damage.
\textcolor{purple}{Figure \ref{fig:protest_main_rotated_panel}a the standard event-study to test for the importance of differential pre-trends, following \citet{dustmann2022reallocation}. Figure \ref{fig:protest_main_rotated_panel}c shows the sensitivity analysis for violations of the parallel-trends assumption using the smoothness restriction of \citet{rambachan2023more}, showing that our findings are robust to violations of parallel trends of 0.75 times the size of the largest observed pre-trend deviation.}



%%%%%%%%%% Agriculturalist Politicians %%%%%%%%%

% \begin{figure}[hbtp!]
%     \centering
%     \caption{\textbf{Effect of switching to an agriculturalist politician on number of fires}}
%     \label{fig:main_agri}
%      \subcaptionbox{By relative time period\label{fig:es_agri_profession}}[\textwidth]{%
%      \centering
%     \includegraphics[width = 0.65\textwidth]{figures/_app_16_polischar_fe12_evst_all_rural_1_rotated.png}
%   }
%    \subcaptionbox{By constituency exposure to smoke\label{fig:agr_downup}}[\textwidth]{%
%     \centering
%     \includegraphics[width = 0.65\textwidth]{figures/Interaction_downwind/Agri_interaction3.png}
%   }

%   \includegraphics[width = 0.65\textwidth]{figures/panelview_self_profession.png}

  
%     \begin{minipage}{\textwidth}
%         {\small \textit{Notes: }\textit{(a)} Event-Study Results from Equation \ref{eq_agriculture} after de-trending for the pre-treatment period linear trend following \citet{dustmann2022reallocation}. The reference category is t=-1, one year before the province election takes place. 
%         The regression includes grid-cohort, and legislature fixed effects, and controls for province linear time trends, average wind speed and direction. Standard errors are clustered at the buffer area $\times$ assembly constituency level.  The control group consists of never-treated observations over the period of analysis—that is, assembly constituencies, whether inside or outside the province holding the election, whose politicians' self-reported occupation is not agriculture-related in the electoral terms just before and after the election. \textit{(b)} Results from Equation \ref{eq_agriculture} with the additional interaction for constituency exposure to smoke. $Down>Up$ is a dummy equal to one in months and grids from which most of the assembly constituency area is exposed to smoke (see Section \ref{ssec:empirical1} for details).
%         %
%         %
%         Plotted coefficients report the effect of being in a constituency exposed to smoke in the pre- vs.\ post-protest period, separately for constituencies that did and did not experience protests. The reference category is non-protest areas, before protests when they are upwind ($Down>Up$=0). 
%         % In the plot, Pre-Non Agricultural Politician corresponds to be in a $Down>Up$ zone in the Pre-Switching period with no Agriculturalist Politician ($Down >Up$), Post-Non Agriculturalist Politician corresponds to be in a $Down>Up$ zone with no agricultural politician after an election period ($Down>Up + Down >Up\ \times Post$), and post-Protest corresponds to switching to an agricultural politician in a $Down>Up$ zone ($Down>Up + Post \times Protest + Post \times Down>Up + Post \times Protest \times Down>Up$).
%         %
%         Table \ref{tab:agri_profession_downwind} reports the corresponding regression coefficients.
%         %
%         The regression includes grid-cohort, relative time-cohort and electoral cycle fixed effects, and controls for province linear time trends, average wind speed and direction. Standard errors are clustered at the buffer area $\times$ assembly constituency level.  
%          } 
%     \end{minipage}
% \end{figure}
% c a b for 
\begin{figure}[hbtp!]
    \centering
    \caption{\textbf{Effect of switching to an agriculturalist politician on number of fires}}
    \label{fig:main_agri}
    
    % ---------- Row 1: (a) and (b) side-by-side ----------
    \subcaptionbox{Treatment status by grid over time\label{fig:panelview_agri}}[0.48\textwidth]{%
        \centering
        \includegraphics[width=\linewidth]{figures/panelview_self_profession.png}
    }
    \subcaptionbox{My Neta: Profession\label{fig:self_profession_report}}[0.48\textwidth]{%
        \centering
        \includegraphics[width=\linewidth]{figures/myneta_example2.png}
    }
    \subcaptionbox{By relative time period\label{fig:es_agri_profession}}[0.48\textwidth]{%
        \centering
        \includegraphics[width=\linewidth]{figures/_app_16_polischar_fe03_evst_main_rural_acpop_1_ori.png}
    }%
    \hfill
    \subcaptionbox{By constituency exposure to smoke\label{fig:agr_downup}}[0.48\textwidth]{%
        \centering
        \includegraphics[width=\linewidth]{figures/_app_19_polischar_fe12_did_downup_inter_plot_rural_acpop.png}
    }
    
    \vspace{1em}
    
    % ---------- Row 2: (c) panelview, larger ----------
    
    
    \begin{minipage}{\textwidth}
        {\small \textit{Notes: }\textit{(a)} 
        The plot reports treatment status over time at the grid level. Each row represents a grid, sorted by first treatment time, while each column represents a month-year. Dark blue cells highlight the periods in which agrid had a self-reported agriculturalist politician in power.
        \textit{(B)} Figure B shows an example of a Self-Declared Agriculturalist Politician report in My Neta Website.  
        \textit{(c)} Event-Study Results from Equation \ref{eq_agriculture}. %after de-trending for the pre-treatment period linear trend following \citet{dustmann2022reallocation}. 
        The reference category is $t=-1$, one year before the election of an agriculturalist politician. 
        The regression includes grid-cohort, and legislature fixed effects, and controls for province linear time trends, average wind speed and direction. Standard errors are clustered at the buffer area $\times$ assembly constituency level. The control group consists of never-treated grids over the period of analysis ---that is, assembly constituencies, whether inside or outside the province holding the election, whose politicians' self-reported occupation is not agriculture-related in the electoral terms just before and after the election. 
        \textit{(d)} Results from Equation \ref{eq_agriculture} with the additional interaction for constituency exposure to smoke. $Down>Up$ is a dummy equal to one in months and grids from which most of the assembly constituency area is exposed to smoke (see Section \ref{ssec:empirical1} for details).
        Plotted coefficients report the effect of being in a constituency exposed to smoke in the pre- vs.\ post-protest period, separately for constituencies that did and did not elect an agriculturalist. The reference category is constituencies with a non-agriculturalist in power, before elections when they are upwind ($Down>Up$=0).
        Table \ref{tab:agri_profession_downwind} reports the corresponding regression coefficients.
        The regression includes grid-cohort, relative time-cohort and electoral cycle fixed effects, and controls for province linear time trends, average wind speed and direction. Standard errors are clustered at the buffer area $\times$ assembly constituency level.
        }
    \end{minipage}
\end{figure}

\paragraph{Erosion of accountability:} Figure \ref{fig:agr_downup} examines whether the election of an agriculturalist politician erodes the accountability incentives documented above by estimating the joint interaction of the two treatments. Before elections, fires fall in areas where politicians’ own constituents bear the pollution burden, consistent with the baseline accountability pattern. After elections, however, the effect depends on who is elected. In constituencies that continue electing non-agriculturalists, fires remain lower when pollution affects the politician’s own constituents. In constituencies that switch to an agriculturalist, instead, fires increase relative to the control group. These results suggest that politicians with personal ties to agriculture undermine the responsiveness to constituent exposure documented in the first part of the paper, consistent with the selection of politicians whose private incentives are misaligned with those of the majority of their constituents.

\section{Conclusions}

Environmental governance is a hard case for democratic accountability: pollution generates concentrated benefits and diffuse harms, and voters rarely prioritize the environment. Against this backdrop, we find that politicians do curb fires, when wind patterns shift smoke toward their own constituents. This is, to our knowledge, among the first large-scale evidence that accountability incentives can generate environmental protection even in a setting with millions of dispersed polluters and limited monitoring.

Accountability, however, does not hold to the competing pressure of concentrated interests and adverse selection. We show that organized lobbying by the farmers' movement increased burning, overriding politicians' responsiveness to their own constituents' welfare. In addition, we document that the election of politicians with private agricultural interests increases burning, even where accountability incentives would otherwise reduce it. Taken together, these results show that democratic accountability can produce environmental regulation, but concentrated interests that benefit from pollution can undermine it.

These findings speak to institutional design. Whether anti-burning policy works in India depends not only on national mandates and enforcement capacity but also on political incentives. 
Interventions that sharpen the link between pollution exposure and political responsibility, such as real-time air quality monitoring and transparent attribution of pollution to specific jurisdictions, could strengthen the accountability channel we document \citep{singh_thachil_2023_airpollution}. But the erosion we find from organized interests and adverse selection suggests that electoral mechanisms alone might not be enough to confront the large challenge of air pollution.

Finally, our results inform a broader question about whether democracies can manage environmental externalities. The evidence from South Asia's air pollution crisis is that the capacity for accountability exists: politicians respond when their constituents pay the costs of pollution. The challenge is that this response is easily overwhelmed by the political economy forces arrayed against it. Sustaining environmental protection in democracies requires not just accountability but institutions that can withstand the pressures that erode it.


\clearpage
\FloatBarrier
\printbibliography



\clearpage
\appendix

\setcounter{page}{1}
\setcounter{figure}{0}
\setcounter{table}{0}
\counterwithin{figure}{section}
\counterwithin{table}{section}
\onehalfspacing



\section*{Appendix}
\section{Data Construction}
\subsection{Combining all the data at a 5 square km grid-cell level}\label{ssec:data_construction}
We combine all the data into a single dataset at the $5x5km$ grid-cell level. Wind and fire grids do not overlap precisely: to reduce errors in the assignment of fires to wind grids, we split the $25km$x$25km$ wind grid-cells into squares of $5x5k$ each. This guarantees that fires are assigned to the wind grid-cell which contains most of the fire grid-cell. Splitting wind grids in smaller parts also allows us to address another concern: some grids might overlap across jurisdiction boundaries. By reducing their size, we minimize the amount of misattribution and we are able to assign each grid to the appropriate jurisdiction. Any remaining error should bias our findings toward zero, as errors are uncorrelated with the dependent and independent variables. In case of overlap, we assign a grid to the jurisdiction that contains a larger area of the grid. For crop data, we assign the pre-period crop production to each $5km^2$ grid-cell to classify what crops are produced in each location. 

\subsection{Coding Farmer Politicians}\label{sssec:myneta_app}
We scraped Myneta.info, an online data repository of election-related information, to collect data on politicians. The website contains information on politicians' profession and their asset declarations.

To identify agricultural ties, we employed a large language model (\emph{GPT}) and fix the temperature to 0 for inference, as this reduces variance in the generation and helps reduce hallucinations. We supplied each profession description to the model and recorded whether it described any kind of agriculture-related job; if so, the politician is classified as having an agricultural profession. Formally, a politician is coded as being an agriculturalist if the profession description refers to a profession related to agricultural.  We are able to construct the indicator for every winning candidate in our sample. Only for 50 out of 22,652\todo{check} politicians the website does not provide any information. 

For validation we carried out two exercises: First, we ran the classification procedure ten times for each observation and assigned the final category using the modal classification across the ten runs. Second, we manually reviewed a stratified random sample of 10 percent of observations (stratifying on LLM positives and negatives of farmer politicians), and found zero errors. 

The following fixed prompt was supplied to GPT for the coding of politicians' agricultural assets and professions: 
\begin{Verbatim}[fontsize=\small]
You are a research assistant interested in political capture in India. You are researching whether politicians with agricultural assets are more prone to taking actions that favor agricultural businesses. Your current task is to assess whether an individual owns agricultural assets.\ You will be provided a table with two columns. The table summarizes the immovable assets owned by the individual. The first column describes the type of immovable assets listed in the corresponding row. The Second column lists the immovable assets of the corresponding type. The individual is considered to own agricultural assets ONLY IF any are listed in the row corresponding to agricultural assets. If the value says \"Nil\", we can infer that the individual does not own any such assets.\ Provide a JSON blob with your assessment, structured as follows. First, the \"explanation\" field should include your reasoning for your assessment on whether the individual owns agricultural assets. Second, the "owns_agricultural_assets" must include a boolean, true if the individual is assessed to own any agricultural assets, and false otherwise.\ You must only use the information included in the row of agricultural assets; any other kind of asset should be ignored.
\end{Verbatim}

We provide additional information on the use of AI in section \ref{ssec:ai}.

\begin{comment}

\begin{Verbatim}[fontsize=\small]
You are a research assistant interested in political capture in India. You are researching whether politicians with agricultural assets are more prone to taking actions that favor agricultural businesses. Your current task is to assess whether an individual owns agricultural assets.\ You will be provided a table with two columns. The table summarizes the immovable assets owned by the individual. The first column describes the type of immovable assets listed in the corresponding row. The Second column lists the immovable assets of the corresponding type. The individual is considered to own agricultural assets ONLY IF any are listed in the row corresponding to agricultural assets. If the value says \"Nil\", we can infer that the individual does not own any such assets.\ Provide a JSON blob with your assessment, structured as follows. First, the \"explanation\" field should include your reasoning for your assessment on whether the individual owns agricultural assets. Second, the "owns_agricultural_assets" must include a boolean, true if the individual is assessed to own any agricultural assets, and false otherwise.\ You must only use the information included in the row of agricultural assets; any other kind of asset should be ignored
\end{Verbatim}
% The classification procedure uses politicians characteristics to generate a prediction of the agricultural assets owned by the election's winner and their family, as well as the profession declared by the winner and their spouse. 
% This procedure uses the gpt-4o-mini model provided by OpenAI to analyze the relevant tables contained in the `winner_details` folder. Two classifications are conducted separately: winner's (and family members's) agrucultural asset ownership, and winner's and spouse's agricultural profession.

% ## Assets classification

% For each candidate, the table named `immovable_assets.csv` is read from disk. This comma-separated table contains information about each of the candidate's own agricultural assets, as well as their spouse's and their dependents's. At most three dependents are analyzed for each of the candidates. Each individual's column is analyzed separately to avoid feeding irrelevant tokens to the language model, which may lead to higher variance in the output. However, all the rows from the table are included, as this helps preserve additional examples of the format by which each asset is listed in each sub-category. The system prompt includes instructions to focus specically on the row that indicates agricultural assets. This prompt also contains instructions about the return format, which is to be a JSON blob containing two fields: "explanation" and "owns_agricultural_assets". This is also achieved using the parser tool of the completions endpoint in the OpenAI API, which allows for the specification of response schemas in JSON format using Pydantic. The "explanation" field is used to generate intermediary tokens that link the input to the classification of ownership of agricultural assets; in this way, there is a lower chance of allucination as these intermediary tokens must be based on the original input table and will determine the classification token later in the generation of the JSON blob. The classification is represented by a boolean integer, with 0 representing no detected ownership of agricultural assets and 1 representing a detection. A temperature of 0 is used for inference, as this reduces variance in the generation and helps reduce hallucinations.
\end{comment}

%\footnote{\red{Ask Rodrigo to add details on the iterative process we followed to obtain no uncertainty on classification.}}  
%\footnote{For 50 winning politicians out of N the website does not provide any information. } 




\subsection{Politicians' Actions}
\subsubsection{Court Cases\label{ssec:courts}} We gathered data from the SHRUG \citep{shrug_judicial} database, which contains all judicial case records from the Indian eCourts platform. From this database, we extracted all judicial cases in which there was a violation of the Air Pollution and Control Act of 1981 which include crop burning offenses from 2012 to 2018. 
%
Because these cases only report the district and the month in which they were filed, we match them to our sample at the district and month-year level from September 2012 to December 2018, corresponding to the start of our main sample period and the end of the available court data, respectively.



\subsubsection{DISHA Document Collection}\label{ssec:disha_app}

Our objective is to retrieve, for each district in our sample, any document or webpage that relates to DISHA meetings, and to match each document to a district and a month-year, using the portal that hosts it and the date recorded inside it. Because DISHA proceedings are posted on separate district government websites, with no common catalogue and no standard address, collecting them at scale requires a purpose-built web-crawler rather than a manual search. The crawl proceeds in two stages: a fast discovery stage that locates the pages where DISHA documents are likely to live, and a content-aware stage that follows those pages to the documents themselves.

\paragraph{Candidate websites.} We exploit the fact that Indian district portals, though not standardized, follow a small number of recurring conventions. For each of the four states in our sample, Punjab, Haryana, Uttar Pradesh, and Bihar, and each district within them, we generate candidate addresses by combining the district name with the sixteen URL templates listed at the end of this section, observed across the \texttt{.nic.in} and \texttt{.gov.in} domains; they include a dedicated \texttt{/disha/} path, a \texttt{/document-category/proceeding/} path, and a notices page queried for the term ``disha''. For multi-word district names we add a spelling variant that removes the hyphens, so that a district such as Charkhi-Dadri is reached whether or not the website encodes the hyphen. This yields sixteen candidate entry points for a single-word district name and up to twenty-four for a hyphenated one. Haryana adds four state-level addresses, because its rural development portal posts DISHA material of its own; no other state in the sample does. We then retain only the addresses that return a valid page.

%\paragraph{Parallel probing.} Rather than requesting these addresses one after another, which would stall on the many that do not resolve, we probe them concurrently. A pool of eight worker threads issues simultaneous requests with a short seven-second timeout, and we retain only the addresses that return a valid page. This parallel probe turns what would otherwise be a long sequence of timeouts into a single fast pass, and it is what makes crawling all districts feasible.

\paragraph{Content-aware, multilingual spidering.} We then read these pages and follow their internal structure to the documents. We parse each page's markup and extract two kinds of links: direct links to files, identified by their PDF, DOC, or DOCX extension and including documents embedded through inline frames; and navigation links whose visible text or address contains a DISHA keyword. The keyword set is deliberately multilingual and includes the Hindi spelling of DISHA alongside English cues such as ``Proceeding'', ``Minutes'', ``Agenda'', and ``District Development Coordination''. Because many websites do not list the documents on the landing page, but behind a ``DISHA'' or ``Proceedings'' tab, or on the later pages of a paginated archive, the crawler follows these keyword links one level deeper, restricting itself to the same website, and repeats the extraction there. This one-step spidering is what allows it to reach documents that a fixed-address scrape would miss, while the depth and same-domain limits keep the navigation bounded and prevent the crawler from wandering across the wider web. The files themselves are then fetched from wherever the page points, which in practice is almost always the government content delivery network on which the district portals store their documents.

\paragraph{Downloading.} Each discovered document is downloaded with safeguards. Before downloading, we issue a lightweight header request to read the file size and skip anything above fifty megabytes; we stream the remaining files in small chunks, enforcing the size cap as we go; we retry once, after a brief pause, on a timeout; and we identify with a standard browser agent and wait one second between downloads to avoid overloading the host. Every document is saved under a structured name that records its state, district, and the meeting date parsed from its link text where one is present, and every attempt, successful or not, is written to a log with its address, status, size, and extracted date.

\paragraph{Coverage.} We attempted 152 districts across the four states and found a live public page for 142 of them; the remaining ten had no publicly accessible page as of the collection date. Eight of the 142 yielded no documents, so 134 districts contributed files, together with the Haryana state-level page. The crawl made 3,072 download attempts, of which 3,059 succeeded, twelve timed out, and one was skipped for exceeding the size cap. The successful downloads comprise 1,251 files in Uttar Pradesh, 1,094 in Bihar, 393 in Haryana, and 321 in Punjab, totalling 6,289 megabytes. Because a small number of documents were saved under a name already used within the same district, and were therefore overwritten, the corpus on disk consists of 3,048 distinct documents.

\paragraph{Matching and coding.} We match each downloaded document to a district and a month-year, using the district portal that hosts it and the date recorded inside the document. Each document is then read and coded in three passes: one recovering the date of the document, one assigning a document type among nine genres, including meeting record, order, report, forward-looking plan, budget table, activity table and procedural instruction, and one assigning a primary topic together with separate relevance scores for agriculture, farmer protest, and air pollution. Because the crawl is deliberately inclusive, capturing every candidate document rather than risking a missed proceeding, no filtering takes place at collection. Documents that are not genuine DISHA proceedings are identified at the coding stage and excluded there, by document type or by topic, depending on the outcome being constructed.

Each district is included in the analysis sample starting from the month-year of its own first available document.\todo{I think we should also give examples of documents or more details on the documents in this section}

\paragraph{URL templates.} Writing \texttt{\{d\}} for the lowercased district name, the sixteen templates are:

\begin{enumerate}
\item \texttt{https://\{d\}.nic.in/disha/}
\item \texttt{https://\{d\}.nic.in/document-category/proceeding/}
\item \texttt{https://\{d\}.nic.in/en/notice/proceeding-of-district-development-coordination-and-monitoring-committee-disha/}
\item \texttt{https://\{d\}.nic.in/en/past-notices/meeting-proceeding/}
\item \texttt{https://\{d\}.nic.in/notices/?s=disha}
\item \texttt{https://\{d\}.nic.in/category/disha/}
\item \texttt{https://\{d\}.gov.in/disha/}
\item \texttt{https://\{d\}.gov.in/notices/disha/}
\item \texttt{https://\{d\}.gov.in/document-category/proceeding/}
\item \texttt{https://\{d\}.nic.in/}
\item \texttt{https://\{d\}.gov.in/}
\item \texttt{https://\{d\}.nic.in/notices/}
\item \texttt{https://\{d\}.nic.in/en/notices/}
\item \texttt{https://\{d\}.nic.in/notice/}
\item \texttt{https://\{d\}.nic.in/documents/}
\item \texttt{https://\{d\}.nic.in/downloads/}
\end{enumerate}

The de-hyphenated spelling variant is generated for the first eight templates. The four Haryana state-level addresses are \texttt{haryanarural.gov.in/}, \texttt{haryanarural.gov.in/agenda-for-disha-meeting-on-11-02-2026/}, \texttt{haryana.gov.in/disha/}, and \texttt{haryana.gov.in/notices/?s=disha}.


 \subsection{AI disclosure}\label{ssec:ai}



We use large language models in two components of the data-construction process. No generative AI tools were used to formulate the research question, select the empirical strategy, estimate the statistical models, or interpret the results.

First, we used GPT API to classify politicians’ self-reported profession descriptions obtained from Myneta.info. Each description was evaluated using a fixed prompt asking whether it referred to an agriculture-related profession that we then manually validate for a small sample. More informaiton about this process is present in Appendix Section \ref{sssec:myneta_app}.

Second, we used Anthropic’s Claude Code agents to construct a dataset from publicly available District Development Coordination and Monitoring Committee (DISHA) documents that were processed separately using OCR, including general-text and Punjabi/Gurmukhi OCR when necessary. Claude reviewed the OCR-extracted text in batches and applied predefined and fixed coding rubrics to: (1) identify each document’s primary topic; (2) assign relevance scores for agriculture, the 2020--2021 farmers’ protests and agricultural reforms, and air pollution or crop-residue burning; (3) infer the document or meeting date, prioritizing the meeting date over drafting or publication dates; and (4) classify the document genre, including meeting summaries, reports, orders, policy plans, instructions, and budget or planning tables. More information about data generation and identification of principal variables is present in Appendix Section \ref{ssec:disha_app}



\FloatBarrier
\section{Politicians' Incentives to Control Fires}
This section presents additional results for the main analysis. First, we present descriptives to the main variables (Table \ref{tab:_desc_stats_main}), including the seasonality of fires and wind direction (Figure \ref{fig:overtime}). Second, we present the distribution of observations by relative time period from treatment (Figure \ref{fig:reltime}), and an illustration of the empirical strategy underlying the pro-sociality test (Figure \ref{fig:prosociality_ex}). We then turn to results, showing that our main estimates, presented in Figure \ref{fig:main_panel}, are robust to alternative definitions of treatment and outcome, to alternative samples, and to alternative explanations.
%
We first vary the definition of treatment (Table \ref{tab:main_by_treat}) and the dependent variable (Table  \ref{tab:main_bydv}), then we subset by year (Table \ref{tab:by_year}) and province (Table \ref{tab:by_province}). We next report whether the effect is driven by farmers' prosociality, rather than politicians (Table \ref{tab:prosociality}). We then replicate our main result using downwind and upwind area, rather than population, to define treatment (Figure \ref{fig:main_panel_area}). Next, we examine heterogeneity by interacting treatment with constituency-level agricultural characteristics (Figure \ref{fig:es_rice}). Finally, we report the event-study specification as in equation \ref{eq2} (Figure \ref{fig:es_main_rotated_panel}a), after de-trending for the pre-
treatment period linear trend (Figure \ref{fig:es_main_rotated_panel}b), along with the Honest DiD sensitivity analysis by \citet{rambachan2023more} (Figure \ref{fig:es_main_rotated_panel}c). 


\subsection{Empirical Strategy: Event-Study}\label{ssec:empirical1_app}

We follow recent applied work in difference-in-differences estimation \citep{rambachan2023more} adopting a neutral transformation to de-trend the event studies \citep{dobkin2018economic, dustmann2022reallocation}. This approach accounts for the possibility that treated and control units' outcomes may have evolved along different paths in the pre-treatment period. In particular, we follow \citet{dustmann2022reallocation} and fit a linear time trend using our estimated coefficients from the pre-treatment period, anchored at zero in the omitted category, and then extrapolate this trend into the post-treatment period. This approach accounts for potential violations of the parallel trends assumption in settings, such as ours, where the differential trend is linear. 

%By plotting the deviations of our estimated coefficients from these predicted trends, we can assess whether deviations from the linear trends after treatment implementation are large enough to explain away the treatment effects we observe, i.e. whether the breaks in outcome trends are substantial for the estimated treatment effect on fires to be reduced to zero. 
\clearpage 
\subsection{Descriptives}


 \begin{table}[htbp!]
     \centering
     \caption{Descriptive Statistics \label{tab:_desc_stats_main}}
     \resizebox{0.8\textwidth}{!}{
      \begin{tabular}{lrrrrrr}
\toprule
 & Mean & SD & Min & Max & Observations & Unique Obs.\\
\midrule
Grid &  &  &  &  & 36,758,771 & 75,326\\
Year &  &  &  &  & 36,758,771 & 11\\
Month &  &  &  &  & 36,758,771 & 12\\
Assembly Constituency (AC) &  &  &  &  & 36,758,771 & 797\\
Province &  &  &  &  & 36,758,771 & 4\\
Number of Fires & 0.156 & 1.005 & 0 & 125 & 36,758,771 & 39\\
Down $>$ Up & 0.103 & 0.303 & 0 & 1 & 36,758,771 & 2\\
Average Wind Speed & 1.342 & 0.368 & 0.419 & 3.175 & 36,758,771 & 491,443\\
Wind Direction & -18.567 & 94.096 & -180 & 180 & 36,758,771 & 497,927\\
Rice Production & 0.452 & 0.498 & 0 & 1 & 36,758,771 & 2\\
\bottomrule
\end{tabular}

      }
      \begin{minipage}{\textwidth}
         {\small \textit{Notes: } This table shows descriptive statistics of the main variable for the politician's incentive analysis }
     \end{minipage}
 \end{table}




\begin{figure}[htbp!]
    \centering
\caption{Patterns of Seasonality in Fires Data}
    \label{fig:overtime}
    \includegraphics[height=4cm]{figures/monthly_fires.png}
    \includegraphics[height=4cm]{figures/monthly_wind_direction.png}
\begin{minipage}{\textwidth}
    {\small \textit{Notes: }(Left) Total number of fires per month. (Right) District-average actual and 10 year rolling average wind direction. All data refer the states in our sample only ---Bihar, Haryana, Punjab, and Uttar Pradesh--- in the period September 2012 - August 2022.}
\end{minipage}
\end{figure}


	\begin{figure}[htbp!]
    	\centering
    	\caption{Rice Production in Northern Indian Subcontinent\label{fig:map_crops} } %%% Anzony -- We need to show rice production in quantiles and indicate in the figure notes we are spplitting in the notes that we are using above/below median analysis
          \includegraphics[height=6cm]{figures/rices_grids_150dpi_q75.pdf} 
          \begin{minipage}{\textwidth}
    {\small \textit{Notes: } Distribution of rice production (in tons) in 2010 across grids, classified by quintile of production. }
\end{minipage}

   % \hfill
   %\begin{minipage}{0.45\textwidth}
   %    \subfloat[\centering \textbf{Wheat crops}]{\includegraphics[width=\textwidth]{figures/wheat_crops.png} }
   %\end{minipage}
    \end{figure}


    \begin{figure}[htbp!]
        \centering
        \caption{Number of observations per relative time period from treatment}
        \label{fig:reltime}
        \includegraphics[height=6cm]{figures/downup_evtime_hist_new.png}
    \end{figure}

    \begin{figure}[htbp!]
        \centering
        \caption{Placebo border (circle) against the true constituency border for one grid}
        \label{fig:prosociality_ex}
        \includegraphics[width=0.5\linewidth]{figures/acs_grids_radius12km.png}
        
    \end{figure}
% \begin{minipage}
%             \footnote{Note: This is Wazirgan Assembly located in Bihar}
%         \end{minipage}

%%%%%%%%%%%%%%%%%%Downup Population%%%%%%%%%%%%%%%
% \begin{table}[htbp!]
%     \centering
%     \caption{Effect of exposing the population \label{tab:main_population}}
%     \resizebox{0.8\textwidth}{!}{
%      \input{tables/_app_main_did_downup_pop_ac_edit}
%      }
%      \begin{minipage}{\textwidth}
%         {\small \textit{Notes: }Results from Equation \ref{eq1} varying the treatment definition. The treatment indicator is a dummy equal to one in months and grids from which most of the district population is exposed to smoke. Controls include average wind direction and wind speed in each grid-cell and month. 
%         The table reports the mean of the dependent variable in each year for grids that are upwind with respect to borders of the constituency. All regressions control for average wind speed and direction. Standard errors are clustered at the grid and AC times month-year level. }
%     \end{minipage}
% \end{table}

%\newpage
\FloatBarrier
\subsection{Translating Fire Reductions into Air Quality and Mortality Estimates}\label{sec:mortality_estimates}

To translate our estimated fire reductions into air quality and mortality impacts, we draw on \citet{lan2022air}, who use the GEOS-Chem adjoint atmospheric chemistry and transport model to quantify the sensitivity of India-wide population-weighted PM2.5 exposure to crop residue burning emissions from 2003 to 2019. Their approach computes how PM2.5 exposure responds to marginal changes in fire emissions at each location and time, making it well-suited to deriving the air quality consequences of the fire reductions we estimate. The model is validated against Central Control Pollution Board ground observations and NASA MODIS satellite-derived aerosol optical depth. 

This translation rests on assuming proportionality between fire counts and emissions: a given percentage reduction in fire counts is treated as equivalent to the same percentage reduction in burning-attributable PM2.5 emissions and, in turn, in PM2.5 concentrations. We also assume that politicians the greater effort exerted by politicians in treated areas is a net effect. However, this could be compensated through reduced effort elsewhere in the constituency. As we note in the main text, this should thus be understood as a benchmark for what would be achievable under full increased political accountability, rather than as estimates of  net gains. We give estimates for Punjab, Haryana, and Uttar Pradesh, the three states which Lan et al. identify as collectively responsible for over two-thirds of India-wide air quality impacts from crop burning. Numbers would be slightly larger adding the fourth state in our analysis sample, Bihar.

For the PM2.5 estimate, Lan et al. report that crop burning across India contributes an annual mean population-weighted PM2.5 of 6.7 $\mu$g m$^{-3}$, of which 64 percent arises from post-monsoon rice burning in October and November. This seasonal contribution of 4.3 $\mu$g m$^{-3}$ to the annual mean, concentrated in approximately 60 days of burning activity, corresponds to an average burning-season PM2.5 from crop residue of approximately 26 $\mu$g m$^{-3}$ across the India-wide population (4.3 × 365/60). Applying our estimated 13 percent fire reduction yields a burning-season PM2.5 reduction of approximately 3.4 $\mu$g m$^{-3}$ at the India national level. Residents of Punjab, Haryana, and Uttar Pradesh experience substantially higher burning-season concentrations due to their proximity to the source. Lan et al. document that burning-attributable PM2.5 in northern India during the burning season reaches 15 to 45 times the WHO annual mean guideline of 5 $\mu$g m$^{-3}$, implying local burning-season concentrations in the range of 75–225 $\mu$g m$^{-3}$, and that Delhi alone consistently exceeded 120 $\mu$g m$^{-3}$ of PM2.5 averaged over the post-monsoon period across their full sample. A 13 percent reduction applied to these local levels implies reductions of roughly 10–29 $\mu$g m$^{-3}$. 

For the mortality estimate, Lan et al. report 69,000 premature deaths per year attributable to ambient PM2.5 exposure from crop burning across India (95\% CI: 57,000–80,000), with Punjab, Haryana, and Uttar Pradesh accounting for approximately 70 percent of this burden, or roughly 48,000 deaths annually. Applying a 13 percent proportional reduction to 48,000 yields approximately 6,200 fewer premature deaths per year. This figure is best understood as a full-accountability benchmark: it represents the hypothetical mortality reduction if political accountability suppressed fires by 13 percent uniformly across all grids and months throughout these states. In practice, our treatment variable (smoke directed toward politicians' own constituents) applies to approximately half of all grid-month observations, given a mean treatment rate of 50 percent in our sample (Table~\ref{tab:_desc_stats_main}). Weighting by the realized treatment share yields an estimate of approximately 3,100 averted deaths per year from the accountability mechanism as currently operating. The full-accountability figure is a benchmark for the potential public health impact of political accountability when fully operative, and when the increase in effort by politicians is not compensated by a decrease in effort elsewhere. Another caveat is that Lan et al.'s mortality estimates are averaged over 2003–2019; more recent years, which see higher burning intensity, would imply a larger baseline and correspondingly larger potential reductions.



\newpage 
\subsection{Robustness}

%%%%%%%%%%%%%%%%%Robustness of Rice Moderator and Rotation %%%%%%%%%%%%%%%


%%%%%%%%%%%%%%%%%Robustness by Treatment definition%%%%%%%%%%%%%%%

 \begin{table}[htbp!]
     \centering
     \caption{Effect of constituency exposure to smoke on number of fires, DID specification \label{tab:main_DID}}
     \resizebox{0.8\textwidth}{!}{
    {
\def\sym#1{\ifmmode^{#1}\else\(^{#1}\)\fi}
\begin{tabular}{l*{4}{c}}
\hline
            &\multicolumn{1}{c}{(1)}         &\multicolumn{1}{c}{(2)}         &\multicolumn{1}{c}{(3)}         &\multicolumn{1}{c}{(4)}             \\
            & \multicolumn{4}{c}{Number of Fires (in 1,000 units)} \\ \hline

Down $>$ Up &       3.027         &      -2.412         &     -37.125\sym{***}&      -5.435\sym{***}\\
            &     (2.085)         &     (1.564)         &     (2.397)         &     (1.644)         \\
\hline
Observations&  36,758,771         &  36,758,771         &  36,758,771         &  36,758,771         \\
N Assembly Constituencies&         797         &         797         &         797         &         797         \\
Month-Year $\times$ Cohort FE&           N         &           Y         &           Y         &           N         \\
Grid $\times$ Cohort FE&           N         &           N         &           Y         &           Y         \\
AC $\times$ Month-Year $\times$ Cohort FE&           N         &           N         &           N         &           Y         \\
Mean DV     &      92.122         &      92.122         &      92.122         &      92.122         \\
\hline
\end{tabular}
}

      }
     \begin{minipage}{\textwidth}
         {\small \textit{Notes: } Results from Equation~\ref{eq1} of the effect of constituency exposure to smoke on the number of fires (in thousands) in a grid-cell and month. The treatment indicator \textit{Down$>$Up} is a dummy equal to one in months and grids from which most of the assembly constituency's population\ is exposed to smoke (i.e. when the constituency's population in the downwind area is larger than the population in the  upwind area from the grid). The table reports the mean of the dependent variable when treated grids are upwind. All regressions control for average wind speed and direction. Standard errors are clustered at the grid and assembly constituency $\times$ month-year level. }
     \end{minipage}
 \end{table}


 \begin{table}[htbp!]
     \centering
     \caption{Effect of constituency exposure to smoke on number of fires, robustness by treatment indicator \label{tab:main_by_treat}}
     \resizebox{0.8\textwidth}{!}{
      \begin{tabular}{lcccccc}
\hline
& (1) & (2) & (3) & (4) & (5) & (6)\\
& \multicolumn{6}{c}{Number of Fires (in 1,000 units) }\\
\hline

Down$>$Up (Pop)&      1.6119         &                     &                     &                     &                     &     -4.4245\sym{*}  \\
            &    (1.3411)         &                     &                     &                     &                     &    (2.5064)         \\
Down$>$Up (Area)&                     &      2.8520\sym{**} &                     &                     &                     &                     \\
            &                     &    (1.3289)         &                     &                     &                     &                     \\
Down$>$Up by 1std&                     &                     &      0.7269         &                     &                     &                     \\
            &                     &                     &    (1.5769)         &                     &                     &                     \\
Downwind over total Population&                     &                     &                     &     -0.0188\sym{***}&                     &                     \\
            &                     &                     &                     &    (0.0050)         &                     &                     \\
Down-Up Percent&                     &                     &                     &                     &     -0.0096         &                     \\
            &                     &                     &                     &                     &    (0.0134)         &                     \\
Down$>$Up (Pop) $\times$ Above Median Rice Production&                     &                     &                     &                     &                     &     -2.0705         \\
            &                     &                     &                     &                     &                     &    (5.6424)         \\
\hline
Observations&  36,878,861         &  36,878,861         &  36,878,861         &  36,740,247         &  36,878,861         &  36,740,247         \\
N Assembly Constituencies&         808         &         808         &         808         &         808         &         808         &         808         \\
Grid $\times$ Cohort FE&           Y         &           Y         &           Y         &           Y         &           Y         &           Y         \\
Assembly $\times$ Month-Year $\times$ Cohort FE&           Y         &           Y         &           Y         &           Y         &           Y         &           Y         \\
Mean DV     &      91.996         &      91.996         &      91.996         &      91.996         &      91.996         &      91.996         \\
\hline
\end{tabular}

      }
     \begin{minipage}{\textwidth}
         {\small \textit{Notes: }Results from Equation~\ref{eq1} considering different treatment indicators: a dummy for whether the downwind population is larger than the upwind population (as in the main specification) in Column 1; a dummy for whether the downwind area, rather than the population, is larger than the upwind area in Column 2; the a dummy for whether the downwind area is larger than the upwind area by one standard deviation or more in column 3; the share of downwind population over total population in Column 4; the percentage difference between down and upwind area in Column 5. The dependent variable is the number of fires in a 5 squared km grid-cell and month. 
         The table reports the mean of the dependent variable when grids are upwind. All regressions control for average wind speed and direction. Standard errors are clustered at the grid and assembly constituency $\times$ month-year level. }
     \end{minipage}
 \end{table}


%%%%%%%%%%%%%%%%%Robustness by DV definition%%%%%%%%%%%%%%%
 \begin{table}[htbp!]
     \centering
     \caption{Effect of constituency exposure to smoke on number of fires, robustness by dependent variable \label{tab:main_bydv}}
     \resizebox{0.8\textwidth}{!}{
      \begin{tabular}{lccc}
\tabularnewline \hline
& (1) & (2) & (3)\\
& Any Fire & Log (N) Fires & Mean Brightness\\
\hline

Down $>$ Up &     -0.0014\sym{***}&     -0.0019\sym{***}&     -0.4550\sym{***}\\
            &    (0.0003)         &    (0.0005)         &    (0.1098)         \\
\hline
Observations&  36,758,771         &  36,758,771         &  36,758,771         \\
N Assembly Constituencies&         797         &         797         &         797         \\
Grid FE $\times$ Cohort&           Y         &           Y         &           Y         \\
Assembly $\times$ Month-Year $\times$ Cohort FE&           Y         &           Y         &           Y         \\
Mean DV     &       0.037         &        0.04         &      12.416         \\
\hline
\end{tabular}

      }
      \begin{minipage}{\textwidth}
         {\small \textit{Notes: }Results from Equation~\ref{eq1} using alternative dependent variables: a dummy for whether any fire took place (Column 1), the natural logarithm of the number of fires plus 1 (Column 2) and the average brightness of fires in a grid-month -- brightness is equal to zero when there are no fires (Column 3).  
         The table reports the mean of the dependent variable when ever treated grids are upwind. All regressions control for average wind speed and direction. Standard errors are clustered at the grid and assembly constituency $\times$ month-year level. }
     \end{minipage}
 \end{table}

 %%%%%%%%%%%%%%%%%Robustness by functional form %%%%%%%%%%%%%%%
%\begin{table}[htbp!]
%    \centering
%    \caption{Effect of constituency exposure to smoke on number Fires, robustness by functional form \label{tab:main_poisson_negbin}}
%    \resizebox{0.8\textwidth}{!}{
%     \input{tables/_app_main_did_downup_area_ac_binomial_family_edited}
%     }
%     \begin{minipage}{\textwidth}
%        {\small \textit{Notes: }Results from Equation~\ref{eq1} using alternative functional forms: a Poisson to accomodate a count dependent variable, and a Negative Binomial to accomodate for the high skew of the dependent variable towards zero. All regressions include fixed-effects at grid and district times month-year level and control for average wind direction and wind speed in each grid and month. 
%        The table reports the mean of the dependent variable in each year for grids that are upwind with respect to borders of the constituency.  All regressions control for average wind speed and direction. Standard errors are clustered at the grid and AC times month-year level. }
%    \end{minipage}
%\end{table}

%%%%%%%%%%%%%%%%%Effect by Year%%%%%%%%%%%%%%%
 \begin{table}[htbp!]
     \centering
     \caption{Effect of constituency exposure to smoke on number of fires, robustness by year \label{tab:by_year}}
     \resizebox{\textwidth}{!}{
      \begin{tabular}{lcccccccccc} \hline
 & (1) & (2) & (3) & (4) & (5) & (6) & (7) & (8) & (9) & (10) \\
 & \multicolumn{10}{c}{Number of Fires (in 1,000 units)}\\
 & 2012/2013 & 2013/2014 & 2014/2015 & 2015/2016 & 2016/2017 & 2017/2018 & 2018/2019 & 2019/2020 & 2020/2021 & 2021/2022 \\ \hline

Down$>$Up (Pop)&       5.381         &       0.673         &       9.912\sym{*}  &      -2.800         &       1.007         &      -3.136         &     -12.267\sym{*}  &      -1.449         &      -4.900         &       1.682         \\
            &     (5.205)         &     (5.654)         &     (5.147)         &     (5.466)         &     (5.742)         &     (5.773)         &     (6.264)         &     (4.138)         &     (6.452)         &     (6.304)         \\
\hline
Observations&   5,552,535         &   5,876,789         &   6,007,388         &   6,194,026         &   6,280,088         &   6,250,728         &   6,479,638         &   6,352,381         &   6,417,605         &   6,201,762         \\
N Assembly Constituencies&         808         &         808         &         808         &         808         &         808         &         808         &         808         &         808         &         808         &         808         \\
Grid FE     &           Y         &           Y         &           Y         &           Y         &           Y         &           Y         &           Y         &           Y         &           Y         &           Y         \\
Assembly $\times$ Month-Year FE&           Y         &           Y         &           Y         &           Y         &           Y         &           Y         &           Y         &           Y         &           Y         &           Y         \\
Mean DV     &     151.139         &     141.492         &     145.732         &     151.219         &     180.108         &     151.672         &     162.715         &     104.452         &     168.627         &     175.076         \\
\hline
\end{tabular}

      }
      \begin{minipage}{\textwidth}
         {\small \textit{Notes: }Results from Equation \ref{eq1} subsetting by year, from each September to the following August. The table reports the mean of the dependent variable in each year when ever treated grids are upwind. All regressions control for average wind speed and direction. Standard errors are clustered at the grid and assembly constituency $\times$ month-year level. }
     \end{minipage}
 \end{table}

%%%%%%%%%%%%%%%%%Effect by Province%%%%%%%%%%%%%%% 
 \begin{table}[htbp!]
     \singlespacing \centering
     \caption{Effect of constituency exposure to smoke on number of fires, robustness by province \label{tab:by_province} }
     \resizebox{0.8\textwidth}{!}{
      \begin{tabular}{lcccc} \hline
 & (1) & (2) & (3) & (4) \\
 & \multicolumn{4}{c}{Number of Fires (in 1,000 units)}\\
 & Bihar & Haryana & Punjab & Uttar Pradesh \\ \hline

Down$>$Up (Pop)&      -0.715         &      12.002\sym{*}  &     -16.098\sym{**} &      -8.333\sym{***}\\
            &     (2.589)         &     (6.282)         &     (7.314)         &     (1.455)         \\
\hline
Observations&   6,140,896         &   4,013,718         &   4,942,212         &  21,661,945         \\
N Assembly Constituencies&         237         &          84         &         107         &         380         \\
Grid FE  $\times$ Cohort&           Y         &           Y         &           Y         &           Y         \\
Assembly $\times$ Month-Year $\times$ Cohort FE&           Y         &           Y         &           Y         &           Y         \\
Mean DV     &      55.155         &     103.195         &     348.495         &      44.393         \\
\hline
\end{tabular}

      }
      \begin{minipage}{\textwidth}
         {\small \textit{Notes: } Results from Equation \ref{eq1} subsetting by province. The table report the mean of the dependent variable in each province when ever treated grids are upwind. All regressions control for average wind speed and direction. Standard errors are clustered at the grid and assembly constituency $\times$ month-year level. }
     \end{minipage}
 \end{table}


\begin{figure}[htbp!]
    \caption{Effect of constituency exposure on number of fires }
    \label{fig:es_main_rotated_panel}
    \centering
    % --- Row 1: Population-weighted (columns carry the referenceable letters) ---
    \par\smallskip
    \subcaptionbox{Event-Study\label{fig:col1}}[0.3\textwidth]{%
        \includegraphics[width=\linewidth]{figures/stacked_event_study_pop_5pre_rural_1_ori.png}}%
         \subcaptionbox{De-Trended Event-study \label{fig:es_main_rotated_app}}[0.3\textwidth]{%
    \includegraphics[width=\linewidth]{figures/stacked_event_study_pop_5pre_rural_1_rotated.png}
  }
    \subcaptionbox{Honest DiD\label{fig:col2}}[0.3\textwidth]{%
        \includegraphics[width=\linewidth]{figures/stacked_event_study_pop_5pre_rural_1_rot_honest2.png}}

    \par\bigskip
    \begin{minipage}{0.9\textwidth}
        \footnotesize \textit{Notes:} (a) Results from Equation \ref{eq2}. (b) Results from Equation \ref{eq2} after de-trending for the pre-treatment period linear trend following \citet{dustmann2022reallocation}. (c) Sensitivity analysis for the average post-treatment effect allowing for violations of parallel trends \citep{rambachan2023more}.
        The reference category is $t=-1$,i.e. the month before the switch.  All regressions include grid $\times$ cohort and assembly constituency $\times$ month-year $\times$ cohort fixed effects and control for average wind speed and direction. Standard errors are clustered at the grid $\times$ cohort and assembly constituency $\times$  month-year $\times$ cohort level.
        %Sensitivity analysis is conducted using the HonestDiD approach of \citet{rambachandran2023} for the average of the post treatment effects allowing for violation of parallel trends as the maximum violation in pre treatment period.. The bounds are computed for the mean of the post-treatment coefficients. Panel (a) presents the HonestDiD bounds for the rotated event study, panel (b) shows the original event study estimates, and panel (c) presents the HonestDiD bounds for the original (non-rotated) event study.
    \end{minipage}
\end{figure}





\begin{figure}
    %%% Main Table %%%%%%%
    \centering
    \caption{Effect of constituency exposure to smoke on number of fires, by area }\label{fig:main_panel_area}
    %%% Main Table and Event Study side by side %%%%%%%
    \vspace{-0.3cm}
    \singlespacing
  % (a) table
  %\subcaptionbox{Main Table\label{tab:main_area}}[0.48\textwidth]{%
    %\begin{adjustbox}{max width=\linewidth}
    %  \input{tables/main_did_downup_ac_rural_acarea} % tabular only
   % \end{adjustbox}
  %}
  % (b) plot
 % \subcaptionbox{Event-study relative to exposure\label{fig:es_main_area_rotated}}[0.48\textwidth]{%
    \includegraphics[width=\linewidth]{figures/stacked_event_study_5pre_rural_1_ori.png}
 % }
  % (c) bottom table
 % \subcaptionbox{Effect of constituency and district exposure to smoke on N Fires\label{tab:main_ac_area_distr}}[\textwidth]{%
 %   \begin{adjustbox}{max width=.7\linewidth}
 %     {\def\sym#1{\ifmmode^{#1}\else\(^{#1}\)\fi}
\begin{tabular}{l*{5}{c}}
\hline
 & (1) & (2) & (3) & (4) & (5) \\
 & \multicolumn{5}{c}{Number of Fires (in 1,000 units)} \\ \hline

Down$>$ Up Bureaucrat&     -24.369\sym{***}&     -22.156\sym{***}&      -2.506\sym{**} &      -7.114\sym{***}&     -18.051\sym{***}\\
            &     (2.958)         &     (1.355)         &     (1.018)         &     (1.002)         &     (1.156)         \\
Down$>$ Up Politician (Pop)&       7.819\sym{**} &      -7.902\sym{***}&       4.762\sym{***}&     -10.448\sym{***}&      -6.530\sym{***}\\
            &     (3.110)         &     (2.077)         &     (1.711)         &     (2.059)         &     (2.170)         \\
Down$>$ Up Pol. (Pop) $\times$ Down$>$ Up Bur.&     -21.678\sym{***}&       3.514         &       1.102         &       9.958\sym{***}&       1.616         \\
            &     (3.991)         &     (2.462)         &     (1.885)         &     (2.003)         &     (2.261)         \\
\hline
Observations&  36,758,771         &  36,758,771         &  36,758,771         &  36,758,771         &  36,758,531         \\
N Assembly Constituencies&         808         &         808         &         808         &         808         &         808         \\
N Districts &         160         &         160         &         160         &         160         &         160         \\
Month-Year FE&           N         &           Y         &           N         &           N         &           N         \\
AC FE       &           N         &           Y         &           N         &           N         &           N         \\
AC $\times$ Month-Year FE&           N         &           N         &           Y         &           Y         &           N         \\
District $\times$ Month-Year FE&           N         &           N         &           N         &           N         &           Y         \\
Grid FE     &           N         &           N         &           N         &           Y         &           Y         \\
Mean DV     &      86.353         &      86.353         &      86.353         &      86.353         &      86.353         \\
\hline
\end{tabular}
}

 %   \end{adjustbox}
 % }
  \vspace{-0.3cm}
     \begin{minipage}{\textwidth}
         {\small \textit{Notes: } 
         %%% Main Table %%%%%%%
        %\textit{(a)} Results from Equation~\ref{eq1} of the effect of constituency exposure to smoke on the number of fires (in thousands) in a grid-cell and month. The treatment indicator \textit{Down$>$Up} is a dummy equal to one in months and grids from which the majority of the assembly constituency's area is exposed to smoke with respect to the grid.  
         %%% Main Event Study %%%%%%%
          Results from Equation \ref{eq2}.
         % after de-trending for the pre-treatment period linear trend following \citet{dustmann2022reallocation}. 
          The treatment indicator \textit{Down$>$Up} is a dummy equal to one in months and grids from which the majority of the assembly constituency's area is exposed to smoke with respect to the grid.  
         The reference category is t=-1, i.e. the month before of the switch. We use a stacked diff-in-diff to generate cohorts based on switching date. We restrict the analyses to a 6-months window around the switch to maximize support. The control grids are never treated and upwind grids around the switching date. 
          The regression includes grid and assembly constituency $\times$ month-year $\times$ cohort fixed effects and controls for average wind speed and direction. Standard errors are clustered at the grid and district $\times$ month-year level.
          %%% Free Riding %%%%%%%
          %\textit{(c)} Results from Equation \ref{eq1} with the additional interaction with a dummy (\textit{Down $>$ Up Bureaucrat}) equal to one in months and grids from which most of the district area is exposed to smoke, i.e. the area of the district downwind is larger than the area upwind from the grid.
     % The tables report the mean of the dependent variable in grids that are upwind with respect to the borders of both the constituency and the district.  
       }
     \end{minipage}
\end{figure}




%%%%%%%%% BY RICE AREA %%%%%%%%%
 \begin{figure}[!h]
 \centering
\caption{Effect of constituency exposure to smoke on number of Fires, by constituency with rice production}\label{fig:es_rice}
{%
    \includegraphics[width=0.48\linewidth]{figures/stacked_event_study_pop_5pre_rural_riceP_ori.png}
  }
  \hfill
       \begin{minipage}{\textwidth}
        {\small \textit{Notes: } Results as in Figure \ref{fig:es_main_rotated} with additional interactions by constituency agricultural characteristics. The figure plots the interaction term between treatment  and an indicator for constituencies with above-median rice production. The reference category is t=-1, i.e. the month before of the switch. All regressions control for average wind speed and direction. Standard errors are clustered at the grid and assembly constituency $\times$ month-year $\times$ corhot level. }
     \end{minipage}
 \end{figure}

 %%%%%%%%%% Farmers' Pro-Sociality %%%%%%%%%%%%
\begin{table}[htbp!]
   \caption{ Effect of exposing the population in a farmer placebo area to smoke on N fires \label{tab:prosociality}} 
   \centering
   \singlespacing
   \resizebox{0.75\textwidth}{!}{
   \begin{tabular}{lccc}
\hline
& \multicolumn{3}{c}{Number of Fires (in 1,000 units)}\\
& Full Sample & Treated for Politicians & Control for Politicians \\
& (1) & (2) & (3) \\
\midrule

Down$>$Up (Placebo)&       4.488\sym{***}&      11.967\sym{***}&       5.284\sym{*}  \\
            &     (1.415)         &     (2.744)         &     (3.203)         \\
\hline
Observations&  35,307,171         &  17,539,917         &  17,691,625         \\
N Assembly Constituencies&         808         &         808         &         808         \\
Grid $\times$ Cohort FE&           Y         &           Y         &           Y         \\
AC $\times$ Month-Year $\times$ Cohort FE&           Y         &           Y         &           Y         \\
Mean DV     &      90.983         &      90.983         &      90.983         \\
\hline
\end{tabular}

   }   
   \begin{minipage}{\textwidth}
    {\small \textit{Notes: }Results from Equation \ref{eq1} considering a different treatment indicator: a dummy equal to one in months and grids from which the population downwind is larger than the population upwind in a 13km radius circle around each grid (median size of an assembly constituency). %The placebo area is aimed at capturing farmers’ prosocial incentives to avoid pollution, independently on the incentives faced by politicians. ANZONY LEFT TO ADD
    Column 1 considers the entire sample; Column 2 restricts to areas that are treated both for the farmer and the politician, i.e. downwind both with respect to the placebo area and the constituency; Column 3 restricts to areas that are treated for the farmer only, i.e. downwind with respect to the placebo area but upwind to the constituency. 
    The table reports the mean of the dependent variable when treated grids are upwind. All regressions control for average wind speed and direction. Standard errors are clustered at the grid and assembly constituency $\times$ month-year level.  }
     \end{minipage}
\end{table}


\newpage
\FloatBarrier
\section{Lobbying}

This section presents additional results for the lobbying analysis. First, we present descriptive statistics for the main variables (Table \ref{tab:_desc_incentive_main}) and illustrate the construction of the empirical strategy (Figures \ref{fig:protest_by_time} and \ref{fig:protest_construction}). Second, we present the de-trended event-study along with the Honest DiD sensitivity analysis following \citet{rambachan2023more} (Figure \ref{fig:protest_main_rotated_panel}). Lastly, we report heterogeneity by constituency exposure to smoke (Table \ref{tab:protest_downup}).

\subsection{Descriptives}

%%%

% \subsection{Descriptives}
 \begin{table}[htbp!]
     \centering
     \caption{Descriptive Statistics \label{tab:_desc_incentive_main}}
     \resizebox{0.8\textwidth}{!}{
      \begin{tabular}{lrrrrrr}
\toprule
 & Mean & SD & Min & Max & Observations & Unique Obs.\\
\midrule
Grid ID &  &  &  &  & 63,240,134 & 76,800\\
Year &  &  &  &  & 63,240,134 & 7\\
Month &  &  &  &  & 63,240,134 & 12\\
Assembly Constituency (AC) &  &  &  &  & 63,240,134 & 808\\
Province &  &  &  &  & 63,240,134 & 4\\
Protest Area &  &  &  &  & 63,240,134 & 1,654\\
Cohort &  &  &  &  & 63,240,134 & 19\\
Legislature &  &  &  &  & 63,240,134 & 5\\
Relative year & -1.71 & 1.478 & -4 & 1 & 63,240,134 & 6\\
Protest & 0.002 & 0.044 & 0 & 1 & 63,240,134 & 2\\
Number of Fires & 143.98 & 968.654 & 0 & 42,000 & 63,240,134 & 39\\
High Rice production (AC level) & 0.451 & 0.498 & 0 & 1 & 63,240,134 & 2\\
\bottomrule
\end{tabular}

      }
      \begin{minipage}{\textwidth}
         {\small \textit{Notes: } This table shows descriptive statistics of the main variable for the protest analysis }
     \end{minipage}
 \end{table}
 

\subsection{Empirical Strategy}
 %%%%%%%%%% Protest Desc %%%%%%%%%
%\begin{table}[htbp!]
%    \centering
%    \caption{Farmers' Protests: Summary Statistics }
%    \label{tab:protest_desc}
%    \resizebox{0.8\textwidth}{!}{\centering\input{tables/_app_desc_10_5km_protest}}
%\end{table}

 %%%%%%%%%%%% Protest across time and space %%%%%%%%%%
%\begin{figure}[htbp!]
%    \centering
%    \caption{Protest occurrence over time}
%    \label{fig:panelview_protest}
%    \includegraphics[width=0.5\linewidth]{figures/protest_panelview_aclvl.png}
%    \begin{minipage}{\textwidth}
%        {\small \textit{Notes: }The figure shows the temporal distribution of treatment assignment, where treatment is defined as the observed protest event. \red{to remove?} }
%    \end{minipage}
%\end{figure}

\begin{figure}[htbp!]
    \centering
    \caption{Number of protests over time}
    \label{fig:protest_by_time}
    \includegraphics[width=0.5\linewidth]{figures/protests_monthly_bars.png}
    %\begin{minipage}{\textwidth}
    %    {\small \textit{Notes: }The histogram plots the distribution of protest occurrence in our sample from June 2020 to December 2021 by month. }
    %\end{minipage}
\end{figure}

% \begin{figure}[htbp!]
%     \centering
%     \caption{Protest occurrence by location}
%     \label{fig:map_protest}
%     \includegraphics[width=0.5\linewidth]{figures/5km_plot.png}
%     \begin{minipage}{\textwidth}
%         {\small \textit{Notes: }The map shows treated (blue) and control (white) grid-cells. A grid-cell is coded as treated if it falls within 5km-radius buffer areas around protest locations. Treatment begins in the month-year of the first protest within each buffer area. Control grid-cells are located outside all buffer areas.}
%     \end{minipage}
% \end{figure}

\begin{figure}[htbp!]
    \centering
    \caption{Buffer area against Assembly Constituency border}
    \label{fig:protest_construction}
    \includegraphics[width=0.5\linewidth]{figures/ac_protest_plot.png}
    \begin{minipage}{\textwidth}
         {\small \textit{Notes: }The map plots the borders of Jind Assembly constituency (black line). Three protests occurred in this constituency during the sample period. Around each protest location, we construct 5km-radius buffer areas (blue shaded) around the protest location (black dot). We classify grid-cells falling within these buffer areas as treated, regardless of whether they fall inside or outside the assembly constituency boundaries. Grid-cells located in overlapping buffer areas (darker shaded) are assigned to the buffer associated with the earliest protest. }
    \end{minipage}
\end{figure}

\clearpage
\subsection{Robustness}

\begin{figure}[htbp!]
    \caption{Effect by relative time period from protest}
    \label{fig:protest_main_rotated_panel}
    \centering
    \subcaptionbox{Event-Study}[0.3\textwidth]{%
        \includegraphics[width=\linewidth]{figures/_app_17_5km_fe12_evst_all_rural_fe03_1_ori.png} 
         }%
        \subcaptionbox{De-Trended Event-Study}[0.3\textwidth]{%
        \centering
        \includegraphics[width=\linewidth]{figures/_app_17_5km_fe12_evst_all_rural_fe03_1_rotated.png}
    }%
    \subcaptionbox{Honest DiD}[0.3\textwidth]{%
        \includegraphics[width=\linewidth]{figures/_app_17_5km_fe12_evst_all_rural_fe03_1_rot_honest2.png}  }%
        \\
    \begin{minipage}{0.9\textwidth}
        \footnotesize \textit{Notes:} (a) Results from Equation \ref{eq_protest}. (b) Results from Equation \ref{eq2} after de-trending for the pre-treatment period linear trend following \citet{dustmann2022reallocation}. (c) Sensitivity analysis for the average post-treatment effect allowing for violations of parallel trends \citep{rambachan2023more}.
        The reference category is $t=-1$,i.e. the month before the protest.   The regressions includes grid-cohort, and province electoral cycle fixed effects, and controls for province linear time trends, average wind speed and direction. Standard errors are clustered at the buffer area $\times$ assembly constituency level. The control group includes never treated grids only. 
    \end{minipage}
\end{figure}


%\begin{table}[htbp!]
%    \caption{ Effect of Protests occurrence on number of fires (DiD) \label{tab:table_protest}}
%    \centering 
%    \resizebox{0.8\textwidth}{!}{
%    \input{tables/_app_12_protest_5km_fe_did_rural}
%    }
%   \begin{minipage}{\textwidth}
%   {\small \textit{Notes: }Results from a simple Difference-in-Differences of the effect of protests occurrence on the number of fires. $Protest$ is a dummy equal to one if a grid-cell falls within a 5km radius buffer around a protest location, and zero otherwise. $Post$ is a dummy equal to one from the month-year when the first protest occurs within the buffer area until the end of the sample, and zero otherwise. The table reports the mean of the dependent variable in grids that never experiences a protest and in grids that experienced protest before protest occurrence. All regressions control  for average wind speed and direction.  Standard errors are clustered at the buffer area $\times$ assembly constituency level.} 
%\end{minipage}
%\end{table}


%%%%%%%% PROTEST TABLE WITH RICE%%%%%%%%%%%%
%\begin{table}[htbp!]
%    \centering
%    \caption{Effect of Protests occurrence on N Fires, by constituency agricultural characteristics }
%    \label{tab:protest_rice_table}
%    \resizebox{0.8\textwidth}{!}{
%    {\def\sym#1{\ifmmode^{#1}\else\(^{#1}\)\fi}
\begin{tabular}{l*{8}{c}}
\hline
 & \multicolumn{8}{c}{Number of Fires (in 1,000 units)} \\
 & \multicolumn{4}{c}{Without Rice Interaction} & \multicolumn{4}{c}{With Rice Interaction} \\
 & (1) & (2) & (3) & (4) & (5) & (6) & (7) & (8) \\ \hline

Post $\times$ Protest&      52.218\sym{***}&     164.771\sym{***}&     154.679\sym{***}&     104.440\sym{***}&       5.875         &      10.830\sym{*}  &      -3.261         &     -11.675\sym{*}  \\
            &    (20.133)         &    (18.923)         &    (18.421)         &    (19.550)         &     (6.689)         &     (6.016)         &     (6.359)         &     (6.559)         \\
Post $\times$ Protest $\times$ High Rice Production&                     &                     &                     &                     &      85.843\sym{***}&     181.665\sym{***}&     187.068\sym{***}&     147.307\sym{***}\\
            &                     &                     &                     &                     &    (27.785)         &    (24.990)         &    (24.750)         &    (25.772)         \\
\hline
Observations&  63,240,134         &  63,240,134         &  63,240,134         &  63,240,134         &  63,240,134         &  63,240,134         &  63,240,134         &  63,240,134         \\
N Assembly Constituencies&         808         &         808         &         808         &         808         &         808         &         808         &         808         &         808         \\
Grid $\times$ Cohort FE&           N         &           Y         &           Y         &           Y         &           N         &           Y         &           Y         &           Y         \\
Relative Time $\times$ Cohort FE&           N         &           N         &           Y         &           Y         &           N         &           N         &           Y         &           Y         \\
Province Trend $\times$ Cohort FE&           N         &           N         &           N         &           Y         &           N         &           N         &           N         &           Y         \\
Mean DV     &     473.109         &     473.109         &     473.109         &     473.109         &     473.109         &     473.109         &     473.109         &     473.109         \\
Mean DV, High Rice Production&                     &                     &                     &                     &     592.523         &     592.523         &     592.523         &     592.523         \\
\hline
\end{tabular}
}

%    }
%    \begin{minipage}{\textwidth}
%        {\small \textit{Notes: }Results from Equation \ref{eq_protest} on the number of fires with the additional interactions for constituency agricultural characteristics.  \textit{Rice Area} is a dummy equal to one for constituencies with above median rice producing areas; \textit{Harvested Rice Area} is equal to one for constituencies with above median harvest rice area; \textit{ Rice Production} is equal to one for constituencies with above median rice production. The table reports the mean of the dependent variables in control grids before the event-time (Mean DV) and in control grids before the event-time in grids with below median rice producing areas (Mean DV2).
%       All regressions control for province linear time trend, average wind speed and direction. 
%   Standard errors are clustered at the buffer area $\times$ assembly constituency level.}
%    \end{minipage}
%\end{table}
%\clearpage
\subsection{Additional Results}
%%%%%%%% PROTEST EVENTSTUDY WITH RICE%%%%%%%%%%%%
\begin{figure}[htbp!]
    \centering
    \caption{Effect of Protests occurrence on N Fires, by constituency agricultural characteristics  }
    \label{fig:es_protest_rice}
      \includegraphics[width=0.7\linewidth]{figures/_app_17_5km_fe12_evst_all_stacked_riceP_5_ori.png}\label{fig:protest_rice_production}\\
        \begin{minipage}{\textwidth}
        {\small \textit{Notes: }Results from Equation \ref{eq_protest} with the additional interactions of a dummy equal to one for constituencies with above median rice production. The reference category is $t=0$, year of the first protest. All regressions include grid-cohort and electoral cycle fixed effects and control for province linear time trends, average wind speed and direction. Standard errors are clustered at the buffer area $\times$ assembly constituency level.}
    \end{minipage}
\end{figure}

%%%%%%%% DOWNUP TIMES PROTEST %%%%%%%%%%
\begin{table}[htbp!]
    \centering
    \caption{Effect of Protest occurrence on number of fires, by constituency exposure to smoke}
    \label{tab:protest_downup}
    \resizebox{\textwidth}{!}{
    {
\def\sym#1{\ifmmode^{#1}\else\(^{#1}\)\fi}
\begin{tabular}{l*{6}{c}}
\hline
            &\multicolumn{1}{c}{(1)}         &\multicolumn{1}{c}{(2)}         &\multicolumn{1}{c}{(3)}         &\multicolumn{1}{c}{(4)}         &\multicolumn{1}{c}{(5)}         &\multicolumn{1}{c}{(6)}         \\
            & \multicolumn{6}{c}{Number of Fires (per 1,000 units)} \\ \hline
\multicolumn{7}{l}{\textbf{Panel A: Estimated Coefficients}} \\ \hline
(1) Post $\times$ Protest&     154.679\sym{***}&     154.679\sym{***}&     104.511\sym{***}&     150.521\sym{***}&     150.521\sym{***}&      99.769\sym{***}\\
            &    (18.421)         &    (18.420)         &    (19.484)         &    (20.749)         &    (20.749)         &    (21.544)         \\
(2) Down $>$ Up&                     &                     &                     &       5.044         &       5.044         &       4.472         \\
            &                     &                     &                     &     (4.490)         &     (4.490)         &     (4.562)         \\
(3) Post $\times$ Down $>$ Up&                     &                     &                     &      -0.987         &      -0.987         &      -1.631         \\
            &                     &                     &                     &     (2.562)         &     (2.562)         &     (2.545)         \\
(4) Protest $\times$ Down $>$ Up&                     &                     &                     &     -25.451         &     -25.451         &     -24.095         \\
            &                     &                     &                     &    (30.100)         &    (30.100)         &    (30.809)         \\
(5) Post $\times$ Protest $\times$ Down $>$ Up&                     &                     &                     &       8.366         &       8.366         &       9.549         \\
            &                     &                     &                     &    (19.840)         &    (19.840)         &    (19.436)         \\




\hline
\multicolumn{7}{l}{\textbf{Panel B: Linear Combinations}} \\ \hline
\multicolumn{2}{l}{Down$>$Up Effect, pre-period in no protest areas (2)}&            &            &       5.044&       5.044&       4.472\\
\multicolumn{2}{l}{}&            &            &     (4.490)&     (4.490)&     (4.562)\\
\multicolumn{2}{l}{Down$>$Up Effect, post-period in no protest areas (3)}&            &            &      -0.987&      -0.987&      -1.631\\
\multicolumn{2}{l}{}&            &            &     (2.562)&     (2.562)&     (2.545)\\
\multicolumn{2}{l}{Down$>$Up Effect, post-period in protest areas (1+3+4+5)}&            &            &132.449\sym{***}&132.449\sym{***}&83.592\sym{**}\\
\multicolumn{2}{l}{}&            &            &    (39.175)&    (39.175)&    (40.026)\\
\multicolumn{2}{l}{Differential downwind effect, post period (1+3+4+5)-(3)}&            &            &133.437\sym{***}&133.437\sym{***}&85.223\sym{**}\\
\multicolumn{2}{l}{}&            &            &    (39.692)&    (39.692)&    (40.581)\\



\hline
Observations&  63,240,134&  63,240,134&  63,240,134&  63,240,134&  63,240,134&  63,240,134\\
N Assembly Constituencies&         808&         808&         808&         808&         808&         808\\
Grid $\times$ Cohort FE&           Y&           Y&           Y&           Y&           Y&           Y\\
Relative Time $\times$ Cohort FE&           Y&           Y&           Y&           Y&           Y&           Y\\
Legislature $\times$ Cohort FE&           N&           Y&           Y&           N&           Y&           Y\\
Province Trend $\times$ Cohort FE&           N&           N&           Y&           N&           N&           Y\\
Mean DV     &     473.109&     473.109&     473.109&     473.109&     473.109&     473.109\\
Mean DV2 (Down$>$Up=1)&            &            &            &     459.315&     459.315&     459.315\\
\hline
\end{tabular}
}

    }
    \begin{minipage}{\textwidth}
        {\small \textit{Notes: } (Columns 1-3) Results from Equation \ref{eq_protest} of the effect of protest occurence on fires. (Columns 4-6) Results from Equation \ref{eq_protest} with the additional interaction for pollution visibility. $Down>Up$ is a dummy equal to one in months and grids from which most of the assembly constituency area is exposed to
smoke (see Section \ref{ssec:empirical1} for details). 
        The table reports the mean of the dependent variable in control grids before the cohort event-time.  All regressions control for average wind speed and direction. 
        Standard errors are clustered at the buffer area $\times$ assembly constituency level. \textcolor{red}{Write why we look at linear combination as triple interaction is not significant}
        }
    \end{minipage}
\end{table}






\newpage
\FloatBarrier

\section{Adverse Selection}

This section presents additional results for the adverse selection analysis. First, we present descriptive statistics for the main variables (Table \ref{tab:_desc_adverse_main}). Second, we present robustness checks using a de-trended event-study along with the Honest DiD sensitivity analysis following \citet{rambachan2023more} (Figure \ref{fig:agri_main_rotated_panel}). Lastly, we report heterogeneity by interacting treatment with constituency-level agricultural characteristics (Figure \ref{fig:es_agri_rice}) and with constituency exposure to smoke (Table \ref{tab:agri_profession_downwind}).

\subsection{Descriptives}

 \begin{table}[htbp!]
     \centering
     \caption{Descriptive Statistics \label{tab:_desc_adverse_main}}
     \resizebox{0.8\textwidth}{!}{
      \begin{tabular}{lrrrrrr}
\toprule
 & Mean & SD & Min & Max & Observations & Unique Obs.\\\\
\midrule
Grid ID &  &  &  &  & 25,625,039 & 61,577\\\\
Year &  &  &  &  & 25,625,039 & 11\\\\
Month &  &  &  &  & 25,625,039 & 12\\\\
Assembly Constituency (AC) &  &  &  &  & 25,625,039 & 689\\\\
Province &  &  &  &  & 25,625,039 & 4\\\\
Election Year &  &  &  &  & 25,625,039 & 9\\\\
Cohort &  &  &  &  & 25,625,039 & 7\\\\
Legislature &  &  &  &  & 25,625,039 & 12\\\\
Relative year & -1.595 &  &  &  & 25,625,039 & 21\\\\
Agricultural Politician &  &  &  &  & 25,625,039 & 2\\\\
Number of Fires (in 1,000 units) & 172.666 &  &  &  & 25,625,039 & 45\\\\
High Rice Production (AC level) & 0.466 &  &  &  & 25,625,039 & 2\\\\
\bottomrule
\end{tabular}

      }
      \begin{minipage}{\textwidth}
         {\small \textit{Notes: } This table shows descriptive statistics of the main variable for the adverse selection results }
     \end{minipage}
 \end{table}
 
%%%%%%%%%% TREATMENT VARIATION AGRICULTURE %%%%%%%%%
%\begin{figure}[htbp!]
%    \centering
%    \caption{Politicians' characteristics: treatment variation}\label{fig:agri_switchers_map}
%    \includegraphics[width = 0.7\textwidth]{figures/8_polisc_switchers.png}
%    \begin{minipage}{\textwidth}
%        {\small \textit{Notes: }The map shows the spatial distribution of treatment status, i.e. switching from electing a non-agriculturalist to electing an agriculturalist politician. In blue assembly constituencies that never switch; in red assembly constituencies that always elect and agriculturalists; in yellow treated units. 
%        \red{Should we change the label of the legend for non-switcher treated? the word treated may be confusing. Different map for each election, stacked dataset switcher and not switchers or Table at the cohort level where for each we say switchers, non switchers} }
%    \end{minipage}
%\end{figure}

%%%%%%%%%% Profession  Summary Table %%%%%%%%%
%   \begin{table}[htbp!]
%       \centering
%       \caption{Politicians' Characteristics}
%       \label{tab:profession_desc} 
%    \resizebox{0.8\textwidth}{!}{
%    \input{tables/_app_desc_polisch_char}
%    }
%\end{table}

\FloatBarrier
\subsection{Robustness}
    \begin{figure}[htbp!]
    \caption{Effect by relative time period from switches to an agriculturalist politician}
    \label{fig:agri_main_rotated_panel}
    \centering

    \subcaptionbox{Event-Study}[0.3\textwidth]{%
        \includegraphics[width=\linewidth]{figures/_app_16_polischar_fe03_evst_main_rural_acpop_1_ori.png} 
         }%
        \subcaptionbox{De-Trended Event-Study}[0.3\textwidth]{%
        \centering
        \includegraphics[width=\linewidth]{figures/_app_16_polischar_fe03_evst_main_rural_acpop_1_rotated.png}
    }%
    \subcaptionbox{Honest DiD}[0.3\textwidth]{%
        \includegraphics[width=\linewidth]{figures/_app_16_polischar_fe03_evst_main_rural_acpop_1_rot_honest2.png}  }%
        \\
       

    \begin{minipage}{0.9\textwidth}
        \footnotesize \textit{Notes:} (a) Results from Equation \ref{eq_agriculture}. (b) Results from Equation \ref{eq_agriculture} after de-trending for the pre-treatment period linear trend following \citet{dustmann2022reallocation}. (c) Sensitivity analysis for the average post-treatment effect allowing for violations of parallel trends \citep{rambachan2023more}.
        The reference category is $t=-1$,i.e. the month before the switch to an agriculturalist.   The regressions includes grid-cohort, and province electoral cycle fixed effects, and controls for province linear time trends, average wind speed and direction. Standard errors are clustered at the buffer area $\times$ assembly constituency level. The control group includes never treated grids only. 
    \end{minipage}
\end{figure}

\FloatBarrier
%%%%%%%%%% Agriculturalist Politicians with Rice Table %%%%%%%%%
\subsection{Additional Results}
%%%%%%%%%% Agriculturalist Politicians with Rice ES %%%%%%%%%
%\begin{table}[htbp!]
%    \centering
%    
%    \caption{Effect of electing an agriculturalist politician,by constituency agricultural characteristics}
%    \label{tab:agri_profession_rice_table}
%    \resizebox{\textwidth}{!}{
%    \centering
%    {\def\sym#1{\ifmmode^{#1}\else\(^{#1}\)\fi}
\begin{tabular}{l*{8}{c}}
\hline
 & \multicolumn{8}{c}{Number of Fires (in 1,000 units)} \\
 & \multicolumn{4}{c}{Without Rice Interaction} & \multicolumn{4}{c}{With Rice Interaction} \\
 & (1) & (2) & (3) & (4) & (5) & (6) & (7) & (8) \\ \hline

Post $\times$ Agriculturalist&      79.200\sym{**} &      29.898\sym{***}&      12.465\sym{***}&      12.483\sym{***}&      -1.249         &       0.023         &       1.418         &       1.462         \\
            &    (32.104)         &     (5.608)         &     (4.525)         &     (4.526)         &    (14.617)         &     (4.026)         &     (4.184)         &     (4.188)         \\
Post $\times$ Agriculturalist $\times$ High Rice Production&                     &                     &                     &                     &     155.962\sym{**} &      61.520\sym{***}&      26.071\sym{***}&      26.009\sym{***}\\
            &                     &                     &                     &                     &    (63.142)         &    (11.313)         &     (9.572)         &     (9.576)         \\
\hline
Observations&   8,268,839         &   8,268,839         &   8,268,839         &   8,268,839         &   8,268,839         &   8,268,839         &   8,268,839         &   8,268,839         \\
N Assembly Constituencies&         695         &         695         &         695         &         695         &         695         &         695         &         695         &         695         \\
Grid $\times$ Cohort FE&           N         &           Y         &           Y         &           Y         &           N         &           Y         &           Y         &           Y         \\
Relative Time $\times$ Cohort FE&           N         &           N         &           Y         &           Y         &           N         &           N         &           Y         &           Y         \\
Province Trend $\times$ Cohort FE&           N         &           N         &           N         &           Y         &           N         &           N         &           N         &           Y         \\
Mean DV     &     105.231         &     105.231         &     105.231         &     105.231         &     105.231         &     105.231         &     105.231         &     105.231         \\
Mean DV, High Rice Production&                     &                     &                     &                     &     194.691         &     194.691         &     194.691         &     194.691         \\
\hline
\end{tabular}
}

%    }
%    \begin{minipage}{\textwidth}
%        {\small \textit{Notes: } Results from Equation \ref{eq_agriculture} with an additional interaction for the constituency agricultural characteristics.  The dummy \textit{Above Median Rice Area} takes value 1 for constituencies with above median rice area. The dummy \textit{Above Median Harvested Rice Area} takes value 1 for constituencies with above median harvested rice area. The dummy \textit{Above Median Rice Production} takes value 1 for constituencies with above median rice production. 
%        The table reports the mean of the dependent variables in control grids before the event-time (Mean DV) and in control grids before the event-time in grids with below median rice producing areas (Mean DV2).
%        All regressions include grid-cohort, month-year and government term-year fixed effects and control for province linear time trends, average wind speed and direction. Standard errors are clustered at the grid and assembly constituency $\times$ month-year level. }
%    \end{minipage}
%    
%\end{table}

\begin{figure}[htbp!]
        \centering
        \caption{Effect of switching to an agriculturalist politician, by constituency agricultural characteristics}\label{fig:es_agri_rice}
      
       \includegraphics[width=0.48\linewidth]{figures/_app_16_polischar_fe12_evst_all_rural_acpop_riceP_5_rotated.png}\label{fig_es_agri_production}\\
       \begin{minipage}{\textwidth}
         {\small \textit{Notes: }Results from Equation \ref{eq_agriculture} with an additional interaction for the constituency agricultural characteristics. The dummy \textit{Above Median Rice Production} takes value 1 for constituencies with above median rice production. 
         The regression includes grid-cohort, and legislature fixed effects, and controls for province linear time trends, average wind speed and direction. Standard errors are clustered at the grid and assembly constituency $\times$ month-year level. }
         \end{minipage}
    \end{figure}


%%%%%%%%%%%%%%%%% Protest Effect on Elections %%%%%%%%%%%%%%%%%%%

%\begin{table}[htbp!]
%    \centering
%    \caption{Effect of Protests Occurrence on Electoral Outcomes}
%    \label{tab:protest_election_table}
%    \resizebox{0.8\textwidth}{!}{
%    \input{tables/protest_on_electoral_outcomes}
%    }
%    \begin{minipage}{\textwidth}
%        {\small \textit{Notes: }Results from Equation 3 examining the effect of protests on electoral outcomes. The treatment dummy $Protest$ is equal to one for assembly constituencies that experienced a protest from the first protest until the end of the sample.  The dependent variable is turnover in columns (1)-(4) and margin of victory in percentage in columns (5)-(8). All regressions include assembly constituency and relative time from protest fixed effects. Standard errors are clustered at the assembly constituency level. }
%    \end{minipage}
%\end{table}

%%%%%%%%%% Agriculturalist Politicians with downwind %%%%%%%%%

\begin{table}[htbp!]
    \centering
    \caption{Effect of switching to an agriculturalist politician on number of fires, by constituency exposure to smoke}
    \label{tab:agri_profession_downwind}
    \resizebox{\textwidth}{!}{
      {
\def\sym#1{\ifmmode^{#1}\else\(^{#1}\)\fi}
\begin{tabular}{l*{6}{c}}
\hline
            &\multicolumn{1}{c}{(1)}         &\multicolumn{1}{c}{(2)}         &\multicolumn{1}{c}{(3)}         &\multicolumn{1}{c}{(4)}         &\multicolumn{1}{c}{(5)}         &\multicolumn{1}{c}{(6)}         \\
            & \multicolumn{6}{c}{Number of Fires (in 1,000 units)} \\ \hline
\multicolumn{7}{l}{\textbf{Panel A: Estimated Coefficients}} \\ \hline
(1) Post $\times$ Switch Agric&      26.660\sym{***}&      11.300\sym{**} &      11.212\sym{**} &      25.908\sym{***}&       9.939\sym{*}  &       9.918         \\
            &     (4.519)         &     (4.843)         &     (4.851)         &     (5.725)         &     (6.021)         &     (6.033)         \\
(2) Down $>$ Up&                     &                     &                     &       4.783\sym{**} &       5.161\sym{***}&       5.073\sym{***}\\
            &                     &                     &                     &     (1.956)         &     (1.945)         &     (1.948)         \\
(3) Post $\times$ Down $>$ Up&                     &                     &                     &      -1.487         &      -2.760         &      -2.766         \\
            &                     &                     &                     &     (2.374)         &     (2.368)         &     (2.373)         \\
(4) Switch Agric  $\times$ Down $>$ Up&                     &                     &                     &      -1.532         &      -1.960         &      -1.315         \\
            &                     &                     &                     &     (4.506)         &     (4.484)         &     (4.487)         \\
(5) Post $\times$ Switch Agric  $\times$ Down $>$ Up&                     &                     &                     &       1.461         &       2.648         &       2.514         \\
            &                     &                     &                     &     (6.434)         &     (6.417)         &     (6.427)         \\




\hline
\multicolumn{7}{l}{\textbf{Panel B: Linear Combinations}} \\ \hline
\multicolumn{2}{l}{(2) Down$>$Up Effect, pre-period of no-switch politicians}&            &            &4.783\sym{**}&5.161\sym{***}&5.073\sym{***}\\
            &            &            &     (1.956)&     (1.945)&     (1.948)\\
\multicolumn{2}{l}{(2+3) Down$>$Up Effect, post-period of no-switch politicians}&            &            &3.296\sym{*}&       2.400&       2.307\\
            &            &            &     (1.937)&     (1.944)&     (1.949)\\
\multicolumn{2}{l}{(1+2+3+4+5) Down$>$Up Effect, post-period of switch politicians}&            &            &29.132\sym{***}&13.028\sym{**}&13.425\sym{**}\\
            &            &            &     (5.666)&     (5.852)&     (5.866)\\
\multicolumn{2}{l}{(1+2+3+4+5)-(2+3) Differential Down$>$Up effect, post period}&            &            &25.837\sym{***}&10.627\sym{*}&11.117\sym{*}\\
            &            &            &     (5.772)&     (5.952)&     (5.966)\\



\hline
Observations&  25,625,039&  25,625,039&  25,625,039&  25,625,039&  25,625,039&  25,625,039\\
N Assembly Constituencies&         689&         689&         689&         689&         689&         689\\
Grid FE     &           Y&           Y&           Y&           Y&           Y&           Y\\
Relative Time FE&           Y&           Y&           Y&           Y&           Y&           Y\\
Legislature FE&           N&           Y&           Y&           N&           Y&           Y\\
Province Trend FE&           N&           N&           Y&           N&           N&           Y\\
Mean DV     &     120.687&     120.687&     120.687&     120.687&     120.687&     120.687\\
Mean DV2 (Down$>$Up=1)&            &            &            &     121.718&     121.718&     121.718\\
\hline
\end{tabular}
}

    }
     \begin{minipage}{\textwidth}
        {\small \textit{Notes: } (Columns 1-3) Results from Equation \ref{eq_agriculture} of the effect of switching to electing a politician with agricultural ties on fires. (Columns
4-6)   Results from Equation \ref{eq_agriculture}  with the additional interaction for constituency exposure to smoke. $Down>Up$ is
a dummy equal to one in months and grids from which most of the assembly constituency area is
exposed to smoke (see Section \ref{ssec:empirical1} for details). The table reports the mean of the dependent variable
in control grids before the cohort event-time.  All regressions control for average wind speed and direction.  Standard errors are clustered at the assembly constituency $\times$ election year level.}
    \end{minipage}
\end{table}


% \begin{figure}[htbp!]
%     \caption{Effect by relative time period from politician switch}
%     \label{fig:es_main_rotated_panel}
%     \centering
%     % --- 4 columns x 4 rows ---
%     \begin{tabular}{cccc}
%         \multicolumn{4}{c}{Effect of switching to a agriculturalist politician on fires} \\
        
%         \includegraphics[width=0.22\textwidth]{figures/8_polisc_selspec17_0_ori.png} &
%         \includegraphics[width=0.22\textwidth]{figures/8_polisc_selspec17_0_honest1.png} &
%         \includegraphics[width=0.22\textwidth]{figures/8_polisc_selspec17_0_rotated.png} &
%         \includegraphics[width=0.22\textwidth]{figures/8_polisc_selspec17_0_honest2.png} \\
%         \small (a) Original & \small (b) Honest (Non-rotated) & \small (c) Rotated & \small (d) Honest (Rotated) \\ 

%                 \multicolumn{4}{c}{High Rice Assembly Constituencies} \\
        
%         \includegraphics[width=0.22\textwidth]{figures/8_polisc_selspec17_1_ori.png} &
%         \includegraphics[width=0.22\textwidth]{figures/8_polisc_selspec17_1_honest1.png} &
%         \includegraphics[width=0.22\textwidth]{figures/8_polisc_selspec17_1_rotated.png} &
%         \includegraphics[width=0.22\textwidth]{figures/8_polisc_selspec17_1_honest2.png} \\
%         \small (a) Original & \small (b) Honest (Non-rotated) & \small (c) Rotated & \small (d) Honest (Rotated) \\ 

%         \multicolumn{4}{c}{High Rice Harvested Assembly Constituencies} \\
        
%         \includegraphics[width=0.22\textwidth]{figures/8_polisc_selspec17_2_ori.png} &
%         \includegraphics[width=0.22\textwidth]{figures/8_polisc_selspec17_2_honest1.png} &
%         \includegraphics[width=0.22\textwidth]{figures/8_polisc_selspec17_2_rotated.png} &
%         \includegraphics[width=0.22\textwidth]{figures/8_polisc_selspec17_2_honest2.png} \\
%         \small (a) Original & \small (b) Honest (Non-rotated) & \small (c) Rotated & \small (d) Honest (Rotated) \\ 

%         \multicolumn{4}{c}{High Rice Production Assembly Constituencies} \\
        
%         \includegraphics[width=0.22\textwidth]{figures/8_polisc_selspec17_3_ori.png} &
%         \includegraphics[width=0.22\textwidth]{figures/8_polisc_selspec17_3_honest1.png} &
%         \includegraphics[width=0.22\textwidth]{figures/8_polisc_selspec17_3_rotated.png} &
%         \includegraphics[width=0.22\textwidth]{figures/8_polisc_selspec17_3_honest2.png} \\
%         \small (a) Original & \small (b) Honest (Non-rotated) & \small (c) Rotated & \small (d) Honest (Rotated) \\ 

        
%         \end{tabular}

% \end{figure}


% \begin{figure}[htbp!]
%     \caption{Effect by relative time period from protest time}
%     \label{fig:es_main_rotated_panel_protest}
%     \centering
%     % --- 4 columns x 4 rows ---
%     \begin{tabular}{cccc}
%         \multicolumn{4}{c}{Effect of Protest on Number of fires} \\
        
%         \includegraphics[width=0.22\textwidth]{figures/protest_5km_spec17_1_ori.png} &
%         \includegraphics[width=0.22\textwidth]{figures/protest_5km_spec17_1_honest1.png} &
%         \includegraphics[width=0.22\textwidth]{figures/protest_5km_spec17_1_rotated.png} &
%         \includegraphics[width=0.22\textwidth]{figures/protest_5km_spec17_1_honest2.png} \\
%         \small (a) Original & \small (b) Honest (Non-rotated) & \small (c) Rotated & \small (d) Honest (Rotated) \\ 

%                 \multicolumn{4}{c}{High Rice Assembly Constituencies} \\
        
%         \includegraphics[width=0.22\textwidth]{figures/protest_5km_spec17_2_ori.png} &
%         \includegraphics[width=0.22\textwidth]{figures/protest_5km_spec17_2_honest1.png} &
%         \includegraphics[width=0.22\textwidth]{figures/protest_5km_spec17_2_rotated.png} &
%         \includegraphics[width=0.22\textwidth]{figures/protest_5km_spec17_2_honest2.png} \\
%         \small (a) Original & \small (b) Honest (Non-rotated) & \small (c) Rotated & \small (d) Honest (Rotated) \\ 

%         \multicolumn{4}{c}{High Rice Harvested Assembly Constituencies} \\
        
%         \includegraphics[width=0.22\textwidth]{figures/protest_5km_spec17_3_ori.png} &
%         \includegraphics[width=0.22\textwidth]{figures/protest_5km_spec17_3_honest1.png} &
%         \includegraphics[width=0.22\textwidth]{figures/protest_5km_spec17_3_rotated.png} &
%         \includegraphics[width=0.22\textwidth]{figures/protest_5km_spec17_3_honest2.png} \\
%         \small (a) Original & \small (b) Honest (Non-rotated) & \small (c) Rotated & \small (d) Honest (Rotated) \\ 

%         \multicolumn{4}{c}{High Rice Production Assembly Constituencies} \\
        
%         \includegraphics[width=0.22\textwidth]{figures/protest_5km_spec17_4_ori.png} &
%         \includegraphics[width=0.22\textwidth]{figures/protest_5km_spec17_4_honest1.png} &
%         \includegraphics[width=0.22\textwidth]{figures/protest_5km_spec17_4_rotated.png} &
%         \includegraphics[width=0.22\textwidth]{figures/protest_5km_spec17_4_honest2.png} \\
%         \small (a) Original & \small (b) Honest (Non-rotated) & \small (c) Rotated & \small (d) Honest (Rotated) \\ 

        
%         \end{tabular}

% \end{figure}



% \begin{figure}[htbp!]
%     \caption{Effect by relative time period from protest time during fire season}
%     \label{fig:es_main_rotated_panel_protest}
%     \centering
%     % --- 4 columns x 4 rows ---
%     \begin{tabular}{cccc}
%         \multicolumn{4}{c}{Effect of Protest on Number of fires} \\
        
%         \includegraphics[width=0.22\textwidth]{figures/protest_5km_spec17_hm_int_1_ori.png} &
%         \includegraphics[width=0.22\textwidth]{figures/protest_5km_spec17_hm_int_1_honest2.png} &
%         \includegraphics[width=0.22\textwidth]{figures/protest_5km_spec17_hm_int_1_rotated.png} &
%         \includegraphics[width=0.22\textwidth]{figures/protest_5km_spec17_hm_int_1_rot_honest2.png} \\
%         \small (a) Original & \small (b) Honest (Non-rotated) & \small (c) Rotated & \small (d) Honest (Rotated) \\ 

%                 \multicolumn{4}{c}{High Rice Assembly Constituencies} \\
        
%         \includegraphics[width=0.22\textwidth]{figures/protest_5km_spec17_hm_int_2_ori.png} &
%         \includegraphics[width=0.22\textwidth]{figures/protest_5km_spec17_hm_int_2_honest2.png} &
%         \includegraphics[width=0.22\textwidth]{figures/protest_5km_spec17_hm_int_2_rotated.png} &
%         \includegraphics[width=0.22\textwidth]{figures/protest_5km_spec17_hm_int_2_rot_honest2.png} \\
%         \small (a) Original & \small (b) Honest (Non-rotated) & \small (c) Rotated & \small (d) Honest (Rotated) \\ 

%         \multicolumn{4}{c}{High Rice Harvested Assembly Constituencies} \\
        
%         \includegraphics[width=0.22\textwidth]{figures/protest_5km_spec17_hm_int_3_ori.png} &
%         \includegraphics[width=0.22\textwidth]{figures/protest_5km_spec17_hm_int_3_honest2.png} &
%         \includegraphics[width=0.22\textwidth]{figures/protest_5km_spec17_hm_int_3_rotated.png} &
%         \includegraphics[width=0.22\textwidth]{figures/protest_5km_spec17_hm_int_3_rot_honest2.png} \\
%         \small (a) Original & \small (b) Honest (Non-rotated) & \small (c) Rotated & \small (d) Honest (Rotated) \\ 

%         \multicolumn{4}{c}{High Rice Production Assembly Constituencies} \\
        
%         \includegraphics[width=0.22\textwidth]{figures/protest_5km_spec17_hm_int_4_ori.png} &
%         \includegraphics[width=0.22\textwidth]{figures/protest_5km_spec17_hm_int_4_honest2.png} &
%         \includegraphics[width=0.22\textwidth]{figures/protest_5km_spec17_hm_int_4_rotated.png} &
%         \includegraphics[width=0.22\textwidth]{figures/protest_5km_spec17_hm_int_4_rot_honest2.png} \\
%         \small (a) Original & \small (b) Honest (Non-rotated) & \small (c) Rotated & \small (d) Honest (Rotated) \\ 

        
%         \end{tabular}

% \end{figure}




 
 
  

\begin{comment} 
\section{Politicians Quotes}
\subsection{Assembly Documents - Stubble burning}


Stubble Burning

------------------------------
PUNJAB -
------------------------------
1. PUNJAB VIDHAN SABHA
 DEBATES 
Wednesday, the 29th November, 2017 
Vol. III -No. 3 
OFFICIAL REPORT
(460.pdf)

During 2017 in the Agriculture Minister november debate, Sardar Sukhpal Singh Khaira Punjab MLA strongly criticized the government for its low efforts on stubble burning managment:
"The government spent ₹97 crore for stubble management, but that money was diverted to other issues. This shows how serious the government is about solving the problem. They talk about allocating ₹5000 per acre for stubble management, but nothing concrete is happening. I hope that we will be able to manage the stubble by the next season."

In this debate the minister remarked that stubble burning only account on a small proportion of the energy grid of Punjab:
"Currently, projects are underway to generate 62.7 megawatts of electricity from biomass, which is being converted from straw. NRG is also involved. In Delhi, two locations are running stubble burning projects alongside NRG. The total stubble burning project in Delhi accounts for only 2%, while 74% of samples are processed. NRG’s process takes 4–5 days. In the meeting, the situation of stubble burning in Delhi was discussed. NRG clearly explained the benefits and drawbacks, stating that we are in touch with the NGT and the government is also concerned about finding a solution..."

---------------------------------------------------------------------
--------------------------------------------------------------------- 
2. PUNJAB VIDHAN SABHA
 DEBATES
 Friday, the 17th January, 2020
 Vol. X- (Special Session)-No. 1
 OFFICIAL REPORT
(1213.pdf)

Debate of Agricultural Prime Minister
Vol. X- (Special Session)-No. 
LAYING OF GOVERNOR’S ADDRESS ON THE TABLE OF THE HOUSE 

"32. My Government recognizes the urgency in resolving the issue
 linked with stubble burning. Under the In-Situ Crop Residue
 Management scheme, during the last financial year, 28609
 residue management machines have been distributed through
 5439 custom hiring centers at 80% subsidy and 12075 to
 individual farmers at 50% subsidy. In addition, during the year
 2019-20 till date approximately 24000 machines has been
 provided to the farmers at subsidized rates."

"72. The State is constantly looking at using more & more paddy
 stubble for economic use to tackle the problem of Paddy
 Stubble burning. Presently, 10 Biomass power projects of  7350
 MWs are in operation and two paddy straw based biomass
 power projects of 36 MWs have recently been commissioned
 to utilize the paddy stubble, These plants are utilizing
 approximatetly 1.10 million tons of paddy straw per annum."
 
"73. Private investment is being facilitated for setting up Bio-CNG/
 Bio-ethanol projects based on paddy stubble.My Government
 has submitted a proposal to the Ministry of New and Renewable
 Energy, Government of India for grant of viability gap funding
 for paddy-straw based Biomass power projects @ 5 Cr. per
 MW for 300 MW capacity projects. "

---------------------------------------------------------------------
---------------------------------------------------------------------
 3. DEMANDS FOR GRANTS(1-14)
 GOVERNMENT OF PUNJAB
 (With Detailed Estimates of Expenditure)
 2021-2022
(2467.pdf)

During the period 2021-2022 Stubble managment has a specific budget
 "Total 02 Compensation for Stubble Management 30,00,00"

---------------------------------------------------------------------
---------------------------------------------------------------------
4. Speech of
Harpal Singh Cheema
Finance Minister, Punjab
Presenting the Budget for the Year 2022-23 to the Punjab Vidhan Sabha
(2367.pdf)

Speech of the finance minister of Punjab (2022-2023) - Harpal Singh Cheema

Harpal Singh Cheema first budget of the Aam Aadmi Party (AAP) government under Chief Minister Bhagwant Singh Mann presented direct measures to promote sowing of paddy with it own budged.

Its noticibly adress the issue in the point 66:
"66. Checking Stubble Burning Practices: We realize that burning 
of paddy straw by farmers is not without a reason as they need to get 
their fields ready for the next crop at the right time to avoid any yield 
loss. This practice causes air pollution, negatively impacting health of 
people across states. Further studies have shown that stubble burning 
is not good for soil health as well. Therefore, this Government would 
only undertake positive targeted actions for controlling it. I propose to 
explore various possibilities and solutions for stubble burning. An 
allocation of ₹200 crore is hereby proposed in this budget. "

With respect to their position on farmers it is a clear support.
"67. Free Power to Farmers:   I want to put all speculations to rest. 
AAP Government pledges to stand by its farmers and shall continue 
providing free power for the agriculture sector. I am proposing a
budgetary allocation of ₹6,947 crore in this FY 2022-23. "


---------------------------------------------------------------------
---------------------------------------------------------------------
5. ADDRESS OF THE GOVERNOR OF PUNJAB SHRI BANWARILAL PUROHIT TO THE STATE LEGISLATURE 
Friday 01 March, 20242
(2519.pdf)

In 2024 the governor of Punjab adress the issue in one of its speech and said

"My Government has been providing 80% subsidy to farmer groups/panchayats and 50% subsidy to individual farmers for purchase of machineries, for both in-situ and ex-situ management of crop residue in the state. From the financial year (FY) 2018-19 to FY 2022-23 a total of 1,17,000 machines has been provided for stubble management. During the financial year 2023-24, approximately 23,000 machines have been provided to farmers under this scheme and similarly in the ensuing Financial Year also this scheme shall be continued. My Government is also aggressively pursuing the development of Compressed Biogas (CBG) projects which shall utilize surplus paddy straw to the tune of 1.87 million tons per annum. Such measures have reduced the fire incidents during the last Kharif harvesting season from 49,922 in the previous season, to 36,623 this season which turns out to be 26%. Being sensitive to the cause of farmers..."

---------------------------------------------------------------------
---------------------------------------------------------------------
------------------------------
Haryana -
------------------------------
---------------------------------------------------------------------
---------------------------------------------------------------------
6. HARYANA VIDHAN SABHA 
RESUME OF BUSINESS TRANSACTED BY THE THIRTEENTH HARYANA VIDHAN SABHA DURING THE TENTH (OCTOBER) SESSION, 2017 WITH APPENDIX CONTAINING DAILY BULLETINS

The Agricultural Minister was noticed about the problems of crop residue burning by 5 MLAs:
"The Speaker informed the House that he had received a Calling Attention Notice No.15 from Sardar Jaswinder Singh Sandhu and 4 other M.L.As. (S/Shri Abhay Singh Chautala, Ram Chand Kamboj, Balwan Singh Daulatpuria and Ved Narang regarding the burning problem of crop residue in the fields and the same has been admitted. "

After what the Minister declared a statement (not avaible on the document)



---------------------------------------------------------------------

7. HARYANA VIDHAN SABHA 
XII  SESSION  
BULLETIN NO. 3 
Tuesday, the 11th September, 2018 (1st Sitting) 
(Bulletin-No.-3-119181st-Sitting.pdf)

The finance minister was asked about several topics, which included parali burning, indicating that the goverment has taken action in this topic
"(iii) Shri Abhay Singh Chautala further stated that the Pollution 
Control Board has imposed penalty @  Rs. 3,000/- per acre on 
those farmers, who have burnt  the husk (Parali) in their fields 
in Jind District. He further alleged that in this case an S.D.O. 
has collected illegally Rs. 3-4 crores by issuing the duplicate 
receipts but the Government has not taken any action against 
him so far.  "

---------------------------------------------------------------------
8. HARYANA VIDHAN SABHA
 Un-Starred  Questions
 to be asked at a meeting of the Haryana Vidhan Sabha to be held on
 Tuesday, the 3rd March, 2020
(Un-Starred-Questions-3rd-March-English.pdf)

 "Shri  Lakshman Napa, M.L.A. : Will the Agriculture and Farmers Welfare
 Minister be pleased to state —
.... whether there is any proposal under consideration of  the Government
 to set up a paddy straw factory at the places of the above said factories
 togetherwith the details thereof  ? "


---------------------------------------------------------------------
9. HARYANA VIDHAN SABHA 
REPORT OF THE COMMITTEE ON THE 
WELFARE OF SCHEDULED CASTES, 
SCHEDULED TRIBES AND BACKWARD CLASSES 
(2021-2022) 
(FOURTEENTH VIDHAN SABHA) 
FORTY FIFTH REPORT
(45th-Report-2021-22.pdf)


In 2021 a subsidy to alternative ways to solve the problem of crop residue:

"It has been also proposed to provide subsidy under state 
scheme (SB-89) to straw management projects like 
briquette/pallet/fodder making @ 50% subsidy (it may be 
capped at 50 lakhs) having maximum project cost of Rs. 
100 lakhs, for establishment of CHC/Crop Residue 
Management Projects. The group of farmer/FPOs/ 
Registered Societies should have at least 40% of 
members from SC category. This can be a credit linked 
project wherein 25% subsidy may be released upon 
verification/purchase of equipment and rest 75% after the 
locking period of 3 years is over. Till then subsidy 
amount will remain with bank"
---------------------------------------------------------------------

10. Proceedings of the Haryana Legislative Assembly 22 December, 2021 Volume–3, Issue–4 Official Record (Emblem of Haryana) Table of Contents Wednesday, 22 December, 2021
(22.12.2021-compressed.pdf)

Private efforts of refinery to create a Ethanol Plant in order to 
burn stubble in a more sustainable way

The minister was asked about the number of cases registered against
farmers for stubble burning with the fines imposed in this cases.
We have the table



--
 \end{comment}

 
 
 
\begin{comment}
 
\subsection{Assembly Documents - Farmers}
Farmer Protest
--------------------------------------
Punjab
--------------------------------------
**
15th Vidhan Sabha 
15th Special) Session 
Revised 
PUNJAB VIDHAN SABHA 
BULLETIN No. 1 
Friday, the 3 September, 2021 

During a set of tributes to recently deceased personalitites several honorables members of the Vidhan Session asked to pay tribute to farmers who lost their lives during the farmer protest in today's list of condolence resolutions. the Speaker (Shri Kuldeep Singh Dhoss) included their names in the list of obituary references. and inform that the condolence messages would be sent to the bereaved families. 
---------------------

* -- * -- * -- * -- * -- * -- * -- * -- * -- * -- * -- * -- * -- *
Haryana 
------------------------------------------------------------------
HARYANA VIDHAN SABHA 
RESUME OF BUSINESS TRANSACTED BY THE THIRTEENTH HARYANA VIDHAN SABHA DURING THE THIRTEENTH (DECEMBER) SESSION, 2018 WITH APPENDIX CONTAINING DAILY BULLETINS HARYANA VIDHAN SABHA SECRETARIAT CHANDIGARH 
DECEMBER, 2018
(December-2018.pdf)


After the Legislative Business, all the Members of the Indian National Lok Dal, present in the House, staged a walk out as protest against not to admit their Adjournment Motions regarding the S.Y.L. Canal issue and to waive off the loan of the farmers and small traders. 
------------------------------------------------------------------
HARYANA VIDHAN SABHA RESUME OF BUSINESS TRANSACTED BY THE FOURTEENTH HARYANA VIDHAN SABHA DURING THE SEVENTH (DECEMBER) SESSION, 2021 WITH APPENDIX CONTAINING DAILY BULLETINS HARYANA VIDHAN SABHA SECRETARIAT CHANDIGARH 
DECEMBER, 2021
(December-2021.pdf)

On being disallowed the Non-Official Resolution regarding Minimum Support Price of crops, all the Members of the Indian National Congress Party, present in the House, staged a walk out as a protest.  
------------------------------------------------------------------

\end{comment}

\begin{comment}
\subsection{Tweets from official accounts}

----------------------------------------------------------------------------
satyakumar_y:
Major steps taken by 
@ArvindKejriwalji to curb pollution in Delhi:

Before Diwali:

- Shut down schools and colleges.
- Ban firecrackers during Diwali.
- Blame Haryana and UP for stubble burning.
- Blame the central government for not allowing him to work.
- Blame

----------------------------------------------------------------------------
satyakumar_y:
Everyone knows that burning fire crackers in Diwali doesn't cause much pollution when compared to massive stubble burning in Punjab. Yet kejriwal chooses to ban fire crackers in Delhi instead of making arrangements for farmers in Punjab to prevent them from stubble burning!
----------------------------------------------------------------------------
satyakumar_y:
I consider the Kisan Andolan to be Pavitra

But,when Andolanjeevis hijack Pavitra Andolans,showcase photos of those jailed for serious offences,does it serve any purpose? 

Not allowing toll plazas to work, destroying telecom towers-does it serve a Pavitra Andolan

#PMInLokSabha
----------------------------------------------------------------------------

\end{comment}
\FloatBarrier
\section{Political Incentives Mobilize Government Machinery}
\begin{table}[htbp]
    \centering
    \caption{Pollution Exposure and Politicians' Actions, all categories}
    \label{tab:court_disha_app}
    \resizebox{\textwidth}{!}{
    \input{tables/action_final_additional_disha}
    }
    \begin{minipage}{\textwidth}
        {\small \textit{Notes:}  The table reports politicians' actions by level of constituency exposure. 
        \textit{Exposed Population Share}  is the total population that upwind grids in a constituency-month would expose, normalized by the constituency's total population, and standardized (z-scored) within the estimation sample.
        We use DISHA documents to report the probability that a document is recorded per 100 in a given month-year, for the following categories: development (column 1), elections (column 2), governance (column 3), and welfare (column 4).
        The table reports the mean of the dependent variable in the sample. All regressions control for wind direction and average wind speed. Standard errors are clustered at the constituency level.  }
    \end{minipage}
\end{table}

 
 
 \end{document}
%%%%%%%%%%%%%%%%% Protest Effect on Elections %%%%%%%%%%%%%%%%%%%

\begin{figure}
    \centering
    \caption{Effect of Protests Occurrence on Electoral Outcomes}
    \label{fig:protest_election_es}
    \centering
     \centering\subfloat[Turnover]{\includegraphics[width=0.45\linewidth]{figures/pr_on_polchar_spec1_rotated.png}\label{fig:protest_turnover}}
       \centering\subfloat[Vote Margin ]{\includegraphics[width=0.45\linewidth]{figures/pr_on_polchar_spec19_rotated.png}\label{fig:protest_vote_margin}}\\
    \begin{minipage}{\textwidth}
        {\small \textit{Notes: }Results from Equation \ref{eq_protest} examining the effect of protests on electoral outcomes by elections before and after the protest event. The left panel shows the effect of protest occurrence on voter turnout. The right panel shows the effect on the margin of votes in percentage points. The analysis is conducted at the assembly constituency level and includes only constituencies that experienced at least one protest. All regressions include assembly constituency and relative time from protest fixed effects. Standard errors are clustered at the assembly constituency level.}
    \end{minipage}
\end{figure}

\section{@Anzony: To Control keep map and drop others}
% Different Specification for Protests

\begin{figure}[htbp!]
\centering
\begin{minipage}[b]{0.45\textwidth}
  \centering
  \includegraphics[width=\linewidth]{figures/2km_plot.png}
  \caption*{2Km}
\end{minipage}
\hfill
\begin{minipage}[b]{0.45\textwidth}
  \centering
  \includegraphics[width=\linewidth]{figures/10km_plot.png}
  \caption*{10Km}
\end{minipage}
\vspace{1em}
\begin{minipage}{0.9\textwidth}
  \footnotesize
  \textbf{Notes:} This figure shows treated and control grids for different ratios around protest treatment:2km and 10km.
\end{minipage}
\end{figure}

\begin{table}[htbp!]
   \caption{\label{tab:app_protest_2km} Protest  around 2km}
   \singlespacing \centering \resizebox{\textwidth}{!}{
\input{tables/protest_2km}}
\begin{minipage}{0.9\textwidth}
  \footnotesize
  \textbf{Notes:} The table reports how the number of fires changes when there is any farmer around 5km in the Assembly Constituency across four subsamples: (i) full sample; (ii) ACs with above-median rice-production area; (iii) ACs in Punjab and Haryana; and (iv) ACs in Punjab and Haryana with above-median rice-production area. We are using the farmer protests from June 2020 to December 2020, and draw a 2km buffer around each protest to define exposed and control areas across Assembly Constituencies (ACs). Exposure to Protest is coded as an absorbing treatment at the AC level. We also report the mean of the dependent variable for grid-cells outside the exposure buffers (non-exposed areas). Estimates come from a panel dataset from September 2012 to August 2022.  We include  grid fixed effects,  month$\times$year fixed effects, assembly fixed effects and the assembly government years fixed effects. Standard errors are clustered at  AC$\times$Exposure Area level. 
\end{minipage}
\end{table}
% \begin{table}[htbp!]
%    \caption{\label{tab:protest_treat_main_pj} Main Specification with Protest as Treatment 2km PJ and Haryana}
%    \centering \resizebox{0.7\textwidth}{!}{
% \input{tables/protest_2km_p_h}}
% \end{table}
\begin{table}[htbp!]
   \caption{\label{tab:app_protest_10km} Protest around 10km}
   \singlespacing \centering \resizebox{\textwidth}{!}{
\input{tables/protest_10km}}
\begin{minipage}{0.9\textwidth}
  \footnotesize
  \textbf{Notes:} The table reports how the number of fires changes when there is any farmer around 5km in the Assembly Constituency across four subsamples: (i) full sample; (ii) ACs with above-median rice-production area; (iii) ACs in Punjab and Haryana; and (iv) ACs in Punjab and Haryana with above-median rice-production area. We are using the farmer protests from June 2020 to December 2020, and draw a 10-km buffer around each protest to define exposed and control areas across Assembly Constituencies (ACs). Exposure to Protest is coded as an absorbing treatment at the AC level. We also report the mean of the dependent variable for grid-cells outside the exposure buffers (non-exposed areas). Estimates come from a panel dataset from September 2012 to August 2022.  We include  grid fixed effects,  month$\times$year fixed effects, assembly fixed effects and the assembly government years fixed effects. Standard errors are clustered at  AC$\times$Exposure Area level. 
\end{minipage}
\end{table}
% \begin{table}[htbp!]
%    \caption{\label{tab:protest_treat_main_pj} Main Specification with Protest as Treatment 10km PJ and Haryana}
%    \centering \resizebox{0.7\textwidth}{!}{
%    \input{tables/protest_10km_pj}
% }
% \end{table}



\begin{figure}[htbp!]
\centering
\begin{minipage}[b]{0.25\textwidth}
  \centering
  \includegraphics[width=\linewidth]{figures/10_pr_cluster_2022_2km.png}
  \caption*{ \scriptsize Full Sample}
\end{minipage}
\hfill
\begin{minipage}[b]{0.25\textwidth}
  \centering
  \includegraphics[width=\linewidth]{figures/10_pr_cluster_2022_2km_spec2.png}
  \caption*{ \scriptsize Full Sample - Assembly FE}
\end{minipage}
\hfill
\begin{minipage}[b]{0.25\textwidth}
  \centering
  \includegraphics[width=\linewidth]{figures/10_pr_cluster_2022_2km_spec3.png}
  \caption*{ \scriptsize Full Sample - Assembly $\times$ Gov. Year FE}
\end{minipage}
\vspace{1em}
\begin{minipage}[b]{0.25\textwidth}
  \centering
  \includegraphics[width=\linewidth]{figures/10_pr_cluster_2022_2km_riceareas.png}
  \caption*{ \scriptsize Above Median Rice Areas}
\end{minipage}
\hfill
\begin{minipage}[b]{0.25\textwidth}
  \centering
  \includegraphics[width=\linewidth]{figures/10_pr_cluster_2022_2km_riceareas_spec2.png}
  \caption*{ \scriptsize Above Median Rice Areas - Assembly FE}
\end{minipage}
\hfill
\begin{minipage}[b]{0.25\textwidth}
  \centering
  \includegraphics[width=\linewidth]{figures/10_pr_cluster_2022_2km_riceareas_spec3.png}
  \caption*{ \scriptsize Above Median Rice Areas - Assembly $\times$ Gov. Year FE}
\end{minipage}
\vspace{1em}
\begin{minipage}[b]{0.25\textwidth}
  \centering
  \includegraphics[width=\linewidth]{figures/10_pr_cluster_2022_2km_pjhy.png}
  \caption*{ \scriptsize Punjab and Haryana}
\end{minipage}
\hfill
\begin{minipage}[b]{0.25\textwidth}
  \centering
  \includegraphics[width=\linewidth]{figures/10_pr_cluster_2022_2km_pjhy_spec2.png}
  \caption*{ \scriptsize Punjab and Haryana - Assembly FE}
\end{minipage}
\hfill
\begin{minipage}[b]{0.25\textwidth}
  \centering
  \includegraphics[width=\linewidth]{figures/10_pr_cluster_2022_2km_pjhy_spec3.png}
  \caption*{ \scriptsize Punjab and Haryana - Assembly $\times$ Gov. Year FE}
\end{minipage}
\vspace{1em}
\begin{minipage}[b]{0.25\textwidth}
  \centering
  \includegraphics[width=\linewidth]{figures/10_pr_cluster_2022_2km_riceareas_pjhy.png}
  \caption*{ \scriptsize Punjab and Haryana Above Median Rice Areas}
\end{minipage}
\hfill
\begin{minipage}[b]{0.25\textwidth}
  \centering
  \includegraphics[width=\linewidth]{figures/10_pr_cluster_2022_2km_riceareas_pjhy_spec2.png}
  \caption*{ \scriptsize Punjab and Haryana Above Median Rice Areas - Assembly FE}
\end{minipage}
\hfill
\begin{minipage}[b]{0.25\textwidth}
  \centering
  \includegraphics[width=\linewidth]{figures/10_pr_cluster_2022_2km_riceareas_pjhy_spec3.png}
  \caption*{ \scriptsize Punjab and Haryana Above Median Rice Areas - Assembly $\times$ Gov. Year FE}
\end{minipage}
\vspace{1em}
\caption*{ Event Studies for Protest around 2Km}
\begin{minipage}{0.9\textwidth}
  
  \textbf{Notes:} \scriptsize This figure shows event studies of the effect of protest on number of fires. We use stack datasets generating cohort for every event. We are capturing the effect after the protests until the end of august 2022. The data is at monthly level, but the plots show effects by year. We use cohort fixed effects to control for repeated observations, and regressions are clustered by the exposure areas generated by the 2km buffer area around the protests. The columns shows the different set of fixed effects, and rows corresponds to the subsamples. 
\end{minipage}
\end{figure}



\begin{figure}[htbp!]
\centering
\begin{minipage}[b]{0.25\textwidth}
  \centering
  \includegraphics[width=\linewidth]{figures/10_pr_cluster_2022_10km.png}
  \caption*{ \scriptsize Full Sample}
\end{minipage}
\hfill
\begin{minipage}[b]{0.25\textwidth}
  \centering
  \includegraphics[width=\linewidth]{figures/10_pr_cluster_2022_10km_spec2.png}
  \caption*{ \scriptsize Full Sample - Assembly FE}
\end{minipage}
\hfill
\begin{minipage}[b]{0.25\textwidth}
  \centering
  \includegraphics[width=\linewidth]{figures/10_pr_cluster_2022_10km_spec3.png}
  \caption*{ \scriptsize Full Sample - Assembly $\times$ Gov. Year FE}
\end{minipage}
\vspace{1em}
\begin{minipage}[b]{0.25\textwidth}
  \centering
  \includegraphics[width=\linewidth]{figures/10_pr_cluster_2022_10km_riceareas.png}
  \caption*{ \scriptsize Above Median Rice Areas}
\end{minipage}
\hfill
\begin{minipage}[b]{0.25\textwidth}
  \centering
  \includegraphics[width=\linewidth]{figures/10_pr_cluster_2022_10km_riceareas_spec2.png}
  \caption*{ \scriptsize Above Median Rice Areas - Assembly FE}
\end{minipage}
\hfill
\begin{minipage}[b]{0.25\textwidth}
  \centering
  \includegraphics[width=\linewidth]{figures/10_pr_cluster_2022_10km_riceareas_spec3.png}
  \caption*{ \scriptsize Above Median Rice Areas - Assembly $\times$ Gov. Year FE}
\end{minipage}
\vspace{1em}
\begin{minipage}[b]{0.25\textwidth}
  \centering
  \includegraphics[width=\linewidth]{figures/10_pr_cluster_2022_10km_pjhy.png}
  \caption*{ \scriptsize Punjab and Haryana}
\end{minipage}
\hfill
\begin{minipage}[b]{0.25\textwidth}
  \centering
  \includegraphics[width=\linewidth]{figures/10_pr_cluster_2022_10km_pjhy_spec2.png}
  \caption*{ \scriptsize Punjab and Haryana - Assembly FE}
\end{minipage}
\hfill
\begin{minipage}[b]{0.25\textwidth}
  \centering
  \includegraphics[width=\linewidth]{figures/10_pr_cluster_2022_10km_pjhy_spec3.png}
  \caption*{ \scriptsize Punjab and Haryana - Assembly $\times$ Gov. Year FE}
\end{minipage}
\vspace{1em}
\begin{minipage}[b]{0.25\textwidth}
  \centering
  \includegraphics[width=\linewidth]{figures/10_pr_cluster_2022_10km_riceareas_pjhy.png}
  \caption*{ \scriptsize Punjab and Haryana Above Median Rice Areas}
\end{minipage}
\hfill
\begin{minipage}[b]{0.25\textwidth}
  \centering
  \includegraphics[width=\linewidth]{figures/10_pr_cluster_2022_10km_riceareas_pjhy_spec2.png}
  \caption*{ \scriptsize Punjab and Haryana Above Median Rice Areas - Assembly FE}
\end{minipage}
\hfill
\begin{minipage}[b]{0.25\textwidth}
  \centering
  \includegraphics[width=\linewidth]{figures/10_pr_cluster_2022_10km_riceareas_pjhy_spec3.png}
  \caption*{ \scriptsize Punjab and Haryana Above Median Rice Areas - Assembly $\times$ Gov. Year FE}
\end{minipage}
\vspace{1em}
\caption*{ Event Studies for Protest around 10km}
\begin{minipage}{0.9\textwidth}
  
  \textbf{Notes:} \scriptsize This figure shows event studies of the effect of protest on number of fires. We use stack datasets generating cohort for every event. We are capturing the effect after the protests until the end of august 2022. The data is at monthly level, but the plots show effects by year. We use cohort fixed effects to control for repeated observations, and regressions are clustered by the exposure areas generated by the 10km buffer area around the protests. The columns shows the different set of fixed effects, and rows corresponds to the subsamples. 
\end{minipage}
\end{figure}

\begin{figure}[htbp!]
\centering
\begin{minipage}[b]{0.3\textwidth}
  \centering
  \includegraphics[width=\linewidth]{figures/self_prof.png}
  \caption*{\scriptsize Agricultural Profession (1)}
\end{minipage}
\begin{minipage}[b]{0.3\textwidth}
  \centering
  \includegraphics[width=\linewidth]{figures/self_prof_spec2.png}
  \caption*{\scriptsize Agricultural Profession (2)}
\end{minipage}
\begin{minipage}[b]{0.3\textwidth}
  \centering
  \includegraphics[width=\linewidth]{figures/self_prof_spec3.jpeg}
  \caption*{\scriptsize Agricultural Profession (3)}
\end{minipage}
\vspace{1cm}
\begin{minipage}[b]{0.3\textwidth}
  \centering
  \includegraphics[width=\linewidth]{figures/self_prof_assets.png}
  \caption*{\scriptsize Agricultural Assets and Profession (1)}
\end{minipage}
\begin{minipage}[b]{0.3\textwidth}
  \centering
  \includegraphics[width=\linewidth]{figures/self_prof_assets_spec2.png}
  \caption*{\scriptsize Agricultural Assets and Profession (2)}
\end{minipage}
\begin{minipage}[b]{0.3\textwidth}
  \centering
  \includegraphics[width=\linewidth]{figures/self_prof_assets_spec3.jpeg}
  \caption*{\scriptsize Agricultural Assets and Profession (3)}
\end{minipage}
\vspace{1cm}
\begin{minipage}{0.9\textwidth}
  \footnotesize
  \textbf{Notes:} These figures present event study estimates based on whether the winning candidate had a profession related to agriculture or held agricultural assets. The analysis uses stacked dataset and is conducted at the grid-month-cohort level, while event time is defined in terms of years before and after the start of the elected government's term, anchored to the month when the winning candidate assumed office. The plots correspond to the columns specification on Table \ref{tab:stacked_pol_dd} where we include  grid$\times$cohort fixed effects,  month$\times$year$\times$cohort fixed effects, assembly fixed effects and the assembly government years fixed effects. Standard errors are clustered at the level of grid and AC$\times$month$\times$year. It uses 1,088,242 observations of our original assets since int he construction of the cohorts we drop observations which do not experiment a change. It represents around 9\% of our data. \red{Keep only profession and all the set of fixed effects for now}
\end{minipage}
\end{figure}

\begin{table}[htbp!]
   \caption{\label{tab:stacked_pol_dd} Effect of Politician Characteristics on Number of Fires}
   \singlespacing
   \centering \resizebox{\textwidth}{!}{
   \input{tables/_edit_pols_char_didtable}}
   \begin{minipage}{0.9\textwidth}
  \footnotesize
  \textbf{Notes:} The table reports how the number of fires changes when the winning legislator has an agricultural link—through (i) an agriculture-related occupation, and (ii) ownership of agricultural assets. Estimates come from a stacked panel that spans the five years before and after each election in every state, aligned to the month the victor assumed office. All specifications include grid$\times$cohort fixed effects and month$\times$year$\times$cohort fixed effects. Standard errors are clustered at the level of grid and AC$\times$month$\times$year. \red{SWITCHING - we keep only profession + professions and assets. The appendix has this result for agricultural assets. This information is in \ref{tab:stacked_pol_dd_assets}  }
\end{minipage}
\end{table}

\begin{figure}[htbp!]
\centering
\begin{minipage}[b]{0.3\textwidth}
  \centering
  \includegraphics[width=\linewidth]{figures/10_pr_cluster_2022_5km.png}
  \caption*{\scriptsize Full Sample}
\end{minipage}
\hfill
\begin{minipage}[b]{0.3\textwidth}
  \centering
  \includegraphics[width=\linewidth]{figures/10_pr_cluster_2022_5km_spec2.png}
  \caption*{ \scriptsize Full Sample - Assembly FE}
\end{minipage}
\hfill
\begin{minipage}[b]{0.3\textwidth}
  \centering
  \includegraphics[width=\linewidth]{figures/10_pr_cluster_2022_5km_spec3.png}
  \caption*{ \scriptsize Full Sample - Assembly $\times$ Gov. Year FE}
\end{minipage}
\vspace{1em}
\begin{minipage}[b]{0.3\textwidth}
  \centering
  \includegraphics[width=\linewidth]{figures/10_pr_cluster_2022_5km_riceareas.png}
  \caption*{ \scriptsize Above Median Rice Areas}
\end{minipage}
\hfill
\begin{minipage}[b]{0.3\textwidth}
  \centering
  \includegraphics[width=\linewidth]{figures/10_pr_cluster_2022_5km_riceareas_spec2.png}
  \caption*{ \scriptsize Above Median Rice Areas - Assembly FE}
\end{minipage}
\hfill
\begin{minipage}[b]{0.3\textwidth}
  \centering
  \includegraphics[width=\linewidth]{figures/10_pr_cluster_2022_5km_riceareas_spec3.png}
  \caption*{\scriptsize Above Median Rice Areas - Assembly $\times$ Gov. Year FE}
\end{minipage}
\vspace{1em}
\begin{minipage}[b]{0.3\textwidth}
  \centering
  \includegraphics[width=\linewidth]{figures/10_pr_cluster_2022_5km_pjhy.png}
  \caption*{Punjab and Haryana}
\end{minipage}
\hfill
\begin{minipage}[b]{0.3\textwidth}
  \centering
  \includegraphics[width=\linewidth]{figures/10_pr_cluster_2022_5km_pjhy_spec2.png}
  \caption*{ \scriptsize Punjab and Haryana - Assembly FE}
\end{minipage}
\hfill
\begin{minipage}[b]{0.3\textwidth}
  \centering
  \includegraphics[width=\linewidth]{figures/10_pr_cluster_2022_5km_pjhy_spec3.png}
  \caption*{\scriptsize Punjab and Haryana - Assembly $\times$ Gov. Year FE}
\end{minipage}
\vspace{1em}
\begin{minipage}[b]{0.3\textwidth}
  \centering
  \includegraphics[width=\linewidth]{figures/10_pr_cluster_2022_5km_riceareas_pjhy.png}
  \caption*{\scriptsize Punjab and Haryana Above Median Rice Areas}
\end{minipage}
\hfill
\begin{minipage}[b]{0.3\textwidth}
  \centering
  \includegraphics[width=\linewidth]{figures/10_pr_cluster_2022_5km_riceareas_pjhy_spec2.png}
  \caption*{\scriptsize Punjab and Haryana Above Median Rice Areas - Assembly FE}
\end{minipage}
\hfill
\begin{minipage}[b]{0.3\textwidth}
  \centering
  \includegraphics[width=\linewidth]{figures/10_pr_cluster_2022_5km_riceareas_pjhy_spec3.png}
  \caption*{ \scriptsize Punjab and Haryana Above Median Rice Areas - Assembly $\times$ Gov. Year FE}
\end{minipage}
\vspace{1em}
\begin{minipage}{0.9\textwidth}
  
  \textbf{Notes:} \footnotesize This figure shows event studies of the effect of protest on number of fires. We use stack datasets generating cohort for every event. We are capturing the effect after the protests until the end of august 2022. The data is at monthly level, but the plots show effects by year. We use cohort fixed effects to control for repeated observations, and regressions are clustered by the exposure areas generated by the 5km buffer area around the protests. The columns shows the different set of fixed effects, and rows corresponds to the subsamples. 
\end{minipage}
\end{figure}


\begin{table}[htbp!]
   \caption{\label{tab:protest_treat_main} Effect of Level of Protests on Number of Fires 5km specification}
   \singlespacing \centering \resizebox{\textwidth}{!}{
\begin{table}[htbp]
   \caption{\label{tab:protest_treat_main} Effect of of Protests around 5km on Number of Fires}
   \centering
   \begin{tabular}{lcccccccccccccccccccccccc}
      \tabularnewline \midrule \midrule
       & \multicolumn{24}{c}{Number of Fires (1000)}\\
                                & (1)        & (2)           & (3)           & (4)           & (5)           & (6)           & (7)         & (8)           & (9)           & (10)          & (11)          & (12)         & (13)         & (14)         & (15)        & (16)          & (17)          & (18)         & (19)      & (20)      & (21)      & (22)          & (23)      & (24)\\  
      \midrule 
      Protest 5km               & 9.329      & 66.91$^{***}$ & 33.96$^{***}$ & 68.20$^{***}$ & 50.96$^{***}$ & 28.22$^{***}$ & 20.82$^{*}$ & 99.49$^{***}$ & 58.36$^{***}$ & 95.61$^{***}$ & 71.40$^{***}$ & 30.94$^{**}$ & 32.08$^{**}$ & 28.11$^{**}$ & 21.91$^{*}$ & 127.5$^{***}$ & 39.10$^{***}$ & 25.21$^{**}$ & 16.91     & 12.41     & 7.248     & 163.3$^{***}$ & 19.61     & 9.995\\   
                                & (7.059)    & (9.795)       & (8.652)       & (9.391)       & (9.242)       & (8.454)       & (11.08)     & (15.55)       & (14.26)       & (15.29)       & (14.54)       & (13.10)      & (13.10)      & (12.43)      & (12.04)     & (14.16)       & (13.89)       & (12.13)      & (20.25)   & (18.80)   & (18.31)   & (21.54)       & (21.50)   & (18.39)\\   
      \hline Mean Y0            & 132.8      & 132.8         & 132.8         & 132.8         & 132.8         & 132.8         & 176.3       & 176.3         & 176.3         & 176.3         & 176.3         & 176.3        & 451.4        & 451.4        & 451.4       & 451.4         & 451.4         & 451.4        & 584.1     & 584.1     & 584.1     & 584.1         & 584.1     & 584.1\\  
      Assembly                  & 853        & 853           & 853           & 853           & 853           & 853           & 410         & 410           & 410           & 410           & 410           & 410          & 207          & 207          & 207         & 207           & 207           & 207          & 92        & 92        & 92        & 92            & 92        & 92\\  
      \midrule
      \emph{Fixed-effects}\\
      Grid                      & Y          & Y             & Y             & Y             & Y             & Y             & Y           & Y             & Y             & Y             & Y             & Y            & Y            & Y            & Y           & Y             & Y             & Y            & Y         & Y         & Y         & Y             & Y         & Y\\  
      month-year                & Y          & Y             & Y             & N             & Y             & Y             & Y           & Y             & Y             & N             & Y             & Y            & Y            & Y            & Y           & N             & Y             & Y            & Y         & Y         & Y         & N             & Y         & Y\\  
      Assembly FE               & N          & Y             & N             & N             & N             & N             & N           & Y             & N             & N             & N             & N            & N            & Y            & N           & N             & N             & N            & N         & Y         & N         & N             & N         & N\\  
      yeargov                   & N          & Y             & N             & N             & N             & N             & N           & Y             & N             & N             & N             & N            & N            & Y            & N           & N             & N             & N            & N         & Y         & N         & N             & N         & N\\  
      Assembly Gov. Year FE     & N          & N             & Y             & N             & N             & Y             & N           & N             & Y             & N             & N             & Y            & N            & N            & Y           & N             & N             & Y            & N         & N         & Y         & N             & N         & Y\\  
      province                  & N          & N             & N             & Y             & Y             & Y             & N           & N             & N             & Y             & Y             & Y            & N            & N            & N           & Y             & Y             & Y            & N         & N         & N         & Y             & Y         & Y\\  
      \midrule \emph{Time Trends}\\
      MonthYear (province)      & N          & N             & N             & Y             & Y             & Y             & N           & N             & N             & Y             & Y             & Y            & N            & N            & N           & Y             & Y             & Y            & N         & N         & N         & Y             & Y         & Y\\  
      \midrule
      \emph{Fit statistics}\\
      \# Grid                   & 100,253    & 100,253       & 100,253       & 100,253       & 100,253       & 100,253       & 41,653      & 41,653        & 41,653        & 41,653        & 41,653        & 41,653       & 22,601       & 22,601       & 22,601      & 22,601        & 22,601        & 22,601       & 9,997     & 9,997     & 9,997     & 9,997         & 9,997     & 9,997\\  
      \# month-year             & 120        & 120           & 120           & --            & 120           & 120           & 120         & 120           & 120           & --            & 120           & 120          & 120          & 120          & 120         & --            & 120           & 120          & 120       & 120       & 120       & --            & 120       & 120\\  
      \# Assembly FE            & --         & 12            & --            & --            & --            & --            & --          & 12            & --            & --            & --            & --           & --           & 6            & --          & --            & --            & --           & --        & 6         & --        & --            & --        & --\\  
      \# yeargov                & --         & 5             & --            & --            & --            & --            & --          & 5             & --            & --            & --            & --           & --           & 5            & --          & --            & --            & --           & --        & 5         & --        & --            & --        & --\\  
      \# Assembly Gov. Year FE  & --         & --            & 44            & --            & --            & 44            & --          & --            & 44            & --            & --            & 44           & --           & --           & 22          & --            & --            & 22           & --        & --        & 22        & --            & --        & 22\\  
      \# province               & --         & --            & --            & 4             & 4             & 4             & --          & --            & --            & 4             & 4             & 4            & --           & --           & --          & 2             & 2             & 2            & --        & --        & --        & 2             & 2         & 2\\  
      Observations              & 12,030,360 & 12,030,360    & 12,030,360    & 12,030,360    & 12,030,360    & 12,030,360    & 4,998,360   & 4,998,360     & 4,998,360     & 4,998,360     & 4,998,360     & 4,998,360    & 2,712,120    & 2,712,120    & 2,712,120   & 2,712,120     & 2,712,120     & 2,712,120    & 1,199,640 & 1,199,640 & 1,199,640 & 1,199,640     & 1,199,640 & 1,199,640\\  
      \midrule \midrule
      \multicolumn{25}{l}{\emph{Clustered (ac\_area\_tr) standard-errors in parentheses}}\\
      \multicolumn{25}{l}{\emph{Signif. Codes: ***: 0.01, **: 0.05, *: 0.1}}\\
   \end{tabular}
\end{table}


}
   \begin{minipage}{0.9\textwidth}
  \footnotesize
  \textbf{Notes:} The table reports how the number of fires changes when there is any farmer around 5km in the Assembly Constituency across four subsamples: (i) full sample; (ii) ACs with above-median rice-production area; (iii) ACs in Punjab and Haryana; and (iv) ACs in Punjab and Haryana with above-median rice-production area. We are using the farmer protests from June 2020 to December 2020, and draw a 5-km buffer around each protest to define exposed and control areas across Assembly Constituencies (ACs). Exposure to Protest is coded as an absorbing treatment at the AC level. We also report the mean of the dependent variable for grid-cells outside the exposure buffers (non-exposed areas). Estimates come from a panel dataset from September 2012 to August 2022.  We include  grid fixed effects,  month$\times$year fixed effects, assembly fixed effects and the assembly government years fixed effects. Standard errors are clustered at  AC$\times$Exposure Area level. 
\end{minipage}
\end{table}

\begin{figure}[htbp!]
\centering
\begin{minipage}[b]{0.9\textwidth}
  \centering
  \includegraphics[width=\linewidth]{figures/ac_level_plot.png}
\end{minipage}

\begin{minipage}{0.9\textwidth}
  \footnotesize
  \textbf{Notes:} This figure shows ACs where protest took place. \red{Add AC boundaries with very thin lines} 
\end{minipage}
\end{figure}

\section{@Anzony: to Check Removed analysis}
\begin{itemize}
    \item Analysis with the double interaction of downup AC with rice, low agricultural workers, harvest months and their combinations. [difficult to interpret the results]
    \item Eventstudy with the triple interaction with downup AC and Downup Distr
    \item Personal Characteristics: DD analysis
\end{itemize}

% Ac X Rice
% check value of y outcome y0 in column (5)
\begin{table}[htbp!]
     \caption{Effect of constituency to smoke on N Fires, by rice producing areas} \singlespacing \centering
      \resizebox{0.8\textwidth}{!}{
      \input{tables/ac_rice}
      }
     \begin{minipage}{\textwidth}
         {\small \textit{Notes: }$Rice$ is defined as having a larger than median area of the grid devoted to rice crop production at AC. The table reports the mean of the dependent variable for grids that are upwind with respect to borders of the constituency in low rice producing areas. All regressions control for average wind speed and direction. Standard errors are clustered at the grid and district times month-year level.}
     \end{minipage}
 \end{table}

%% Event Study from Treatment to Control using -5 to +6
\begin{figure}[htbp!]
\caption{Event Study Downup $\times$ Rice using -5 pre treatment observations}
    \label{fig:combined_plot2}
    \centering
    \begin{subfigure}[b]{0.45\textwidth}
        \centering
        \includegraphics[width=\textwidth]{figures/NOT_ROTATED_eventStudy_acxrice5.png}
        \caption[]{ Not Rotated}
        \label{fig:treat2control_not_rotated}
    \end{subfigure}%
    \begin{subfigure}[b]{0.45\textwidth}
        \centering
        \includegraphics[width=\textwidth]{figures/ROTATED_eventStudy_acxrice5.png}
        \caption[]{Rotated}
        \label{fig:treat2control_rotated}
    \end{subfigure}
    
    \begin{subfigure}[b]{0.45\textwidth}
        \centering
        \includegraphics[width=\textwidth]{figures/HONEST_eventStudy_estimates_acxrice5.png}
        \caption[]{Honest DiD}
        \label{fig:treat2control_honest}
    \end{subfigure}
     \begin{minipage}{\textwidth}
         {\small \textit{Notes: }This plot shows the differential effect for Rice areas. we are polotting the original event study and the rotated version. This version takes into account the pre treatment trend and extrapolates to post treatment result. Finally, we follow rambachan and roth generating the plot of bounded estimation of the post treatment. we set bounds on how large post-treatment deviations could be, we have to assume a very large bound to take them as null effect, relative to the biggest differences pre-treatment.  }
     \end{minipage}
\end{figure}

\newpage

% Ac X Low Agricultural Workers
\begin{table}[htbp!]
     \caption{Effect of constituency and district exposure to smoke on N Fires, by areas with low levels of working population in agriculture } \singlespacing \centering
      \resizebox{0.8\textwidth}{!}{
          \input{tables/ac_agriwork}
      }
     \begin{minipage}{\textwidth}
         {\small \textit{Notes: }$Low \quad Agriculture \quad Workers$ is defined as having a share of the working population in agriculture lower than median share in the sample. The table reports the mean of the dependent variable for grids that are upwind with respect to borders of  the constituency in constituencies with high level of working population in agriculture. All regressions control for average wind speed and direction. Standard errors are clustered at the grid and AC times month-year level. The last column report the effect calculating the variable using village dataset. }
     \end{minipage}
 \end{table}

%% Event Study from Treatment to Control using -5 to +6
\begin{figure}[htbp!]
\caption{Event Study Downup $\times$ Low Minority Agricultural Workers using -5 pre treatment observations}
    \label{fig:combined_plot2}
    \centering
    \begin{subfigure}[b]{0.45\textwidth}
        \centering
        \includegraphics[width=\textwidth]{figures/NOT_ROTATED_eventStudy_acxminagri5.png}
        \caption[]{ Not Rotated}
        \label{fig:treat2control_not_rotated}
    \end{subfigure}%
    \begin{subfigure}[b]{0.45\textwidth}
        \centering
        \includegraphics[width=\textwidth]{figures/ROTATED_eventStudy_acxminagri5.png}
        \caption[]{Rotated}
        \label{fig:treat2control_rotated}
    \end{subfigure}
    
    \begin{subfigure}[b]{0.45\textwidth}
        \centering
        \includegraphics[width=\textwidth]{figures/HONEST_eventStudy_estimates_acxminagri5.png}
        \caption[]{Honest DiD}
        \label{fig:treat2control_honest}
    \end{subfigure}
     \begin{minipage}{\textwidth}
         {\small \textit{Notes: }This plot shows the differential effect for ACs below median number of agricultural workers. we are plotting the original event study and the rotated version. This version takes into account the pre treatment trend and extrapolates to post treatment result. Finally, we follow rambachan and roth generating the plot of bounded estimation of the post treatment. we set bounds on how large post-treatment deviations could be, we have to assume a very large bound to take them as null effect, relative to the biggest differences pre-treatment.  }
     \end{minipage}
\end{figure}

\newpage

%did table
\begin{table}[htbp!]
     \caption{Effect of constituency exposure to smoke on N Fires \label{tab:ac_myneta}} \singlespacing \centering
      \resizebox{0.7\textwidth}{!}{
     \input{tables/ac_myneta}
     }
     \begin{minipage}{\textwidth}
         {\small \textit{Notes: }The table reports the mean of the dependent variable in grids that are upwind in constituencies where the politician has no agricultural assets or job related to agriculture.  All regressions control for average wind speed and direction. Standard errors are clustered at the grid and AC times month-year level. \red{Should we get rif of this or move it to the appendix?} }
     \end{minipage}
 \end{table}
 % % Choropleth Maps We comment thsi since this is the highway analysis
\begin{figure}[htbp!]
    \centering
    \includegraphics[width=0.8\linewidth]{figures/demonstrations_choropleth.png}
    \caption{Protests at AC level}
    \label{fig:protest_ac}
\vspace{1em} % Add some vertical space between the figures and the note
\begin{minipage}{0.9\textwidth}
  \footnotesize
  \textbf{Notes:} This figure shows the number of protests at the Assembly Constituency (AC) level. The raw data comes from \href{https://acleddata.com/}{ACLED} and contains two key descriptive variables: actors and notes. These variables are used to identify farmer-related protests. The selection strategy involves searching for any conjugation or variation of the root word "farm" within either of these two fields. If such a match is found, the protest is included in our sample. Each protest event is geolocated using latitude and longitude, which allows us to spatially intersect the events with AC boundary shapefiles to assign them to specific constituencies. \red{Make value 0 = white and the colors darker + Add AC boundaries with very thin lines}
\end{minipage}
\end{figure}
 % \begin{figure}[htbp!]
% \caption{ACs where the Elected Legislator has Agricultural Assets}
%     \label{fig:combined_plot_assets}
%     \centering
%     \begin{subfigure}[b]{0.3\textwidth}
%         \centering
%         \includegraphics[width=\textwidth]{figures/NOT_ROTATED_eventStudy_count_self_agri_assets.png}
%         \caption[]{Not Rotated}
%         \label{fig:assets_not_rotated}
%     \end{subfigure}%
%     \begin{subfigure}[b]{0.3\textwidth}
%         \centering
%         \includegraphics[width=\textwidth]{figures/ROTATED_eventStudy_count_self_agri_assets.png}
%         \caption[]{ Rotated}
%         \label{fig:assets_rotated}
%     \end{subfigure}
%     \begin{subfigure}[b]{0.3\textwidth}
%         \centering
%         \includegraphics[width=\textwidth]{figures/HONEST_eventStudy_estimates_count_self_agri_assets.png}
%         \caption[]{ Honest DiD}
%         \label{fig:assets_honest}
%     \end{subfigure}
% \end{figure}

% \begin{figure}[htbp!]
%     \caption{ACs where the Elected Legislator's Profession is related with Agriculture}
%     \label{fig:combined_plot_profession}
%     \centering
%     \begin{subfigure}[b]{0.3\textwidth}
%         \centering
%         \includegraphics[width=\textwidth]{figures/NOT_ROTATED_eventStudy_count_self_agri_prof.png}
%         \caption[]{Not Rotated}
%         \label{fig:profession_not_rotated}
%     \end{subfigure}
%     \begin{subfigure}[b]{0.3\textwidth}
%         \centering
%         \includegraphics[width=\textwidth]{figures/ROTATED_eventStudy_count_self_agri_prof.png}
%         \caption[]{Rotated}
%         \label{fig:profession_rotated}
%     \end{subfigure}
%     \begin{subfigure}[b]{0.3\textwidth}
%         \centering
%         \includegraphics[width=\textwidth]{figures/HONEST_eventStudy_estimates_count_self_agri_prof.png}
%         \caption[]{Honest DiD}
%         \label{fig:profession_honest}
%     \end{subfigure}
% \end{figure}

% \begin{table}[htbp!]
%    \caption{\label{tab:protest_treat_main2fe} Effect of Level of Protests on Number of Fires with Electoral Cycle Fixed Effects}
%    \centering \resizebox{\textwidth}{!}{
\begin{table}[htbp]
   \caption{\label{tab:protest_treat_main} Effect of Level of Protests on Number of Fires}
   \centering
   \begin{tabular}{lcccccccccccc}
      \tabularnewline \midrule \midrule
       & \multicolumn{12}{c}{Number of Fires}\\
                                            & (1)           & (2)           & (3)           & (4)           & (5)           & (6)           & (7)           & (8)           & (9)           & (10)          & (11)          & (12)\\  
      \midrule
      \emph{Variables}\\
      Protest Absr $=$ 1                    & 0.1185        & 0.1185        & 0.1199        & 0.1913        & 0.1913        & 0.1941        & 0.1034        & 0.1034        & 0.1040        & 0.0772        & 0.0771        & 0.0794\\   
                                            & (0.0185)      & (0.0185)      & (0.0219)      & (0.0315)      & (0.0315)      & (0.0379)      & (0.0384)      & (0.0384)      & (0.0385)      & (0.0636)      & (0.0636)      & (0.0632)\\   
      DownupAC                              &               & -0.0090       & -0.0088       &               & -0.0094       & -0.0090       &               & -0.0147       & -0.0145       &               & -0.0094       & -0.0086\\   
                                            &               & (0.0028)      & (0.0018)      &               & (0.0038)      & (0.0027)      &               & (0.0084)      & (0.0087)      &               & (0.0113)      & (0.0117)\\   
      Protest Absr $\times$ DownupAC        &               &               & -0.0028       &               &               & -0.0057       &               &               & -0.0012       &               &               & -0.0046\\   
                                            &               &               & (0.0064)      &               &               & (0.0101)      &               &               & (0.0085)      &               &               & (0.0118)\\   
      \hline Mean Y0                        & 0.1253        & 0.1253        & 0.1253        & 0.1650        & 0.1650        & 0.1650        & 0.4432        & 0.4432        & 0.4432        & 0.6176        & 0.6176        & 0.6176\\  
      \midrule
      \emph{Fixed-effects}\\
      Grid                                  & $\checkmark$  & $\checkmark$  & $\checkmark$  & $\checkmark$  & $\checkmark$  & $\checkmark$  & $\checkmark$  & $\checkmark$  & $\checkmark$  & $\checkmark$  & $\checkmark$  & $\checkmark$\\   
      MonthYear                             & $\checkmark$  & $\checkmark$  & $\checkmark$  & $\checkmark$  & $\checkmark$  & $\checkmark$  & $\checkmark$  & $\checkmark$  & $\checkmark$  & $\checkmark$  & $\checkmark$  & $\checkmark$\\   
      Province $\times$ Election Cycle      & $\checkmark$  & $\checkmark$  & $\checkmark$  & $\checkmark$  & $\checkmark$  & $\checkmark$  & $\checkmark$  & $\checkmark$  & $\checkmark$  & $\checkmark$  & $\checkmark$  & $\checkmark$\\   
      \midrule
      \emph{Fit statistics}\\
      \# Grid                               & 100,253       & 100,253       & 100,253       & 57,570        & 57,570        & 57,570        & 22,601        & 22,601        & 22,601        & 12,808        & 12,808        & 12,808\\  
      \# MonthYear                          & 120           & 120           & 120           & 120           & 120           & 120           & 120           & 120           & 120           & 120           & 120           & 120\\  
      \# Province $\times$ Election Cycle   & 12            & 12            & 12            & 12            & 12            & 12            & 6             & 6             & 6             & 6             & 6             & 6\\  
      Observations                          & 12,030,360    & 12,030,360    & 12,030,360    & 6,908,400     & 6,908,400     & 6,908,400     & 2,712,120     & 2,712,120     & 2,712,120     & 1,536,960     & 1,536,960     & 1,536,960\\  
      \midrule \midrule
                                            &               &               &               &               &               &               &               &               &               &               &               &  \tabularnewline 
      \\
   \end{tabular}
\end{table}


}   
% \end{table}

% \subsection{Ovelapping Jurisdictions}

% While the results show a reduction in fires when pollution is directly and visibly affecting the politician's jurisdiction, this incentive shrinks when the fires would pollute also the bureaucts's jurisdiction. this incentive weakens when the smoke also affects the bureaucrat’s jurisdiction. As shown in Table~\ref{tab:main_ac_distr}, we observe a 6\% reduction in effort, consistent with a coordination failure in which each actor hopes the other will absorb the cost of enforcement. As a result, the reduction in fires attributable to politicians' incentives effectively disappears in downwind areas shared with bureaucrats.

\begin{figure}[htbp!]
\centering
\begin{minipage}[b]{0.25\textwidth}
  \centering
  \includegraphics[width=\linewidth]{figures/9_ac_treatment_level.png}
  \caption*{ \scriptsize Full Sample}
\end{minipage}
\hfill
\begin{minipage}[b]{0.25\textwidth}
  \centering
  \includegraphics[width=\linewidth]{figures/9_ac_treatment_level_spec2.png}
  \caption*{ \scriptsize Full Sample - Assembly FE}
\end{minipage}
\hfill
\begin{minipage}[b]{0.25\textwidth}
  \centering
  \includegraphics[width=\linewidth]{figures/9_ac_treatment_level_spec3.png}
  \caption*{ \scriptsize Full Sample - Assembly $\times$ Gov. Year FE}
\end{minipage}
\vspace{1em}
\begin{minipage}[b]{0.25\textwidth}
  \centering
  \includegraphics[width=\linewidth]{figures/9_ac_treatment_level_riceareas.png}
  \caption*{ \scriptsize Above Median Rice Areas}
\end{minipage}
\hfill
\begin{minipage}[b]{0.25\textwidth}
  \centering
  \includegraphics[width=\linewidth]{figures/9_ac_treatment_level_riceareas_spec2.png}
  \caption*{ \scriptsize Above Median Rice Areas - Assembly FE}
\end{minipage}
\hfill
\begin{minipage}[b]{0.25\textwidth}
  \centering
  \includegraphics[width=\linewidth]{figures/9_ac_treatment_level_riceareas_spec3.png}
  \caption*{ \scriptsize Above Median Rice Areas - Assembly $\times$ Gov. Year FE}
\end{minipage}
\vspace{1em}
\begin{minipage}[b]{0.25\textwidth}
  \centering
  \includegraphics[width=\linewidth]{figures/9_ac_treatment_level_pjhy.png}
  \caption*{ \scriptsize Punjab and Haryana}
\end{minipage}
\hfill
\begin{minipage}[b]{0.25\textwidth}
  \centering
  \includegraphics[width=\linewidth]{figures/9_ac_treatment_level_pjhy_spec2.png}
  \caption*{ \scriptsize Punjab and Haryana - Assembly FE}
\end{minipage}
\hfill
\begin{minipage}[b]{0.25\textwidth}
  \centering
  \includegraphics[width=\linewidth]{figures/9_ac_treatment_level_pjhy_spec3.png}
  \caption*{ \scriptsize Punjab and Haryana - Assembly $\times$ Gov. Year FE}
\end{minipage}
\vspace{1em}
\begin{minipage}[b]{0.25\textwidth}
  \centering
  \includegraphics[width=\linewidth]{figures/9_ac_treatment_level_riceareas_pjhy.png}
  \caption*{ \scriptsize Punjab and Haryana Above Median Rice Areas}
\end{minipage}
\hfill
\begin{minipage}[b]{0.25\textwidth}
  \centering
  \includegraphics[width=\linewidth]{figures/9_ac_treatment_level_riceareas_pjhy_spec2.png}
  \caption*{ \scriptsize Punjab and Haryana Above Median Rice Areas - Assembly FE}
\end{minipage}
\hfill
\begin{minipage}[b]{0.25\textwidth}
  \centering
  \includegraphics[width=\linewidth]{figures/9_ac_treatment_level_riceareas_pjhy_spec3.png}
  \caption*{ \scriptsize Punjab and Haryana Above Median Rice Areas - Assembly $\times$ Gov. Year FE}
\end{minipage}
\vspace{1em}
\begin{minipage}{0.9\textwidth}
  \textbf{Notes:} \scriptsize This figure shows event studies of the effect of protest on number of fires. We use stack datasets generating cohort for every event. The treated areas  are generated based on the AC where any protest took place. After the pixel gets treated it remains treated for the whole period of analysis. We are capturing the effect until the end of august 2022. The data is at monthly level, but the plots show effects by year. We use cohort fixed effects to control for repeated observations, and regressions are clustered by AC. The columns shows the different set of fixed effects, and rows corresponds to the subsamples. 
\end{minipage}
\end{figure}





% \begin{table}[htbp!]
%    \caption{\label{tab:protest_treat_main_pj} Main Specification with Protest as Treatment at AC Level}
%    \centering \resizebox{0.7\textwidth}{!}{ \input{tables/8_ac_treatment_level}
% }
% \end{table}
% \begin{table}[htbp!]
%    \caption{\label{tab:protest_treat_main_pj} Main Specification with Protest as Treatment at Ac Level PJ and Haryana}
%    \centering \resizebox{0.7\textwidth}{!}{ \input{tables/8_ac_treatment_level_pjhy}
%  }
% \end{table}





\section{Results Using Area-Based Treatment Definition}
\label{sec:app_downup_ac}

\begin{figure}[htbp!]
  \centering
  \caption{Effect of constituency exposure to smoke on number of fires (Area-Based Treatment)}\label{fig:app_main_panel_area}
  \subcaptionbox{Main Table}[0.48\textwidth]{%
    \begin{adjustbox}{max width=\linewidth}
      {
\def\sym#1{\ifmmode^{#1}\else\(^{#1}\)\fi}
\begin{tabular}{l*{4}{c}}
\hline
            &\multicolumn{1}{c}{(1)}         &\multicolumn{1}{c}{(2)}         &\multicolumn{1}{c}{(3)}         &\multicolumn{1}{c}{(4)}         \\
            & \multicolumn{4}{c}{Number of Fires (in 1,000 units)} \\ \hline

Down $>$ Up &       5.677\sym{***}&       1.067         &     -34.243\sym{***}&      -2.914\sym{*}  \\
            &     (2.082)         &     (1.536)         &     (2.349)         &     (1.638)         \\
\hline
Observations&  36,866,365         &  36,866,365         &  36,866,365         &  36,866,365         \\
N Assembly Constituencies&         798         &         798         &         798         &         798         \\
Month-Year $\times$ Cohort FE&           N         &           Y         &           Y         &           N         \\
Grid $\times$ Cohort FE&           N         &           N         &           Y         &           Y         \\
AC $\times$ Month-Year $\times$ Cohort FE&           N         &           N         &           N         &           Y         \\
Mean DV     &      91.271         &      91.271         &      91.271         &      91.271         \\
\hline
\end{tabular}
}

    \end{adjustbox}
  }
  \subcaptionbox{Event-study relative to exposure}[0.48\textwidth]{%
    \includegraphics[width=\linewidth]{figures/stacked_event_study_5pre_rural_1_rotated.png}
  }
  \subcaptionbox{Effect of constituency and district exposure to smoke}[\textwidth]{%
    \begin{adjustbox}{max width=.7\linewidth}
      \input{tables/_main_3_bureau_polisc_did_rural}
    \end{adjustbox}
  }
\end{figure}

\begin{figure}[htbp!]
  \centering
  \caption{Effect of constituency exposure to smoke on number of fires, by constituency with rice production (Area-Based Treatment)}\label{fig:app_es_rice_area}
  \centering\includegraphics[width=0.9\linewidth]{figures/main_event_study_rural_riceP_rotated.png}
\end{figure}

\begin{figure}[htbp!]
    \caption{Effect by relative time period from exposure (Area-Based Treatment)}
    \label{fig:app_es_main_rotated_panel_area}
    \centering
    \begin{tabular}{ccc}
        \includegraphics[width=0.300\textwidth]{figures/stacked_event_study_5pre_rural_1_ori.png} &
        \includegraphics[width=0.300\textwidth]{figures/stacked_event_study_5pre_rural_1_rotated.png} &
        \includegraphics[width=0.300\textwidth]{figures/stacked_event_study_5pre_rural_1_rot_honest2.png} \\
        \small (a) Event-Study & \small (b) De-Trended Event-Study & \small (c) Honest DiD \\
    \end{tabular}
\end{figure}

\begin{table}[htbp!]
  \centering
  \caption{Effect of constituency exposure to smoke on number of fires, robustness by treatment indicator (Area-Based Treatment)}\label{tab:app_main_by_treat_area}
  \resizebox{0.8\textwidth}{!}{
    \input{tables/_app_6_main_did_treat_definition_rural}
  }
\end{table}

\begin{table}[htbp!]
  \centering
  \caption{Effect of constituency exposure to smoke on number of fires, robustness by dependent variable (Area-Based Treatment)}\label{tab:app_main_bydv_area}
  \resizebox{0.8\textwidth}{!}{
    \input{tables/_app_7_main_did_downup_area_ac_dv_rural}
  }
\end{table}

\begin{table}[htbp!]
  \centering
  \caption{Effect of constituency exposure to smoke on number of fires, robustness by year (Area-Based Treatment)}\label{tab:app_by_year_area}
  \resizebox{\textwidth}{!}{
    \input{tables/_app_8_main_did_by_year_rural}
  }
\end{table}

\begin{table}[htbp!]
  \singlespacing \centering
  \caption{Effect of constituency exposure to smoke on number of fires, robustness by province (Area-Based Treatment)}\label{tab:app_by_province_area}
  \resizebox{0.8\textwidth}{!}{
    \input{tables/_app_9_main_did_by_state_rural}
  }
\end{table}

\begin{figure}[htbp!]
    \centering
    \caption{Effect of farmers' protests occurrence on number of fires (Area-Based Treatment)}\label{fig:app_main_protest_area}
    \subcaptionbox{By relative time period}[0.48\textwidth]{%
        \centering
        \includegraphics[width=\linewidth]{figures/_app_17_5km_fe12_evst_all_rural_fe03_1_rotated.png}
    }%
    \hfill
    \subcaptionbox{By constituency exposure to smoke}[0.48\textwidth]{%
        \centering
        \includegraphics[width=\linewidth]{figures/Interaction_downwind/_app_downup_rel_protest.png}
    }
\end{figure}

\begin{table}[htbp!]
    \centering
    \caption{Effect of Protest Occurrence on number of fires, by constituency exposure to smoke (Area-Based Treatment)}\label{tab:app_protest_downup_area}
    \resizebox{0.8\textwidth}{!}{
    {
\def\sym#1{\ifmmode^{#1}\else\(^{#1}\)\fi}
\begin{tabular}{l*{6}{c}}
\hline
            &\multicolumn{1}{c}{(1)}         &\multicolumn{1}{c}{(2)}         &\multicolumn{1}{c}{(3)}         &\multicolumn{1}{c}{(4)}         &\multicolumn{1}{c}{(5)}         &\multicolumn{1}{c}{(6)}         \\
            & \multicolumn{6}{c}{Number of Fires (in 1,000 units)} \\ \hline

Post $\times$ Protest&     154.607\sym{***}&     154.607\sym{***}&     104.445\sym{***}&     146.538\sym{***}&     146.538\sym{***}&      95.866\sym{***}\\
            &    (18.415)         &    (18.415)         &    (19.483)         &    (20.777)         &    (20.777)         &    (21.599)         \\
Down $>$ Up &                     &                     &                     &       7.206\sym{*}  &       7.206\sym{*}  &       6.398         \\
            &                     &                     &                     &     (4.357)         &     (4.357)         &     (4.428)         \\
Post $\times$ Down $>$ Up&                     &                     &                     &      -0.507         &      -0.507         &      -1.285         \\
            &                     &                     &                     &     (2.565)         &     (2.565)         &     (2.543)         \\
Protest $\times$ Down $>$ Up&                     &                     &                     &     -33.154         &     -33.154         &     -31.561         \\
            &                     &                     &                     &    (29.599)         &    (29.599)         &    (30.288)         \\
Post $\times$ Protest $\times$ Down $>$ Up&                     &                     &                     &      16.458         &      16.458         &      17.512         \\
            &                     &                     &                     &    (18.974)         &    (18.974)         &    (18.580)         \\
\hline
Observations&  63,278,278         &  63,278,278         &  63,278,278         &  63,278,278         &  63,278,278         &  63,278,278         \\
N Assembly Constituencies&         808         &         808         &         808         &         808         &         808         &         808         \\
Grid FE     &           Y         &           Y         &           Y         &           Y         &           Y         &           Y         \\
Relative Time FE&           Y         &           Y         &           Y         &           Y         &           Y         &           Y         \\
Legislature FE&           N         &           Y         &           Y         &           N         &           Y         &           Y         \\
Province Trend FE&           N         &           N         &           Y         &           N         &           N         &           Y         \\
Mean DV     &     473.109         &     473.109         &     473.109         &     473.109         &     473.109         &     473.109         \\
Mean DV2 (Down$>$Up=1)&                     &                     &                     &     458.614         &     458.614         &     458.614         \\
\hline
\end{tabular}
}

    }
\end{table}

\begin{figure}[htbp!]
    \caption{Effect by relative time period from protest (Area-Based Treatment)}
    \label{fig:app_es_protest_panel_area}
    \centering
    \begin{tabular}{cc}
        \includegraphics[width=0.450\textwidth]{figures/_app_17_5km_fe12_evst_all_rural_fe03_1_ori.png} &
        \includegraphics[width=0.450\textwidth]{figures/_app_17_5km_fe12_evst_all_rural_fe03_1_rot_honest2.png} & \\
        \small (a) Event-Study & \small (b) Honest DiD  \\
    \end{tabular}
\end{figure}

\begin{figure}[htbp!]
    \centering
    \caption{Effect of switching to an agriculturalist politician on number of fires (Area-Based Treatment)}\label{fig:app_main_agri_area}
    \subcaptionbox{By relative time period}[0.48\textwidth]{%
        \centering
        \includegraphics[width=\linewidth]{figures/_app_16_polischar_fe03_evst_main_rural_acpop_1_rotated.png}
    }%
    \hfill
    \subcaptionbox{By constituency exposure to smoke}[0.48\textwidth]{%
        \centering
        \includegraphics[width=\linewidth]{figures/Interaction_downwind/_app_downup_rel_polischar.png}
    }
\end{figure}

\begin{figure}[htbp!]
    \centering
    \caption{Effect of switching to an agriculturalist politician, by constituency agricultural characteristics (Area-Based Treatment)}\label{fig:app_es_agri_rice_area}
    \includegraphics[width=0.9\linewidth]{figures/_app_16_polischar_fe12_evst_all_rural_riceP_5_rotated.png}
\end{figure}

\begin{figure}[htbp!]
    \caption{Effect by relative time period from politician switch (Area-Based Treatment)}
    \label{fig:app_es_agri_panel_area}
    \centering
    \begin{tabular}{cc}
        \includegraphics[width=0.450\textwidth]{figures/_app_16_polischar_fe03_evst_main_rural_acpop_1_ori.png} &
        \includegraphics[width=0.450\textwidth]{figures/_app_16_polischar_fe03_evst_main_rural_acpop_1_rot_honest2.png} & \\
        \small (a) Event-Study & \small (b) Honest DiD  \\
    \end{tabular}
\end{figure}

\begin{table}[htbp!]
    \centering
    \caption{Effect of switching to an agriculturalist politician, by constituency exposure to smoke (Area-Based Treatment)}\label{tab:app_agri_profession_downwind_area}
    \resizebox{0.75\textwidth}{!}{
      {
\def\sym#1{\ifmmode^{#1}\else\(^{#1}\)\fi}
\begin{tabular}{l*{2}{c}}
\hline
            &\multicolumn{1}{c}{(1)}         &\multicolumn{1}{c}{(2)}         \\
            & \multicolumn{2}{c}{Number of Fires (in 1,000 units)} \\ \hline

Post $\times$ Agriculturalist&      12.465\sym{***}&      12.465\sym{***}\\
            &     (4.525)         &     (4.525)         \\
Down $>$ Up &                     &                     \\
            &                     &                     \\
Post $\times$ Down $>$ Up&                     &                     \\
            &                     &                     \\
Agriculturalist $\times$ Down $>$ Up&                     &                     \\
            &                     &                     \\
Post $\times$ Agriculturalist $\times$ Down $>$ Up&                     &                     \\
            &                     &                     \\
\hline
Observations&   8,268,839         &   8,268,839         \\
N Assembly Constituencies&         695         &         695         \\
Grid FE     &           Y         &           Y         \\
Relative Time FE&           Y         &           Y         \\
Legislature FE&           N         &           Y         \\
Province Linear Time Trend FE&           N         &           N         \\
Mean DV     &     105.231         &     105.231         \\
Mean DV (Down $>$ Up=1)&                     &                     \\
\hline
\end{tabular}
}

    }
\end{table}

\end{document}

to do:
tab 1: change the spacing + add number of ACs
es 1: make text larger and make it to the right where you have white background
change the Omitted time period to -1 and send it to the chat both not rotated and rotated - see Fenizia again
+ write up the interpretation 

ROTATE ALL EVENSTUDY AND THEN MOVE ALL THE NOT ROTATED AND HONEST IN APPENDIX (NR, R, HONEST AS A PANEL) AND MOVE THE ROTATED IN EMPIRICAL STRATEGY 

Add very carefully in the notes what the mean DV is + number of fires + x 1,000

remove downup AC in the first table

AC: for politicians
district: for bureaucrats
number of districts
add spacing between rows, but not the SE

move placebo area to appendix -> change the label ; population based -> write what the different columns are in the headers (e.g. full sample, downwind
move column 1 to a robustness as it is an alternative definition of fires. add the table as in Nature rob bytreat

politician characteristics: how to improve the es?
drop assets, assets and profession

try unstacked version of agri profession

PROTESTS: 5km area -> assembly FE


For the eventstudies we have 3 main figures usually in this order. 
\begin{itemize}
    \item Panel (A) effect for the pre and post treatment cofficients with a line that reports a linear trend estimated using the pre periods. 
    \item Panel (B) Reports the deviations from the eventstudy coefficient and the linear time trend
    \item Panel (C) reports the results of the sensitivity analysis using the honest approach to parallel trends of \citet{rambachan2023more}. The x-axis reports different values of \textit{M}. The approach imposes that the imposing that the post-treatment violation of parallel trends is no more than M larger than the maximum violation of parallel trends in the pre-treatment period. The value of M = 0.2, for example, imposes that the post-treatment violation is no bigger than 0.2 the worst violation of parallel trends in the pre-period. 
    \item The strategy used follows the approach by \citet{fenizia2024organized}
\end{itemize}
%\begin{equation}
%Y_{it} = \sum_{k \neq 0} \beta_k \cdot \text{EventTime}_{it}^{k} + \mathbf{X}_{it}'\gamma + \alpha_i + \delta_{acy} + \varepsilon_{it}
%\label{eq:event_study_exclude0}
%\end{equation}

define the area represented by an elected member in


1. cluster for personal incentives ac times election year
2. plot row 6 or 7
3. for all of them 
3 by 3 and table

To DO: try the classic DD not stacked, can we do it?

% \begin{table}[htbp!]
%    \caption{\label{tab:protest_treat_main} Effect of Level of Protests on Number of Fires 5km specification with Election Cycle}
%    \centering \resizebox{\textwidth}{!}{\input{tables/10_pr_cluster_2022_5kmtable_electionfe}}   
% \end{table}

% \begin{table}[htbp!]
%    \caption{\label{tab:protest_treat_main_pj} Main Specification with Protest as Treatment 5km}
%    \centering \resizebox{0.7\textwidth}{!}{
% \input{tables/protest_5km}
% }
% \end{table}
% \begin{table}[htbp!]
%    \caption{\label{tab:protest_treat_main_pj} Main Specification with Protest as Treatment 5km PJ and Haryana}
%    \centering \resizebox{0.7\textwidth}{!}{
% \input{tables/protest_5km_pj_hy}}
% \end{table}
